\chapter{\texorpdfstring{The engine table}{The engine table}}
\label{chap:engine-table}

\textbf{Purpose:} The machine, part by part: what each part computes, the exact algebra it holds to, what it does, what it wants, every hypothesis tried there with its result, and the next move.

\textbf{Layout:} orior's \texttt{src/\allowbreak{}engine} is in category folders (arithmetic, analysis, codecs, compiler, formats, nbody, prg\_\allowbreak{}sch, quantum, render, runtime, sims), with the tests under \texttt{test/\allowbreak{}engine}; item 13 of the list at the end maps every old folder to its new one. The benches are the root \texttt{bench/\allowbreak{}}, the tests the root \texttt{test/\allowbreak{}}, and every build writes only into the root \texttt{build/\allowbreak{}}: each build into its own \texttt{build/\allowbreak{}<date>\_\allowbreak{}<time>\_\allowbreak{}<name>/\allowbreak{}}, removed when a later build starts more than a day on, with the finished program copied up into \texttt{build/\allowbreak{}} and \texttt{build/\allowbreak{}logs/\allowbreak{}} and \texttt{build/\allowbreak{}tsv/\allowbreak{}} never removed. The six steps (ingest, seal, lift, measure, partition, relate) and every file's step are in engine/README.md.

\textbf{Scope:} \texttt{engine/\allowbreak{}}: exact arithmetic, the residual operator, the component tree, the moments, the correlation operator, the marginal, matching, the entropy history, the files (their size against the floor is tracked in compression\_\allowbreak{}table.md), the record machine, errors, the seal (integrity; its theory is \hyperref[chap:obsignatio-seal]{obsignatio\_\allowbreak{}seal.md}), the scheduler, the sift, render, the memory operations, the device pools, the period reading, the sims, the root universal, the root noise of a box, and the ask and the state. The programs that run on it are separate concerns with their own tables: cell tracking in cell\_\allowbreak{}tracking\_\allowbreak{}table.md, and RSNA knee MRI (its own project). Numbers are in \hyperref[chap:ledger]{ledger.md}. Statuses follow \hyperref[chap:readme]{README.md}. The protocol is \hyperref[chap:query-protocol-table]{query\_\allowbreak{}protocol\_\allowbreak{}table.md} and the plan engine\_\allowbreak{}plan.md.

\section{\texorpdfstring{The rule the machine is held to}{The rule the machine is held to}}

\textbf{The engine is optimized for no scale.} A size, a spacing, an order, a window, a width, a cell or a voxel count is never written into the machine. Each comes in with the request or is read from the data. The machine is exact at every scale its words can hold, and where a word is too narrow it says so (a request error, or \texttt{needed\_\allowbreak{}bits}) and never rounds. Tuning the machine to one program's scale is the error this table exists to catch, and the audit below lists what is still in.

The machine and each program are optimized against each other only through the interface at the end of this table. Merged, neither can be optimized on its own.

\section{\texorpdfstring{The engine in math}{The engine in math}}

The whole engine algorithm in math, from every workbook, with zero terms missed. Sources: every workbook file (this table, \hyperref[chap:vertical-time-compression]{vertical\_\allowbreak{}time\_\allowbreak{}compression.md}, \hyperref[chap:two-crystals]{two\_\allowbreak{}crystals.md}, \hyperref[chap:imprint-key-cycle]{imprint\_\allowbreak{}key\_\allowbreak{}cycle.md}, \hyperref[chap:keys-explained]{keys\_\allowbreak{}explained.md}, \hyperref[chap:kolmogorov-arnold]{kolmogorov\_\allowbreak{}arnold.md}, \hyperref[chap:noise-sieve-tower]{noise\_\allowbreak{}sieve\_\allowbreak{}tower.md}, \hyperref[chap:obsignatio-seal]{obsignatio\_\allowbreak{}seal.md}, \hyperref[chap:tessera-scheduler]{tessera\_\allowbreak{}scheduler.md}, compression\_\allowbreak{}table.md) and the code (\texttt{tower.\allowbreak{}cu}, \texttt{keymath.\allowbreak{}cu}, \texttt{cycle.\allowbreak{}c}, \texttt{pi\_\allowbreak{}tower.\allowbreak{}cu}). Each line carries the status its source gives it. What is not tied to the code is listed at the end.

\textbf{Objects.} A sample is u : Λ → [0, 2\textasciicircum{}16), Λ = [0,T)×[0,Z)×[0,Y)×[0,X), axes (t, z, y, x). A frame is u\_\allowbreak{}t = u(t, ·). The atom is all of the thing under inspection, held as one integer: A = Σ\_\allowbreak{}i v\_\allowbreak{}i·2\textasciicircum{}\{16i\}, lane i at place 2\textasciicircum{}\{16i\}. Every value is in ℤ; nothing is rounded.

\textbf{E0. Numbers} (M1, A0).

\begin{itemize}
\item Width: 2\textasciicircum{}k·32 bits, k ≥ 0, no ceiling (256 limbs in the engine build, fitted to the record file).
\item a·b: long multiplication, then Karatsuba from 32 limbs, then Schönhage–Strassen mod 2\textasciicircum{}n + 1 from 8,192 limbs (every twiddle a shift).
\item a = q·b + r, q = trunc(a/b), |r| < |b|, sgn r = sgn a: Knuth's D, and Newton's reciprocal once divisor and quotient reach 8,192 limbs. b = 0 is refused.
\item Exact a/b: b = 2\textasciicircum{}s·o with o odd. The quotient is ((a ≫ s)·o⁻¹) mod 2\textasciicircum{}\{32·limbs\}, with o⁻¹ from x ← x(2 − o·x); it is refused unless q·b = a.
\item gcd: Lehmer's (Knuth L); gcd(0, x) = |x|, gcd(0, 0) = 0.
\item A decimal (−1)\textasciicircum{}neg·digits·10\textasciicircum{}ex maps to ⌊|d|·2\textasciicircum{}\{−e2\}⌋ in W bits, with e2, \texttt{terminates}, \texttt{fits} and \texttt{needed\_\allowbreak{}bits}.
\end{itemize}

\textbf{E1. The record machine} (M10, A13): the evaluator every engine program runs on.

\begin{itemize}
\item A program is P = (s\_\allowbreak{}1, …, s\_\allowbreak{}n), s\_\allowbreak{}i = (op\_\allowbreak{}i, l\_\allowbreak{}i, r\_\allowbreak{}i, μ\_\allowbreak{}i), with l\_\allowbreak{}i, r\_\allowbreak{}i < i, and outputs O ⊆ \{1..n\}, no output twice.
\item A lane ℓ reads up to 3 records ρ\_\allowbreak{}μ = in\_\allowbreak{}μ[idx\_\allowbreak{}μ(ℓ)] (idx the identity when none is given). Registers v\_\allowbreak{}1..v\_\allowbreak{}n ∈ ℤ, v\_\allowbreak{}i = op\_\allowbreak{}i(v\_\allowbreak{}\{l\_\allowbreak{}i\}, v\_\allowbreak{}\{r\_\allowbreak{}i\}):

{\small
\setlength{\tabcolsep}{6pt}
\begin{longtable}{>{\RaggedRight\arraybackslash}p{\dimexpr0.2896\linewidth-2\tabcolsep\relax}>{\RaggedRight\arraybackslash}p{\dimexpr0.3313\linewidth-2\tabcolsep\relax}>{\RaggedRight\arraybackslash}p{\dimexpr0.3791\linewidth-2\tabcolsep\relax}}
\toprule
op & v\_\allowbreak{}i & width w\_\allowbreak{}i (keymath) \\
\midrule
\endhead
FIELD & bits [o, o+b) of ρ\_\allowbreak{}μ, unsigned & b \\
FIELD\_\allowbreak{}SIGNED & the same, signed & b \\
CONSTANT & c < 2\textasciicircum{}64 & bitlen(c) \\
SUM, DIFFERENCE & a ± b & the fewer of max(w\_\allowbreak{}a, w\_\allowbreak{}b) + 1 and the linear form's bound \\
PRODUCT & a·b & w\_\allowbreak{}a + w\_\allowbreak{}b; by a constant, the fewer of that and the linear form's bound \\
ABSOLUTE & |a| & w\_\allowbreak{}a \\
COMPARE & sgn(a − b) & 1 \\
LADDER (b > 0) & sgn(a)·\#\{k ∈ [1, 91] : F\_\allowbreak{}k·b ≤ |a|\}, F\_\allowbreak{}1 = F\_\allowbreak{}2 = 1 & bitlen(91) \\
QUOTIENT & trunc(a/b) & w\_\allowbreak{}a − ⌊log2 c⌋ (≥ 1) for a constant c, else w\_\allowbreak{}a \\
REMAINDER & a − b·trunc(a/b) & min(w\_\allowbreak{}a, w\_\allowbreak{}b) \\
GCD & gcd(a, b) ≥ 0 & max(w\_\allowbreak{}a, w\_\allowbreak{}b) \\
EXACT\_\allowbreak{}QUOTIENT & a/b, b | a, else the lane is refused & as QUOTIENT \\
XOR, AND & bitwise on two's complement without end & max + 1; AND with one never-negative operand = that operand's width, with two = the narrower; XOR of two never-negative = max \\
WRAP (w ≥ 4) & ((a + 2\textasciicircum{}\{w−1\}) mod 2\textasciicircum{}w) − 2\textasciicircum{}\{w−1\} & min(w\_\allowbreak{}a, w) \\
TABLE & T[|a| mod 2\textasciicircum{}b], 1 ≤ b ≤ min(32, w\_\allowbreak{}a) & out\_\allowbreak{}bits \\
\bottomrule
\end{longtable}
}
\item The linear form (A16): every register is c + Σ c\_\allowbreak{}i·v\_\allowbreak{}\{a\_\allowbreak{}i\} over atoms a\_\allowbreak{}i. SUM and DIFFERENCE add their operands' forms, PRODUCT by a constant (a form with no atoms) scales the other's, CONSTANT is its value, and a WRAP that passes its register through keeps that register's form. Every other register, FIELD included, is an atom. The bound is bitlen(|c| + Σ |c\_\allowbreak{}i|(2\textasciicircum{}\{w\_\allowbreak{}\{a\_\allowbreak{}i\}\} − 1)). A coefficient or constant past 2\textasciicircum{}62 makes the register an atom.
\item Never negative (N): FIELD, CONSTANT, ABSOLUTE, GCD and TABLE; SUM, PRODUCT, QUOTIENT, EXACT\_\allowbreak{}QUOTIENT or XOR of two in N; REMAINDER whose numerator is in N; AND with either in N; a WRAP that passes one in N through.
\item Refused at imprint: an operand that is not earlier, w\_\allowbreak{}i > 32·256, a WRAP below 4, a table index wider than its source. Refused per lane: b = 0, or an inexact EXACT\_\allowbreak{}QUOTIENT. A refused lane adds 1 to a device count, and a count ≠ 0 refuses the whole sweep as a request error (\texttt{cycle.\allowbreak{}c}).
\item Out: ⊕\_\allowbreak{}\{i∈O\} v\_\allowbreak{}i, each as two's complement over w\_\allowbreak{}i + 1 bits, the sum below 2\textasciicircum{}31 bits.
\item Layout: register i lives on [i, last reader of i]. With reuse, first fit in the file, ≤ 256 limbs; without it, one register per step at the running total. The kernel's file is 64 limbs when that holds the layout, else 256. A dividing program carries 5·WIDE + 4 limbs of scratch.
\item Sweep: R\_\allowbreak{}ℓ = P(ρ(ℓ)) for every ℓ, grid-stride, ≤ 4096 × 256 threads, one launch. The host runs P on the exact integer, and device = host word for word is the port's proof.
\item A member's field lies in its first 2\textasciicircum{}32 bits.
\item Laws:

\begin{itemize}
\item (h∘g)∘f = h∘(g∘f), g∘f ≠ f∘g. A stack of floors is one program. Time per lane ∝ n, independent of depth. Stacked against chained measured 13 to 22 times faster.
\item WRAP to w is π\_\allowbreak{}w : ℤ₂ → ℤ/2\textasciicircum{}w, window W\_\allowbreak{}w = [−2\textasciicircum{}\{w−1\}, 2\textasciicircum{}\{w−1\}); π\_\allowbreak{}w ∘ π\_\allowbreak{}\{w+1\} = π\_\allowbreak{}w.
\item For F over \{SUM, DIFFERENCE, PRODUCT, XOR, AND\}: π\_\allowbreak{}w ∘ F = F\_\allowbreak{}w ∘ π\_\allowbreak{}w (proved).
\item These do not factor through π\_\allowbreak{}w: ABSOLUTE, COMPARE, QUOTIENT, REMAINDER, GCD, LADDER, TABLE, ⌊v/2\textasciicircum{}k⌋ (through π\_\allowbreak{}\{w+k\}), and EXACT\_\allowbreak{}QUOTIENT by an even c.
\item For c odd with c | v: WRAP(v/c, w) = WRAP(WRAP(v, w)·(c⁻¹ mod 2\textasciicircum{}w), w). The kernel grows c⁻¹ by Newton: 6, 12, 24, 48 bits in a word, then doubling.
\item Mod an odd p\textasciicircum{}v, REMAINDER is the projection ℤ\_\allowbreak{}p → ℤ/p\textasciicircum{}v, and SUM, DIFFERENCE and PRODUCT commute with it. The CRT join y\_\allowbreak{}2 + 2\textasciicircum{}8·(((y\_\allowbreak{}p − y\_\allowbreak{}2)·2\textasciicircum{}\{−8\}) rem m) is the run mod 2\textasciicircum{}8·m. XOR breaks mod 3.
\item τ\_\allowbreak{}k(x) = trunc(x/2\textasciicircum{}k): τ\_\allowbreak{}j∘τ\_\allowbreak{}k = τ\_\allowbreak{}\{j+k\}; order is kept weakly; τ\_\allowbreak{}k(x + y) − τ\_\allowbreak{}k(x) − τ\_\allowbreak{}k(y) ∈ \{−1, 0, 1\}.
\end{itemize}
\item Keys, the linear programs (imprint\_\allowbreak{}key\_\allowbreak{}cycle, keys\_\allowbreak{}explained):

\begin{itemize}
\item key = P(δ), δ the impulse. For linear shift-invariant P, P(x) = key ∗ x.
\item w·x = Σ over the set bits b of w of x ≪ b.
\item With lanes wide enough that none carries, A·K holds the smoothed lanes, and K = K\_\allowbreak{}z·K\_\allowbreak{}y·K\_\allowbreak{}x.
\item A lane below 2\textasciicircum{}m under a row of sum S stays below 2\textasciicircum{}\{m + bitlen(S − 1)\}.
\item The CRC-64 key: the advance over 2\textasciicircum{}k zero bytes is A\textasciicircum{}\{2\textasciicircum{}k\}, 48 matrices (24,576 bytes), and crc(a ⊕ b) = crc(a) ⊕ crc(b) ⊕ crc(0).
\item The mirror fold repeats with period 2N: one period is laid as a table of offsets.
\end{itemize}
\item Tables: (g∘f)[x] = g[f[x]]. Two never-negative u, v with b\_\allowbreak{}u + b\_\allowbreak{}v ≤ 32 index as u·2\textasciicircum{}\{b\_\allowbreak{}v\} + v. A table is 2\textasciicircum{}b·⌈out\_\allowbreak{}bits/32⌉·4 bytes.
\item The lane index as a register and the latch, first lane = min\{ℓ : cond(ℓ)\}, are built, with the broadcast. \texttt{LANE} (\texttt{ENGINE\_\allowbreak{}RECORD\_\allowbreak{}LANE}, 18, \texttt{engine\_\allowbreak{}config.\allowbreak{}h}) reads the number ℓ the sweep runs the lane as, 64 bits and never negative. The broadcast: a member holding one record is read by every lane that has no index, and one shared record with x = base + ℓ enumerates a range with nothing stored a lane. The latch is \texttt{cycle\_\allowbreak{}record\_\allowbreak{}latch} (\texttt{cycle\_\allowbreak{}record\_\allowbreak{}launch.\allowbreak{}cu}): a scan in each thread, a warp shuffle tree and one atomic minimum a block, returning \texttt{CYCLE\_\allowbreak{}LATCH\_\allowbreak{}NONE} (\texttt{cycle.\allowbreak{}h}) where no lane's condition holds. \texttt{cycle\_\allowbreak{}record\_\allowbreak{}latch\_\allowbreak{}host} (\texttt{cycle.\allowbreak{}c}) is the same on the host. A latch through the engine's own entry is not built. \texttt{record\_\allowbreak{}lane\_\allowbreak{}test} is 17 checks, 0 failed.
\item The sum across lanes is built: \texttt{cycle\_\allowbreak{}record\_\allowbreak{}sum} (\texttt{cycle\_\allowbreak{}record\_\allowbreak{}launch.\allowbreak{}cu}) adds one output of every record over each run of consecutive lanes, exactly, as two's complement. Each 32-bit limb is summed into its own 64-bit column by the lanes of a thread and one atomic add a run, the negative fields are counted beside, and the host takes the carries. It merges nothing: every column is exact, the whole sum is kept, and each lane's record stays where it lies. \texttt{cycle\_\allowbreak{}record\_\allowbreak{}sum\_\allowbreak{}host} (\texttt{cycle.\allowbreak{}c}) adds in order and is its port check. \texttt{record\_\allowbreak{}sum\_\allowbreak{}test} is 98 checks, 0 failed: sums of ℓ, -ℓ, ℓ² and -ℓ² over 2\textasciicircum{}22 lanes meet their closed forms, and random fields up to nine limbs meet the host word for word.
\end{itemize}

\textbf{E2. Ingest} (M9). Source → (u, σ). A signed DICOM pixel p ↦ p + 2\textasciicircum{}15. Slices are ordered by an exact key from their positions (E0). σ = the headers, deflated by our LZ77 and Huffman.

\textbf{E3. The seal} (M13, A10).

\begin{itemize}
\item K\_\allowbreak{}ℓ = BLAKE3-derive-key(“obsignatio aeterna 2026-09-23 “ ‖ ℓ), H\_\allowbreak{}ℓ(m) = BLAKE3-keyed(K\_\allowbreak{}ℓ, m). The levels ℓ: row, plane, volume, lanes, chunk, stream, side stored, side inflated, side, members, sample, set, file, universal.
\item ρ(t, z, y) = H\_\allowbreak{}row(u[t, z, y, ·]), π(t, z) = H\_\allowbreak{}plane(ρ(t, z, ·)), υ(t) = H\_\allowbreak{}volume(π(t, ·)), Λ = H\_\allowbreak{}lanes(υ(·)).
\item κ(c) = H\_\allowbreak{}chunk(LE64(ℓ\_\allowbreak{}c) ‖ bits\_\allowbreak{}c); Σ = H\_\allowbreak{}stream(LE64(o\_\allowbreak{}0 … o\_\allowbreak{}\{C−1\}) ‖ LE64(B) ‖ κ(0) ‖ … ‖ κ(C−1)).
\item σ = H\_\allowbreak{}side(H\_\allowbreak{}side stored(deflated) ‖ H\_\allowbreak{}side inflated(inflated)), μ = H\_\allowbreak{}members(the six tables ‖ the names).
\item Ω = H\_\allowbreak{}sample(LE64(extent, C, B, lane offset, members, side bytes, deflated bytes, name bytes, lane nodes) ‖ Λ ‖ Σ ‖ σ ‖ μ). Θ = H\_\allowbreak{}set(Ω\_\allowbreak{}1 ‖ … ‖ Ω\_\allowbreak{}n), in the order named.
\item Check order: κ, Σ, σ\_\allowbreak{}s, μ, then σ\_\allowbreak{}i, σ, then ρ … Λ, then Ω.
\item Locality: a change at (t, z, y, x) changes exactly ρ, π, υ, Λ, Ω and Θ on its path; a change in chunk c changes κ(c), Σ, Ω and Θ. Any other node agrees except with probability 2\textasciicircum{}−256.
\item The device fold (pair each level, carry the odd node up) is BLAKE3's left-heavy tree. H is not associative, and the shape is therefore fixed.
\item The witness: the decoded voxels against the source on the device, then a second independent read, pixel for pixel.
\item Ruled and not built: H\_\allowbreak{}universal over every file the engine touches plus the engine's own hash.
\end{itemize}

\textbf{E4. The lift} (M9, M20, A8, A14).

\begin{itemize}
\item One level along one axis of length N, on the active block:

\begin{itemize}
\item d\_\allowbreak{}j = x\_\allowbreak{}\{2j+1\} − ⌊(x\_\allowbreak{}\{2j\} + x\_\allowbreak{}\{2j+2\})/2⌋, j < ⌊N/2⌋;
\item s\_\allowbreak{}i = x\_\allowbreak{}\{2i\} + ⌊(d\_\allowbreak{}\{i−1\} + d\_\allowbreak{}i + 2)/4⌋, i < ⌈N/2⌉;
\item an edge repeats its neighbor: x\_\allowbreak{}\{2j+2\} := x\_\allowbreak{}\{2j\} at the end, d\_\allowbreak{}\{−1\} := d\_\allowbreak{}0, d\_\allowbreak{}\{⌊N/2⌋\} := d\_\allowbreak{}\{i−1\};
\item low-pass coefficients stored on [0, ⌈N/2⌉), high-pass coefficients on [⌈N/2⌉, N);
\item ⌊v/2\textasciicircum{}k⌋ is taken toward −∞ (\texttt{tower\_\allowbreak{}floor\_\allowbreak{}shift}).
\end{itemize}
\item A level runs the axes 0..3 whose extent is ≥ 2, in order. The extents then go e ↦ ⌈e/2⌉, and the next level lifts only the all-low-pass corner. There are L levels, until every extent is 1.
\item T⁻¹: x\_\allowbreak{}\{2i\} = s\_\allowbreak{}i − ⌊(d\_\allowbreak{}\{i−1\} + d\_\allowbreak{}i + 2)/4⌋, then x\_\allowbreak{}\{2j+1\} = d\_\allowbreak{}j + ⌊(x\_\allowbreak{}\{2j\} + x\_\allowbreak{}\{2j+2\})/2⌋. Levels and axes run in reverse.
\item The ruleset, the wave transform's math. A line splits into its low-pass coefficients L (the evens) and high-pass coefficients H (the odds), and a lifting step moves one band by a rounded sum over the other:

\begin{itemize}
\item target\_\allowbreak{}a ← target\_\allowbreak{}a + σ·⌊(r + Σ\_\allowbreak{}k w\_\allowbreak{}k·other\_\allowbreak{}\{a+o\_\allowbreak{}k\}) / 2\textasciicircum{}s⌋, σ = ±1, an index past either end of \texttt{other} taken at that end;
\item a ruleset R = (ρ\_\allowbreak{}1, …, ρ\_\allowbreak{}m) runs its steps in order along every line of every level. T⁻¹ runs them last first with each σ turned. A step moves only its target, reading only the other band, and every ruleset is therefore invertible exactly, whatever its weights;
\item the kernels' 5/3 is ρ\_\allowbreak{}1 = (H, σ = −1, o = (0, 1), w = (1, 1), r = 0, s = 1), then ρ\_\allowbreak{}2 = (L, σ = +1, o = (−1, 0), w = (1, 1), r = 2, s = 2);
\item two rulesets the kernels do not run are the S transform, (H, −1, (0), (1), 0, 0) then (L, +1, (0), (1), 0, 1), and a four-tap predict, (H, −1, (−1, 0, 1, 2), (−1, 9, 9, −1), 8, 4), with the 5/3's update;
\item on the record machine a weight other than ±1 is a PRODUCT by a constant, which keymath's linear forms carry (A16), and each constant is one step, laid where it is first read.
\end{itemize}
\item Edges: E\_\allowbreak{}k(w) = w − (w mod 2\textasciicircum{}b) + π\_\allowbreak{}k(w mod 2\textasciicircum{}b), π\_\allowbreak{}k a permutation of [0, 2\textasciicircum{}b), 1 ≤ b ≤ 20.

\begin{itemize}
\item Lift = E\_\allowbreak{}L ∘ T\_\allowbreak{}L ∘ … ∘ E\_\allowbreak{}1 ∘ T\_\allowbreak{}1 ∘ E\_\allowbreak{}0. Lower is the inverse, last laid first.
\item Two edges on one floor fold into π\_\allowbreak{}b ∘ π\_\allowbreak{}a; a level between them blocks the fold.
\item A non-bijective f rides as (x, y) ↦ (x, y ⊕ f(x)) (Bennett), at the price of b\_\allowbreak{}x + b\_\allowbreak{}y ≤ 20 bits.
\end{itemize}
\item |c| ≥ 2\textasciicircum{}30 refuses the lift.
\item The lattice lifted is u's 16-bit lanes or a lattice of ints (\texttt{device\_\allowbreak{}values} on the lift's request), such as a residual left below zero (E16). An int with |v| ≥ 2\textasciicircum{}30 refuses the lift, as a coefficient does.
\item Stored (\texttt{compression.\allowbreak{}cu}):

\begin{itemize}
\item zigzag: z(c) = 2c for c ≥ 0, 2|c| − 1 for c < 0;
\item per block of 64: g = bitlen(⌊Σz/64⌋); k is the least-cost k ∈ [max(g − 2, 0), min(g + 1, 31)], with cost(k) = 5 + Σ\_\allowbreak{}z (⌊z/2\textasciicircum{}k⌋ < 24 ? ⌊z/2\textasciicircum{}k⌋ + 1 + k : 24 + 32);
\item written as k in 5 bits, then per value q = ⌊z/2\textasciicircum{}k⌋ ones and a zero and the low k bits, or 24 ones and the 32-bit value;
\item 64 blocks a chunk, and chunks sit at stored bit offsets: any floor and any chunk is read directly (the elevator E(f, t)).
\end{itemize}
\item Held when T⁻¹(decode(stream)) = u lane for lane, and u equals a second read of the source.
\item Properties (proved along one line on 64 samples at L = 4 unless marked):

\begin{itemize}
\item T is a bijection of ℤ\textasciicircum{}n, a homeomorphism of ℤ₂\textasciicircum{}n (derived), and keeps Haar measure.
\item Reach: T's level-ℓ low-pass coefficients read 3ℓ bits down, its high-pass coefficients 3ℓ − 2; T⁻¹ reads L + 2. So |T(x) − T(y)|₂ ≤ 2\textasciicircum{}\{3L\}|x − y|₂.
\item Linear part M, entries in ℤ[1/2]: T(x + 2\textasciicircum{}\{3L\}z) = T(x) + M·2\textasciicircum{}\{3L\}z. A constant c moves only the level-L low-pass coefficients, by c. Negation and doubling do not pass. det M = +1.
\item Haar: every output quantum at level w holds exactly 2\textasciicircum{}\{3Ln\} input quanta at level w + 3L.
\item ℝ: 2\textasciicircum{}\{−k\}·T(trunc(2\textasciicircum{}k x)) → Mx. ℤ\_\allowbreak{}p, p odd: T does not extend; M does, in SL\_\allowbreak{}n(ℤ\_\allowbreak{}p).
\item Widths: keymath's linear forms (A16) give every register at level ℓ, low-pass and high-pass, w + ℓ + 1 bits: the first level adds 2, each after it 1 (derived by hand; the ring below matches). Values obey |s| ≤ 1.5B + 1 and |d| ≤ 2B (+0.585 and +1 bits a level).
\item Ring: ring\_\allowbreak{}ℓ = ring\_\allowbreak{}0 + n + 2n(1 − 2\textasciicircum{}\{−ℓ\}) for ℓ ≥ 1 (measured: 1,536, 1,664, 1,696, 1,712 and 1,720 at n = 64, w = 24). The mirror floor still adds n/2\textasciicircum{}ℓ (1,720, 1,712, 1,696 and 1,600), and the heap mirrors exactly across the crystal.
\item Code length: Σ\_\allowbreak{}r ⌈log₂(W\_\allowbreak{}r + 1)⌉ + Σ\_\allowbreak{}r m\_\allowbreak{}r ≥ K(x | program) − c.
\item Identity by null permutation: a lane is identified when heap(T(x)) < heap(T(σ\_\allowbreak{}i x)) for every i ≤ d. Noise passes at a rate ≤ 1/(d + 1). T(σx) = T(x) ⇔ σx = x.
\item The KA step: x ← x + Φ(Σ\_\allowbreak{}p φ\_\allowbreak{}p(x\_\allowbreak{}p)), with Φ(s) = ⌊s/2\textasciicircum{}k⌋.
\item Matching on a floor. Floor 2 of a 64³ frame is its 16³ corner after two levels, m = n/64 values. Laid once as 16 bit planes, it gives every position holding x as ∧\_\allowbreak{}j (x\_\allowbreak{}j ? P\_\allowbreak{}j : ¬P\_\allowbreak{}j), 32 positions a word, reading no value and at most 16m/32 = n/128 plane words; a word stops once its mask is empty. \textbf{Proved} (\texttt{floor\_\allowbreak{}match}, 11 checks, 0 failed): the engine's floor 2, the crystal's corner lowered as its own tower, equals the host's two-level lifting at 4,096 of 4,096; 1,024 values on floor 2 and 1,024 not, every match set equal to the host's scan; the most any query read was 1,340 plane words, against n/2 = 131,072 samples. The index costs one read of a, the lift, paid once.
\end{itemize}
\item As record floors, emitted by the engine itself (\texttt{tower\_\allowbreak{}record\_\allowbreak{}lift} and \texttt{tower\_\allowbreak{}record\_\allowbreak{}lower}, each floor focused from a ruleset; A13):

\begin{itemize}
\item ⌊v/2\textasciicircum{}k⌋ = EXACT\_\allowbreak{}QUOTIENT(v − AND(v, 2\textasciicircum{}k − 1), 2\textasciicircum{}k);
\item E(v) = (v − AND(v, 2\textasciicircum{}k − 1)) + TABLE\_\allowbreak{}π(AND(v, 2\textasciicircum{}k − 1));
\item F ∘ T⁻¹ is one stack, and T⁻¹ ∘ T = id;
\item T⁻¹GT is G conjugated, (T⁻¹G₂T)(T⁻¹G₁T) = T⁻¹G₂G₁T, and a G on level-ℓ high-pass coefficients costs ℓ levels each way;
\item a linear F: |F(T⁻¹c) − F(W⁻¹c)| ≤ L·‖F‖∞ along one line;
\item the wrap at the mirror: WRAP(v, b + 1) on each rebuilt register, b its forward twin's width.
\end{itemize}
\end{itemize}

\textbf{E5. Measure} (M2, M8, M18), per frame.

\begin{itemize}
\item Residual: L\_\allowbreak{}n = [1 1]\textasciicircum{}\{∗n\} per axis, undivided; B\_\allowbreak{}s ∗ B\_\allowbreak{}b = B\_\allowbreak{}\{s+b\}; R = 2\textasciicircum{}\{|b|\}·L\_\allowbreak{}s u\_\allowbreak{}t − L\_\allowbreak{}\{s+b\} u\_\allowbreak{}t.

\begin{itemize}
\item ΣB\_\allowbreak{}s = 2\textasciicircum{}\{|s|\}, and a constant cancels.
\item The heat identity: 4\textasciicircum{}m δ − H\textasciicircum{}m = Σ\_\allowbreak{}\{k<m\} 4\textasciicircum{}\{m−1−k\} H\textasciicircum{}k (4δ − H), with H = [1 2 1] and 4δ − H = −Δ².
\item Parity: 2δ − [1 1] = [1 −1], half a voxel off. An odd b is therefore refused and an odd s is held with its offset (ruling (a)).
\item The key is 2\textasciicircum{}g·B\_\allowbreak{}narrow − B\_\allowbreak{}wide·B\_\allowbreak{}narrow, in 9 limbs.
\item Readings of any declared width w (the unit sweep's planes): R leaves in ⌈(w + Σs + Σb + 1)/32⌉ limbs, and a reading ≥ 2\textasciicircum{}w is refused.
\item The positive set: P(x) = [R(x) > 0].
\end{itemize}
\item The history:

\begin{itemize}
\item f\_\allowbreak{}j(x, W) = Σ\_\allowbreak{}\{t∈W\} bit\_\allowbreak{}j(u\_\allowbreak{}t(x) ⊕ u\_\allowbreak{}\{t−1\}(x)), |W| = 11 transitions, ≤ 16 windows;
\item held: Σ\_\allowbreak{}W f\_\allowbreak{}j ≡ bit\_\allowbreak{}j(u\_\allowbreak{}0(x) ⊕ u\_\allowbreak{}\{F−1\}(x)) (mod 2) at every x and j;
\item the anchor count a\_\allowbreak{}j(x) = Σ\_\allowbreak{}t bit\_\allowbreak{}j(u\_\allowbreak{}t(x)); the anchor bits are ∧\_\allowbreak{}t u\_\allowbreak{}t and ∧\_\allowbreak{}t ¬u\_\allowbreak{}t;
\item the cloud: D\_\allowbreak{}W(x) = Σ\_\allowbreak{}j f\_\allowbreak{}j(x, W)·2\textasciicircum{}j, and O(W, W′) = Σ\_\allowbreak{}x D\_\allowbreak{}W(x)·D\_\allowbreak{}\{W′\}(x);
\item the knf section s(x) = (D\_\allowbreak{}W(x))\_\allowbreak{}W less its mean over W; E = Σ\_\allowbreak{}x Σ\_\allowbreak{}\{a∈\{z,y,x\}\} ⟨s(x), s(x + e\_\allowbreak{}a)⟩ on the torus; the departure is d·E − Σ\_\allowbreak{}i E(σ\_\allowbreak{}i x).
\end{itemize}
\item The period, per axis of extent n, M = N/n:

\begin{itemize}
\item L = ⌊n/2⌋; the starts are the positions below n − L on the axis, Q = (n − L)·M;
\item a(ℓ) = \#\{x : u(x) = u(x + ℓ·e)\};
\item a candidate p has 2 ≤ p, 2p + 1 ≤ L, a(p) > max(a(p ± 1)) and a(2p) > max(a(2p ± 1));
\item h(p) = min(a(p) − max(a(p ± 1)), a(2p) − max(a(2p ± 1)));
\item the null draw k: each line along the axis permuted by a keyed Fisher–Yates (swap s with h mod (s + 1), 4 keys from the content signum and c = draw·rank + axis), h\_\allowbreak{}k = max\_\allowbreak{}p h;
\item P = min\{p : h(p)/Q > max\_\allowbreak{}k h\_\allowbreak{}k/Q\}, a 128-bit cross-multiply; P = 0 when none clears.
\item False-period rate ≤ 1/(d + 1).
\end{itemize}
\end{itemize}

\textbf{E6. Partition} (M3, M4, A2, A3).

\begin{itemize}
\item Admitted: R(x) > 0. The code c(x) = (dense rank of R(x) over the admitted) + 1, and 0 elsewhere.
\item Adjacency: 6 (a face joins x and x + e\_\allowbreak{}a, a ∈ \{z, y, x\}, both admitted).
\item The face key is (c of the weaker end, then ¬(3x + a) in the lowest limb). A stronger weaker end wins and then the lower face name. The tree is the maximum spanning forest by these keys (\texttt{max\_\allowbreak{}tree\_\allowbreak{}contract}, \texttt{max\_\allowbreak{}tree.\allowbreak{}cu}):

\begin{itemize}
\item Borůvka rounds: each component takes its strongest outgoing face;
\item of a mutual pair the higher root hooks onto the lower;
\item pointer jumping between rounds;
\item the chosen faces are the forest.
\end{itemize}
\item Components at every level from the forest: count(r) = \#\{x : c(x) ≥ r\} − \#\{forest faces with min(c) ≥ r\}, one scan from the top. The cut is r* = min\{r : count(r) = max count\}. The bodies are the components of \{c ≥ r*\}, labeled by min-label spreading along the forest. The count found must equal count(r*), else it is a logic error.
\item A node n has a level λ(n), and ν(x) is x's own node.

\begin{itemize}
\item The component at level r holding x is the ancestor a of ν(x) with λ(a) ≥ r > λ(parent(a)).
\item A subtree is a DFS range [in(a), out(a)).
\item x and x′ share a component at r exactly when r ≤ λ(LCA(ν(x), ν(x′))).
\end{itemize}
\item Ties: every face's key carries the face's name in its lowest limb, and no two keys tie.
\item Moments: μ(n) = Σ\_\allowbreak{}\{ν(x)=n\} (1, x, x xᵀ), and each node's are Σ over sub(n). The centered K = m·M − S·Sᵀ is never divided. ℓ ≤ 2 harmonics = (m, S, M).
\item Built and proved on the host (A2): the probe partition from k − 1 LCA levels, and J(v)[n, n′] with C\_\allowbreak{}\{r,r′\} a rectangle sum over two DFS ranges. Nothing calls them yet, and the slide still bisects over levels (\texttt{slide\_\allowbreak{}score.\allowbreak{}cu} :97 and :207 in biohub). Theory and not built: moments by one DFS prefix.
\end{itemize}

\textbf{E7. Relate} (M5, M6, M7).

\begin{itemize}
\item Agreement: A(v) = Σ\_\allowbreak{}x P\_\allowbreak{}t(x)·P\_\allowbreak{}\{t′\}(x + v), by an NTT mod 998,244,353.

\begin{itemize}
\item Each axis of extent n is padded to the least power of 2 ≥ 2n − 1, and every lag −(n − 1)..(n − 1) is linear, never wrapped.
\item It is refused when |Λ\_\allowbreak{}t| ≥ the prime or the padded total > 2\textasciicircum{}31 − 1, and the longest padded axis is ≤ 2\textasciicircum{}23.
\item v* = argmax A, ties to the least Σ\_\allowbreak{}a w\_\allowbreak{}a·v\_\allowbreak{}a², then the least index.
\end{itemize}
\item \texttt{C\_\allowbreak{}\{t,t′\}[a, b](v)} = Σ\_\allowbreak{}x [L\_\allowbreak{}t(x) = a]·[P\_\allowbreak{}\{t′\}(x + v)]·[L\_\allowbreak{}\{t′\}(x + v) = b]. Its reductions:

\begin{itemize}
\item the argmax over v;
\item the climb (\texttt{climb\_\allowbreak{}machine.\allowbreak{}cu}): S(v′) is the climber's count at lag v′, over v′ ∈ v + \{−1, 0, 1\}³. It moves to the v′ ≠ v with S(v′) > S(v), the largest, ties to the least w\_\allowbreak{}z·z² + y² + x². It stops when none exceeds S(v), and holds S(v). The landing is a majority vote (one pass picks, one counts) of the leaf under each of the climber's voxels at the lag and along a spiral around it, the most agreed kept;
\item the transpose, \texttt{C\_\allowbreak{}\{t′,t\}[b, a](−v)} = \texttt{C\_\allowbreak{}\{t,t′\}[a, b](v)};
\item box sums at nested widths from one summed-area table;
\item the tent;
\item the moments, at v = 0 with weight φ ∈ \{1, x, x xᵀ\}.
\end{itemize}
\item Marginal:

\begin{itemize}
\item Z = Σ\_\allowbreak{}α Π\_\allowbreak{}s ω(s, α(s)), α injective with ∅ allowed, is the permanent of [ω(s, t) | diag ω(s, ∅)];
\item P(s → t) = Z\_\allowbreak{}\{s→t\}/Z, and it factors over the support's components;
\item bounds: ≤ 2\textasciicircum{}16 arrangements, ≤ 64 targets, 16 limbs.
\end{itemize}
\item Matching: M* = argmax over one-to-one M of Σ\_\allowbreak{}\{(a,b)∈M\} w(a, b), in exact integers.
\item Body programs (velocity, division, contact side, fingerprint, print pair) are E1 programs over packed body records, ≤ 3 members.
\item Null draws on the identity line: frame f stands on floor ⌊f/11⌋ at place f mod 11. Its draw on floor g is 11g + (f mod 11), the floors taken by least O.
\end{itemize}

\textbf{E8. The scheduler} (M14, A11).

\begin{itemize}
\item The quantities: device C, in use U(τ); job j has d\_\allowbreak{}j, its standing s\_\allowbreak{}j (the bytes its process held as it asked: its context and whatever it kept, by pid, or by the process's own report under WSL), r\_\allowbreak{}j (only growing) and u\_\allowbreak{}j(τ) by pid.
\item O(τ) = U(τ) − Σ\_\allowbreak{}j u\_\allowbreak{}j(τ); H(τ) = C − O(τ) − Σ\_\allowbreak{}j max(r\_\allowbreak{}j, u\_\allowbreak{}j(τ)).
\item A job wants w\_\allowbreak{}k = max(s\_\allowbreak{}k + d\_\allowbreak{}k, p\_\allowbreak{}σ) and needs n\_\allowbreak{}k = w\_\allowbreak{}k − s\_\allowbreak{}k, since U already counts s\_\allowbreak{}k. The p\_\allowbreak{}σ term is committed and unapproved; Doug's rule is r\_\allowbreak{}k = d\_\allowbreak{}k.
\item Admit k when n\_\allowbreak{}k ≤ H, with r\_\allowbreak{}k = w\_\allowbreak{}k and u\_\allowbreak{}k = s\_\allowbreak{}k until its first sweep.
\item Growth: u\_\allowbreak{}j > r\_\allowbreak{}j sets r\_\allowbreak{}j = u\_\allowbreak{}j and warns. On release, p\_\allowbreak{}σ = max\_\allowbreak{}τ u\_\allowbreak{}j(τ) under the signum σ = BLAKE3(request).
\item Asked: d\_\allowbreak{}k > p\_\allowbreak{}σ, the declaration alone. The job is held, overridable, and lost after its holding time.
\item Backfill with the head kept:

\begin{itemize}
\item the shadow S = least time with H + Σ\_\allowbreak{}\{e\_\allowbreak{}j ≤ S\} max(r\_\allowbreak{}j, u\_\allowbreak{}j) ≥ n\_\allowbreak{}h, e\_\allowbreak{}j = start\_\allowbreak{}j + t\_\allowbreak{}σj;
\item the spare X = that sum − n\_\allowbreak{}h;
\item k goes now if n\_\allowbreak{}k ≤ H and (now + t\_\allowbreak{}σk ≤ S or n\_\allowbreak{}k ≤ X);
\item a signum with no history is never backfilled.
\end{itemize}
\item Every deadline sits in one min-heap on time to change.
\item A ticket is (pid, start time).
\end{itemize}

\textbf{E9. Errors} (M12). Every entry returns a value or refuses with (kind ∈ \{request, resource, logic\}, module, site, status, execaddr, evacaddr, imagebase, ≤ 16 frames). First write wins. A logic error withholds the result.

\textbf{E10. The sift, the ladder, memory, render} (M11, M15, M16, M17).

\begin{itemize}
\item Sift dispatch: in order when 100·total² against 85·distinct·Σcount² says so, ≥ 143 bits. Anchors are 1 when a period P is found, else 4. A periodic corpus filters no lower than 1/P. H₂ is the collision entropy.
\item β(n) = \#\{Fibonacci rungs ≤ n after the first\}, 92 rungs below 2\textasciicircum{}64. LADDER(a, b) reads β of |a|/b.
\item The memory arms answer exactly what a byte loop answers; render's host bytes equal the device's.
\end{itemize}

\textbf{E11. The ask and the state} (M21, A15).

\begin{itemize}
\item p\_\allowbreak{}i = Tr(ρE\_\allowbreak{}i), E\_\allowbreak{}i = (I + v\_\allowbreak{}i·σ)/4, which gives p\_\allowbreak{}i = (1 + r·v\_\allowbreak{}i)/4.
\item With Σv\_\allowbreak{}i = 0 and Σv\_\allowbreak{}i v\_\allowbreak{}iᵀ = c·I, r = (4/c)·Σ p\_\allowbreak{}i v\_\allowbreak{}i.
\item Crossing: (1 + r\_\allowbreak{}a)/2 = Σ p\_\allowbreak{}i (½ + (2/c)v\_\allowbreak{}\{i,a\}).
\item Valid ⇔ |r|² ≤ 1.
\item ℚ(√3): (a + b√3)(c + d√3) = (ac + 3bd) + (ad + bc)√3, the sign from a² against 3b².
\end{itemize}

\textbf{E12. The limits and the ordered machine} (two\_\allowbreak{}crystals).

\begin{itemize}
\item The inverse limit ℤ₂ = lim← ℤ/2\textasciicircum{}w; the direct limit of ℤ/2\textasciicircum{}w under x ↦ 2x is ℤ[1/2]/ℤ, its dual. ℤ is dense in ℤ₂, and ℤ₂ \textbackslash{} ℤ is uncountable.
\item The solenoid Σ₂ = (ℝ × ℤ₂)/ℤ. The adeles 𝔸/ℚ ≅ (ℝ × Ẑ)/ℤ, Ẑ = Π\_\allowbreak{}p ℤ\_\allowbreak{}p, and |x|\_\allowbreak{}∞·Π\_\allowbreak{}p |x|\_\allowbreak{}p = 1.
\item Orders: F\textasciicircum{}k for k ∈ ℤ, F\textasciicircum{}\{−k\} = (F⁻¹)\textasciicircum{}k. The counter n ↦ n + 1 is a shear.
\item ±ω: each bit ← lim sup. On a finite window that is the OR over the cycle, and +ω = −ω.
\end{itemize}

\textbf{E13. Programs on the engine} (M19).

\begin{itemize}
\item The camera law:

\begin{itemize}
\item S = background + ramp·x + wave + Σ bodies' brightness;
\item shot noise, a count in which each of the S expected electrons adds 0, 1, 2 or 4 with chances 9, 8, 6 and 1 in 24, whose first four cumulants are each S, Poisson's (\texttt{sim\_\allowbreak{}poisson\_\allowbreak{}four\_\allowbreak{}cumulants}); the symmetric law Binomial(4S, ½) − S is kept beside it. Read noise Binomial(4r², ½) − 2r²;
\item value = offset + fixed pattern + gain·S + read;
\item every draw is splitmix64(key ⊕ counter).
\end{itemize}
\item Ω (\texttt{chaitin\_\allowbreak{}omega}): Ω = Σ over closed BLC terms t with a normal form of 2\textasciicircum{}\{−|t|\}.

\begin{itemize}
\item The bracket through L bits: [Σ\_\allowbreak{}halted 2\textasciicircum{}\{−|t|\}, that + Σ\_\allowbreak{}open 2\textasciicircum{}\{−|t|\} + the mass past L, directed].
\item Halt: a normal form. Loop: Brent's watcher. Forever: the spine grows. Halts: simply typed.
\item Codes, de Bruijn: 00M is λM, 01MN is MN, 1\textasciicircum{}k0 is variable k. Terms are enumerated by rank: unrank(length, depth, index) over the counts c(n, k), the variable, then the lambda, then applications by the left part's length.
\item On the engine (\texttt{chaitin\_\allowbreak{}omega.\allowbreak{}cu}, the run by the engine): one normal-order step of every live term is one sweep, a term is one record a token, and there is no step or token budget. Per round:

\begin{itemize}
\item phase one, a thread a term: find the leftmost outermost redex (λM)N; a normal form settles;
\item phase two: each reduct token is a plan (depth D, bound k, body, argument) gathered from its source token by the index. The record program computes (λM)N → M[N], with a body variable i > D + 1 ↦ i − 1, an argument variable i > k ↦ i + D, and every other token copied. It is imprinted at the round's widths;
\item phase three: the same decisions as \texttt{omega\_\allowbreak{}run}, namely halt at a normal form, loop by Brent's watcher, grows forever by the spine proof;
\item settled terms leave as one record; kept terms move into the next round's rooms;
\item a job (≤ 2\textasciicircum{}22 terms of one length) is admitted once fewer than 2\textasciicircum{}21 are live;
\item a term the device cannot hold is parked as outgrown, and a stop leaves the live terms open;
\item the run is one tessera job.
\end{itemize}
\end{itemize}
\item π (\texttt{pi\_\allowbreak{}tower}), BBP at hex position d:

\begin{itemize}
\item term lane i ≤ d, for j ∈ \{1, 4, 5, 6\} with m = 8i + j: r = 16\textasciicircum{}\{d−i\} mod m, by r ← r\textasciicircum{}16·16\textasciicircum{}\{nibble\} mod m one nibble at a time;
\item f\_\allowbreak{}j = ⌊r·2\textasciicircum{}W/m⌋, and F\_\allowbreak{}i = (4f\_\allowbreak{}1 − 2f\_\allowbreak{}4 − f\_\allowbreak{}5 − f\_\allowbreak{}6) mod 2\textasciicircum{}W;
\item tail t < T: ⌊2\textasciicircum{}W/(m·16\textasciicircum{}t)⌋, the same combination; pairs F ← (F\_\allowbreak{}\{2k\} + F\_\allowbreak{}\{2k+1\}) mod 2\textasciicircum{}W;
\item S = ΣF mod 2\textasciicircum{}W, error ≤ 4(N + T) + 1 units, and the leading bits where S ± E agree are certified;
\item cost: d + 1 lanes, linear in d, ⌈(d + 1)/70,656⌉ sweeps on the RTX 3070;
\item neither tower is in it.
\end{itemize}
\item Host only: ψ (\texttt{ka\_\allowbreak{}psi}, the exponent β(r) = (nʳ − 1)/(n − 1) held symbolically) and Goodstein (hereditary trees, ordinals below ε₀).
\end{itemize}

\textbf{E14. The floor the coder answers to} is its own theory, the compression workbook: compression\_\allowbreak{}table.md, C0 to C13.

\textbf{E15. The governing functional} (noise\_\allowbreak{}sieve\_\allowbreak{}tower). F[U] = Π₋₄(⊕\_\allowbreak{}\{∂Φ\_\allowbreak{}new\} LUT(Key(W\_\allowbreak{}local \& L↓))):

\begin{itemize}
\item L↓ is the program;
\item W\_\allowbreak{}local is the atom;
\item \& applies the key's masks;
\item Key is the imprint;
\item LUT is the table step;
\item ⊕ is the cycle's fold;
\item Π₋₄ is the proof on the product.
\end{itemize}

\textbf{E16. The root noise of a box} (the noise detector, \texttt{noise\_\allowbreak{}root\_\allowbreak{}box}).

\begin{itemize}
\item A box B = \hyperref[chap:noise-vector-integration-table]{lo, hi) on each of t, z, y and x: an object's box, which the program names (the cell workbook's). Each term the noise vector table ([noise\_\allowbreak{}vector\_\allowbreak{}integration\_\allowbreak{}table.md}) names as shared keeps one value for each place along its kept axes and shares it along the rest. Rows (row 8) share x; columns (9) share y; planes (10) share y and x; the fixed pattern (18) shares t and z; the stack term (25) shares z.
\item The pattern is p\_\allowbreak{}τ(k) = ⌊the mean of u over the voxels of B at k on τ's kept axes⌋, exact. The residual is r\_\allowbreak{}τ = u − p\_\allowbreak{}τ, spread back along the shared axes, and u = r\_\allowbreak{}τ + p\_\allowbreak{}τ exactly (\texttt{noise\_\allowbreak{}root\_\allowbreak{}return}).
\item The price c(v) is the coder's bits for T(v) (E4) on a lattice of ints. Term τ saves g\_\allowbreak{}τ = c(u\_\allowbreak{}B) − c(r\_\allowbreak{}τ) − c(p\_\allowbreak{}τ), signed.
\item The root is argmax g\_\allowbreak{}τ over the terms with g\_\allowbreak{}τ > 0, the first on a tie, and none when no term saves. The yank keeps the root's residual and its pattern.
\item The price comes in as a callback (\texttt{Noise\allowbreak{}Cost}), since the detector reaches no module but its own; the sim prices through \texttt{tower\_\allowbreak{}lift} and \texttt{compression\_\allowbreak{}encode}.
\item The patterns are estimated here, as floor means over the box. The detector's other passes measure each term's size instead of its pattern.
\item \textbf{Proved} (\texttt{noise\_\allowbreak{}root}, 72 checks, 0 failed, a tessera job). The scene is 16 frames of 16 × 64 × 64 around one standing body of 300 e. The background is 100 e at offset 300, with read variance 4 and the Poisson shot, which gives a voxel an independent variance of 104. The box is all 16 frames × z 2 to 13 × y 16 to 47 × x 16 to 47, 196,608 voxels. Each term is planted alone at variance 64, 256 and 1,024, 16 runs with the null, and the checks hold:

\begin{itemize}
\item the null names no root;
\item in every planted run, the planted term is the root or there is none, never another term;
\item at 1,024 every term is the root;
\item the root's residual and pattern return the box exactly.
\end{itemize}
\item \textbf{Measured}, the bits the root saves, per mille of the box's (the box's own bits run 1,194,939 to 1,362,651). At variance 64, planes save 7 and the fixed pattern 5, and rows, columns and the stack term pay for no pattern (their savings are −5,409, −6,638 and −54,497 bits). At 256: rows 31, columns 30, planes 28 and the fixed pattern 24, while the stack term still pays for none (−17,214 bits). At 1,024: rows 89, columns 80, planes 59, the fixed pattern 50 and the stack term 32.

\begin{itemize}
\item The stack term's pattern is the largest, 16 × 32 × 32 values, against 32 × 32 for the fixed pattern and 16 × 12 for the planes. That it pays last for that reason is a reading; the patterns' own bits were not printed.
\item The noise vector table puts the coherent terms on the kcr set at 1.0\% to 3.0\% of the static background's frame-difference bits (a Gaussian estimate). A real box may hold no root (a reading, still unmeasured).
\end{itemize}
\item \textbf{Audited against the code} (\texttt{noise\_\allowbreak{}root\_\allowbreak{}extent}, \texttt{noise\_\allowbreak{}root\_\allowbreak{}box} and \texttt{noise\_\allowbreak{}root\_\allowbreak{}return}, noise\_\allowbreak{}detector.cu). Each term's shared axes, the floor mean with an equal count of voxels in every pattern cell, the residual and its exact return, the signed saving, the root as the first strict gain over 0, and each pattern priced at its own extent all hold as stated. The code also bounds what E16 leaves open. A box holds at most 2\textasciicircum{}31 voxels, and a price past 2\textasciicircum{}61 bits is refused. With no root, the residual is the box itself and no pattern is left. \texttt{noise\_\allowbreak{}root\_\allowbreak{}return} does not check that its sums fit an int; only the pair \texttt{noise\_\allowbreak{}root\_\allowbreak{}box} leaves is exact.
\item The engine's own price is \texttt{engine\_\allowbreak{}lattice\_\allowbreak{}bits} (\texttt{engine.\allowbreak{}h}): a lattice of ints lifted by \texttt{tower\_\allowbreak{}lift} and coded by \texttt{compression\_\allowbreak{}encode}, the coder's bits, with the \texttt{Noise\allowbreak{}Cost} signature a driver passes. It is built (the engine build compiles it) and not yet run: nothing calls it until a driver part does.
\item Not built: a driver part over the cell workbook's object boxes, and the pattern stored in the crystal beside the residual. The ruling is that e16 is variable and should or into the construction set. Read here as follows; this is a reading, not the ruling's words.

\begin{itemize}
\item An object's box spans a variable number of frames. How a box's span is chosen is not yet written. The boxes the cell program holds are a body's in one frame (the climb machine's extents, which the box history reads).
\item The pattern is stored beside the residual, or held in the construction set instead. The construction set is found where the noise terms are all members of the photon shot set. Either way the crystal's format changes.
\end{itemize}
\end{itemize}

\textbf{E17. The ask} (M24, A18): every reading the engine takes of a part.

\begin{itemize}
\item An ask a = (address, qualifier). Unbound it returns its cost t(a); bound by β it returns bit(a) = [qualifier holds]·[t(a) ≤ β], β read off the baseline of unbound costs and never written.
\item The gate: S = \{c : c(k) = r(k) for every k ∈ K\}. The order of K sets the price and never S.
\item The order: n links, n + 1 a power of two, ask r covers link c where (r + 1) \& (c + 1) has an odd count of ones. (n + 1)·x = 4·Sᵀb − 2·(Σb)·1, from 4·SᵀS − (n + 1)·J = (n + 1)·I.
\item Contention: with the count taken out, P\_\allowbreak{}sy²·(N + 1) > 4·P\_\allowbreak{}yy·P\_\allowbreak{}ss over the N sweep asks, every term an exact integer.
\item The rank: m\_\allowbreak{}c = min\_\allowbreak{}j t\_\allowbreak{}c,j and s = max\_\allowbreak{}c (max\_\allowbreak{}j t\_\allowbreak{}c,j − m\_\allowbreak{}c). The least c₁ is chosen where m\_\allowbreak{}c₂ − m\_\allowbreak{}c₁ > s, and otherwise the survivors are one cost.
\item A form asked of the part: the slot closes a body that answers N where the form takes the way back and 1 where it falls through.
\end{itemize}

\textbf{Tied to the code since the first draft:} the tree's adjacency and forest (\texttt{max\_\allowbreak{}tree.\allowbreak{}cu}), the Rice coder (\texttt{compression.\allowbreak{}cu}), the climb (\texttt{climb\_\allowbreak{}machine.\allowbreak{}cu}), the refused lane (\texttt{cycle\_\allowbreak{}record\_\allowbreak{}kernel.\allowbreak{}cu}), Ω's engine rounds (\texttt{chaitin\_\allowbreak{}omega.\allowbreak{}cu}) and the NTT's padding (\texttt{shift\_\allowbreak{}agreement.\allowbreak{}cu}). Nothing listed is left untied; an audit against the code is pending.

\section{\texorpdfstring{The table}{The table}}

\tablerecord{\textbf{M1. Exact numbers}}
\begin{description}
\item[{the algebra it holds to}] Every quantity is an exact integer or an exact rational held as a pair; nothing is rounded and no float is formed.
\item[{does}] \texttt{Anchor\allowbreak{}Exact\allowbreak{}Integer} (orior's) is the integer, in \texttt{src/\allowbreak{}cu/\allowbreak{}types/\allowbreak{}integers/\allowbreak{}exact\_\allowbreak{}integer.\allowbreak{}h} and the \texttt{exact\_\allowbreak{}integer\_\allowbreak{}*.\allowbreak{}c} beside it. The engine build compiles it from this tree by default (\texttt{cell\_\allowbreak{}tracking/\allowbreak{}build\_\allowbreak{}engine.\allowbreak{}sh} and \texttt{build\_\allowbreak{}driver.\allowbreak{}sh}; \texttt{ANCHOR\_\allowbreak{}EXACT\_\allowbreak{}ROOT} still overrides), and this tree is the source orior's engine is to be copied from (orior). The header carries their full doc text. The width is any power of two from one limb (32 bits) up, with no ceiling, named in limbs (\texttt{ANCHOR\_\allowbreak{}EXACT\_\allowbreak{}LIMBS}) or in bits (\texttt{ANCHOR\_\allowbreak{}EXACT\_\allowbreak{}BITS}). The default is 128 limbs, 4,096 bits. Static asserts refuse a width that is not a power of two, is below 32 bits, or where the two names disagree. The declared floor \texttt{ANCHOR\_\allowbreak{}EXACT\_\allowbreak{}DIGITS} stays 1,024 digits. A width below 4,096 bits declares its own, and one that does not is refused at compile time. Width-sized working copies sit on the stack up to \texttt{ANCHOR\_\allowbreak{}EXACT\_\allowbreak{}STACK\_\allowbreak{}LIMBS} (4,096 limbs) and come from the heap past it. A copy the heap cannot hold returns \texttt{ANCHOR\_\allowbreak{}EXACT\_\allowbreak{}WILL\_\allowbreak{}NOT\_\allowbreak{}FIT}, and no stack limits a width. Multiplication is a ladder: long multiplication, then Karatsuba from 32 limbs (\texttt{ANCHOR\_\allowbreak{}EXACT\_\allowbreak{}KARATSUBA\_\allowbreak{}LIMBS}), then the Schönhage–Strassen transform once the shorter factor reaches 8,192 limbs (\texttt{ANCHOR\_\allowbreak{}EXACT\_\allowbreak{}TRANSFORM\_\allowbreak{}LIMBS}). The transform works modulo 2\textasciicircum{}n + 1, where every twiddle is a shift, and nothing is rounded. An integer cost model picks each split's depth, and a rung whose workspace cannot be held steps down to the one below. \texttt{anchor\_\allowbreak{}exact\_\allowbreak{}divide} is Knuth's algorithm D in base 2\textasciicircum{}32. Once the divisor and the quotient both reach 8,192 limbs (\texttt{ANCHOR\_\allowbreak{}EXACT\_\allowbreak{}NEWTON\_\allowbreak{}LIMBS}), it switches to Newton's reciprocal: the reciprocal is grown at half precision recursively, taken one Newton step and corrected exactly, and the quotient is then one product on the ladder. Either way the quotient rounds toward zero and the remainder carries the numerator's sign: numerator = quotient · divisor + remainder. \texttt{anchor\_\allowbreak{}exact\_\allowbreak{}divide\_\allowbreak{}exact} is for a division known to leave nothing. The divisor's twos shift out, and the rest multiplies by the inverse of the odd part modulo 2\textasciicircum{}(32 · limbs), then is masked. The inverse comes from Newton's step x(2 − dx), 2 − t being t's two's complement plus two, and each step works only to the precision it doubles to. The quotient is multiplied back, and a remainder is refused as \texttt{ANCHOR\_\allowbreak{}EXACT\_\allowbreak{}NOT\_\allowbreak{}EXACT}. \texttt{anchor\_\allowbreak{}exact\_\allowbreak{}gcd} is Lehmer's (Knuth's algorithm L); a bit-at-a-time binary gcd stalls the 4,194,304-bit test. A zero divisor returns \texttt{ANCHOR\_\allowbreak{}EXACT\_\allowbreak{}BY\_\allowbreak{}ZERO}. \texttt{anchor\_\allowbreak{}exact\_\allowbreak{}multiply\_\allowbreak{}transform} and \texttt{anchor\_\allowbreak{}exact\_\allowbreak{}divide\_\allowbreak{}newton} run their rungs at any size, for tests and timing. The sims' rationals (\texttt{sim\_\allowbreak{}rational.\allowbreak{}h}) use the gcd and the exact division to stay in lowest terms at any width. The key machine divides too, as four record operations (M10). \texttt{decimal\_\allowbreak{}double} expands a decimal string, digits × 10\textasciicircum{}ex, into the exact leading bits of its binary value in the candidate's full width, with its binary exponent. It says whether the value fits and terminates, and how many bits it needs when it does not (A0).
\item[{wants}] Every decimal the machine is handed (a spacing, a position, a cosine) enters through \texttt{decimal\_\allowbreak{}double}, and geometry is exact.
\item[{tried, and what it gave}] \textbf{measured} (the scratchpad \texttt{decimal\_\allowbreak{}double\_\allowbreak{}test}, a checker in exact arithmetic): 0 failures. 1e2200 needs 5,109 bits and reports fits = 0 at 4,096; the machine says “grow this” and does not round. \textbf{Proved}, the ladder and Newton's division (\texttt{test/\allowbreak{}engine/\allowbreak{}arithmetic/\allowbreak{}no\_\allowbreak{}rounding/\allowbreak{}exact\_\allowbreak{}transform\_\allowbreak{}test}, 4 checks, 0 failed at each of 128 limbs, 4,096 limbs and 4,194,304 bits). The ladder and the transform alone each equal a long multiplication kept in the test, on 24, 54 and 63 products: balanced and unbalanced, keyed and all ones, across every rung boundary up to half the width. Each product divides back to its factors, exactly and with remainder, and the gcd holds a factor. Newton's division rebuilds its numerator, with the remainder below the divisor and signed as the numerator, and equals the dispatched division on 24, 54 and 63 divisions with divisors up to half the width. \textbf{Measured}, the rungs' timings (the ledger's “The exact integer: width, ladder and division”). \textbf{Proved}, the width's guards: 1, 8 and 64 limbs compile with declared floors (9, 76 and 616 digits), and 8 limbs with no floor is refused at compile time. \textbf{Proved}, the division on its own (\texttt{test/\allowbreak{}engine/\allowbreak{}arithmetic/\allowbreak{}no\_\allowbreak{}rounding/\allowbreak{}exact\_\allowbreak{}divide\_\allowbreak{}test}, 9 checks, 0 failed, at 128 limbs, 4,096 bits). Long division rebuilds its numerator, with the remainder below the divisor and the signs of truncation, on 200,000 trials of edge-shaped limbs up to 8 wide and on 4,000 trials at widths drawn up to 127 limbs. It also holds on a case built to run the add-back step (an instrumented copy confirmed the branch runs). Exact division returns the quotient of every product on 20,000 trials, odd and even divisors, and refuses one past each product. The gcd divides both values, holds their shared factor and leaves coprime cofactors on 4,000 trials, and it equals Euclid's gcd run on the long division on 4,000 more; gcd(0, x) = |x| and gcd(0, 0) = 0. A zero divisor is refused. With lowest terms, \texttt{ka\_\allowbreak{}psi} is still 152 of 152 and \texttt{ask\_\allowbreak{}state} 33 of 33 on the merged source.
\item[{status}] built; the engine build takes it from this tree
\item[{next}] Read the DICOM spacings and positions through it (M9). orior's engine copied from this tree (item 7 below).
\end{description}

\tablerecord{\textbf{M2. The residual operator}}
\begin{description}
\item[{the algebra it holds to}] R = 2\textasciicircum{}\{|b|\} L\_\allowbreak{}s − L\_\allowbreak{}\{s+b\} over one binomial ladder L\_\allowbreak{}n = [1 1]\textasciicircum{}n ∗ I, separable per axis (A1). The kernel sums to zero exactly: a constant added to I leaves R unchanged.
\item[{does}] \texttt{residual\_\allowbreak{}program} builds the four steps; \texttt{keymath} imprints one key, \texttt{key\_\allowbreak{}schedule} lays it out, \texttt{cycle} runs it in one launch. \texttt{unit\_\allowbreak{}sweep} runs the same as centered [1 2 1] steps per axis in shared memory, working only the live limbs. \texttt{engine\_\allowbreak{}residual} runs either or both; both proved compares every lane, and \textbf{a disagreement is a logic error}: the result is withheld. The orders come in with the request; every step is error-wired (M12). The unit sweep also takes its input as limb planes of a declared width (\texttt{device\_\allowbreak{}planes}, \texttt{input\_\allowbreak{}bits}). A value past that width is refused as a request error, and the output's limbs follow from the width and the orders. \texttt{engine\_\allowbreak{}residual\_\allowbreak{}planes} hands it host planes and returns the residual with its limbs. The 16-bit volume path and the key are unchanged.
\item[{wants}] Orders from the data, per request, with nothing rounded: a per-axis coherence reading gives an exact integer period P, and order = period. Parity by the algebra instead of a blanket guard: an odd background order misaligns the two terms by half a voxel and is refused; an odd smooth order shifts both terms by the same half voxel and is held with that offset recorded (Doug's ruling (a)). Every (s, b) from one sweep that keeps its ladder levels (A1).
\item[{tried, and what it gave}] The key against the step-by-step residual: \textbf{proved}, 0 of 10,485,760,000 lanes differ; 31 ms a frame against 24.5 for the passes (\textbf{measured}). The unit sweeps against the key: \textbf{proved}, 0 lanes differ, and about half the key's time a frame (\textbf{measured}; numbers in the cell workbook's on\_\allowbreak{}the\_\allowbreak{}engine.md, since they were taken on a program's samples). The planes (\texttt{test/\allowbreak{}engine/\allowbreak{}analysis/\allowbreak{}unit\_\allowbreak{}sweep/\allowbreak{}unit\_\allowbreak{}sweep\_\allowbreak{}planes\_\allowbreak{}test}, a tessera job): \textbf{proved}, 336 checks, 0 failed, against the 16-bit volume path by linearity. The readings are built from 16-bit parts at chosen shifts, at 16, 34, 64 and 116 bits in. On every lane their residual equals the parts' residuals shifted and summed. That is 4,357 lanes a case, 4,356 of them not zero, over three sets of orders, \{2,34,34\}/\{6,96,96\} at 10 and 13 limbs among them. Sixteen-bit planes equal the volume path. A constant added moves no lane. A bit past the width refuses, and the sweep gives the same lanes after it. Eight malformed requests refuse.
\item[{status}] proved (exact); the sweeps 2× faster; any input width proved (the planes)
\item[{next}] Ruling (a) at all three guards (\texttt{residual.\allowbreak{}cu}, \texttt{keymath.\allowbreak{}cu}, \texttt{unit\_\allowbreak{}sweep.\allowbreak{}cu}, and \texttt{narrow\_\allowbreak{}bits} at :259), with the half-voxel offset as a field of the result. Taking the orders from the per-axis coherence reading, now built as M18 (\texttt{period\_\allowbreak{}read}, A12); the knee pulls it. The ladder-keeping sweep.
\end{description}

\tablerecord{\textbf{M3. The component tree}}
\begin{description}
\item[{the algebra it holds to}] The dense level code, the components of every upper level set, nested as a tree. Every question about levels is a question about tree nodes: membership at level r is an ancestor query, two voxels share a component at r exactly when r ≤ level(LCA), and every subtree is one range of the DFS order (A2).
\item[{does}] \texttt{max\_\allowbreak{}tree\_\allowbreak{}objects}: the level code, the components at one cut (the level with the most, proved against the forest's own count), the per-component moments (M4). \texttt{max\_\allowbreak{}tree\_\allowbreak{}slide}: the caller's probe voxels, and each recorded level's probe partition, proved against the tree's count at that level; a disagreement is a logic error. \texttt{max\_\allowbreak{}tree\_\allowbreak{}overlap}: two frames' trees at one residual level each, the caller's probes and probe links, and the census of the bipartite support at the lag: one partner, unpaired, mutual, forked, per side. \texttt{max\_\allowbreak{}tree\_\allowbreak{}keep\_\allowbreak{}frames} holds a frame's code for the next. No cell is known to it: probes are points, links are point pairs.
\item[{wants}] Every level at once instead of one cut or sampled cuts: the probe partition at every level from k − 1 LCA levels, the census at every level from one own-node pair table summed over DFS ranges (A2).
\item[{tried, and what it gave}] The slide's counts: \textbf{proved}, 0 of 26,454 differ from the tree. The overlap and census on 6bba\_\allowbreak{}48816121 and 44b6\_\allowbreak{}0113de3b: \textbf{measured} (numbers in the cell table, row 1, since they are graded by a program's key).
\item[{status}] built, error-wired
\item[{next}] The LCA and DFS-range forms (A2), which remove the bisection and the per-level recount.
\end{description}

\tablerecord{\textbf{M4. The moments}}
\begin{description}
\item[{the algebra it holds to}] Per component, exact: the mass, the three index sums, the six second moments. Together these are the harmonics up to ℓ = 2 (ℓ = 0 the mass, ℓ = 1 the sums, ℓ = 2 the moments). Every node's moments are its subtree's own-node moments summed (A3).
\item[{does}] \texttt{max\_\allowbreak{}tree.\allowbreak{}cu} accumulates them on the device for the components at the cut; \texttt{grow\_\allowbreak{}leaves} carries them out (\texttt{ENGINE\_\allowbreak{}BODY\_\allowbreak{}WORDS} 11: peak, mass, 3 sums, 6 moments).
\item[{wants}] Every node's moments at every level, by a prefix over the DFS order (A3).
\item[{tried, and what it gave}] Carried into every record program and the viewer; proofs are the programs' (device equal to host).
\item[{status}] built, error-wired
\item[{next}] A3's prefix, which a level-free program (the slide's chains) needs.
\end{description}

\tablerecord{\textbf{M5. The correlation operator C}}
\begin{description}
\item[{the algebra it holds to}] \texttt{C\_\allowbreak{}\{t,t′\}[a, b](v)} = Σ\_\allowbreak{}x [L\_\allowbreak{}t(x) = a] · [P\_\allowbreak{}\{t′\}(x + v)] · [L\_\allowbreak{}\{t′\}(x + v) = b]: one count of labeled positive voxels under a shift, read by different reductions (A4).
\item[{does}] \texttt{shift\_\allowbreak{}agreement}: the argmax of the unlabeled sum over every v, by an NTT mod 998,244,353, exact because every count is below the prime. \texttt{climb\_\allowbreak{}machine}: a local argmax per label from a start lag, 26 neighbors a step; \texttt{-\allowbreak{}-\allowbreak{}box} reads C at the 27 shifts around every climbed lag; \texttt{-\allowbreak{}-\allowbreak{}core} the nested widths; extents per label. \texttt{body\_\allowbreak{}overlap}: C at one lag.
\item[{wants}] One C per frame pair over a window of shifts, then only reads: sums, local argmax, transpose, box sums by one summed-area table (A4). The window is the only open design choice, and it must come from the request instead of a scale. The search reads against the integrated noise floor (M8): a count at a shift is weighed by how far it departs from what the noise functionals predict there. The climb and the null draws then spend work only where the field departs from the floor.
\item[{tried, and what it gave}] The box (\texttt{-\allowbreak{}-\allowbreak{}box}): \textbf{proved}. C is read at the 27 shifts around every climbed lag and at the drift, for every forward pair and every null draw. A cell meeting more labels than the kernel's table writes its raw pieces and the host merges them, with no cap. At the climbed lag, C sums to the climb's held count on every climber; the climbed lag is the highest of its 27 neighbors on every climber; at the drift, C equals \texttt{body\_\allowbreak{}overlap} on every label (numbers in the cell workbook's on\_\allowbreak{}the\_\allowbreak{}engine.md, since they were taken on a program's sample). The core per width: \textbf{proved} at width 0 against C at the climbed lag, 35,755 of 35,755 (numbers in the cell table, row 13).
\item[{status}] built; the box proved
\item[{next}] Make C the one pass, with the window from the request.
\end{description}

\tablerecord{\textbf{M6. The exact marginal}}
\begin{description}
\item[{the algebra it holds to}] Z = the sum over injective arrangements of sources onto targets (or a null each) of the product of their integer weights: a permanent. It factors exactly over the connected components of the bipartite support graph: a source's probability over its component is its probability over everything (A5).
\item[{does}] \texttt{marginal}: Z and Z\_\allowbreak{}\{b→·\} for each source over a local set, device equal to the host, compared by cross-multiplying. The weights are the caller's.
\item[{wants}] The set taken as the source's whole support component, and that makes the probability the global one exactly; the component's size reported instead of truncated.
\item[{tried, and what it gave}] Timings and outcomes are the program's (cell table, row 7).
\item[{status}] built
\item[{next}] The component as the set (A5); the limits below become measured costs.
\end{description}

\tablerecord{\textbf{M7. One-to-one matching}}
\begin{description}
\item[{the algebra it holds to}] The heaviest one-to-one set on exact integer costs.
\item[{does}] \texttt{heaviest\_\allowbreak{}matching}.
\item[{wants}] none
\item[{tried, and what it gave}] Used by the cell program's \texttt{-\allowbreak{}-\allowbreak{}print-\allowbreak{}match} (cell table, row 7).
\item[{status}] built
\item[{next}] none
\end{description}

\tablerecord{\textbf{M8. The entropy history}}
\begin{description}
\item[{the algebra it holds to}] Per voxel and bit, the flips in each window of transitions; the parity of a bit's total flips equals that bit of the net change, at every voxel (A7). The window-overlap cloud folds into the same pass.
\item[{does}] \texttt{entropy\_\allowbreak{}history}: the flips, the parity proof, the cloud, one pass on the device, error-wired (broken parity is a logic error). A sample of fewer than two frames is a request error (no transition to count).
\item[{wants}] The window from the request or the data, not the machine (see the audit). The noise categories as functionals (the plan in compression\_\allowbreak{}table.md): fixed pattern from the bits that never flip, shot and read from the photon transfer curve (Σ (I\_\allowbreak{}t − I\_\allowbreak{}\{t−1\})² against Σ I\_\allowbreak{}t per signal level, exact sums), quantization as one step of the lane, and the spatially correlated terms per row, column and place. Integrated, they give every voxel's floor and so every departure from it.
\item[{tried, and what it gave}] \textbf{proved} exact on 2,000 voxels × 9 windows. Measurements of what it shows are the cell program's (cell table, row 12). Its identity by spatial null permutation (\texttt{knf\_\allowbreak{}identity}, M19; A13, the two crystals): \textbf{proved}, the host's walk equals the engine's history at all 262,144 voxels of a 64³ lattice over 16 windows, and of 8,192 single flipped bits 1,500 leave it unchanged, each where the window rule says: the history is many-to-one; the 48 exact cube motions carry it and keep its cloud.
\item[{status}] built, error-wired
\item[{next}] Take the window as a request field. The pairwise test between sections: one body's departure curve against another's (A13, the two crystals).
\end{description}

\tablerecord{\textbf{M9. The files}}
\begin{description}
\item[{the algebra it holds to}] Every file the machine writes is proved by reading it back: pixels voxel for voxel against a second read of the source, and every part by its seal (M13; the crystal's seal, the others' CRC-64 until they are sealed too). Integer lifting is exactly invertible (A8).
\item[{does}] \textbf{\texttt{.\allowbreak{}kcr}}, the crystal (the Kolmogorov crystal; \texttt{.\allowbreak{}iapx} crystals are not read, sets are re-ingested). Layout, read front to back: head (12 words), the seal (6 roots, every lane node, every chunk leaf), offsets and stream, the deflated side bytes, the member tables and names, EOF. \textbf{dicom} is a source kind: one sample is one SeriesUID, slices ordered by an exact key from their positions (M1), signed pixels sign-extended and lifted by 2\textasciicircum{}15 into the u16 lane (a slice whose raw words differ keeps its pixel bytes), every header kept whole in the side bytes (deflated by our own LZ77 and Huffman coder). The tower (\texttt{tower\_\allowbreak{}lift}, integer lifting to the floors, of the lanes or of a lattice of ints, E4), \texttt{compression} (a Rice coder per block), \texttt{krep} (the container that writes and reads the \texttt{.\allowbreak{}kcr}), and the rebuild proof on every ingest. Their device buffers are held as pools (\texttt{device\_\allowbreak{}pool}). Each pool is one allocation whose slices are laid at 256-byte steps, and the pool is rounded up to the 2 MiB page once. The tower's pool holds its coefficients, scratch, flag and mismatch count, and its edge table grows apart. Compression holds two pools: its chunk bits, chunk offsets and scan scratch, and its stream. \texttt{tower\_\allowbreak{}hold\_\allowbreak{}bytes} and \texttt{compression\_\allowbreak{}hold\_\allowbreak{}bytes} give a job those bytes before any device work. The stream is sized by the values once they are measured, and it is left to the kept peak. \texttt{.\allowbreak{}oapx}: the entropy history. \texttt{.\allowbreak{}bapx}: the bodies' table. \texttt{flatten}: the set as one magnitude vector per label in one file. The ingest readers: zarr v2/v3/N5 with every codec, TIFF, HDF5, npy/npz, NRRD, NIfTI, \texttt{.\allowbreak{}stack}. \textbf{zip}: its own module (\texttt{src/\allowbreak{}cu/\allowbreak{}includes/\allowbreak{}codecs/\allowbreak{}zip}), the directory read whole with no size cap, members read and CRC-32 checked, an index by containing folder, and one archive held across samples. \texttt{npy} reads through it. Every part is error-wired: an IO failure carries errno, bad contents in the machine's own files are logic errors, a malformed source is a request error.
\item[{wants}] The export: the crystal written back out as any source format (the dicom zip byte for byte from the side bytes, npy, NIfTI, TIFF, zarr). Names for the other files (flattened, history, bodies, key) are open with Doug. The floor: the residue coded against the integrated noise functionals (M8) over exact counts. Exact arithmetic coding ends within 2 bits a stream of the length those functionals give, and enumerative coding in the engine's exact integers hits it exactly; a fixed-width range coder rounds each interval and is only the fast approximation. That length is the most any coder can take once the functionals are the real categories, with the model's own bits counted as a two-part code (F5 and F3′ in compression\_\allowbreak{}table.md). It replaces the Rice coder's chosen block, k and escape.
\item[{tried, and what it gave}] \texttt{.\allowbreak{}kcr} on the RSNA knee: \textbf{proved}, 1 × 34 × 960 × 960 unsigned, 25,218,496 bytes (40.2\% of raw) sealed \texttt{7ac2cb89…12eb2a}; 1 × 24 × 640 × 640 signed, 11,319,416 bytes (57.5\%) sealed \texttt{3e70b8cb…8d9966}; rebuilt voxel for voxel, pixel for pixel and node for node. The RSNA test\_\allowbreak{}series: 15 of 15 held, 192,020,272 of 599,191,552 bytes (32.0\%), set root \texttt{98b25a42…d7af979}, prove 696 ms with every seal node held. Deflate: within about 1\% of zlib-9 on DICOM headers, round-tripped through our inflate and zlib. A missing sample reports module 10, site 178, status 2 (errno ENOENT). \texttt{flatten}: \textbf{measured}, slowed by its residual slab sized to half the free device memory, with the same bytes out one frame at a time (numbers in the cell workbook's on\_\allowbreak{}the\_\allowbreak{}engine.md, since they were taken on a program's sample). Codec variants: the residue table in the compression workbook's ledger. The pools (\texttt{test/\allowbreak{}engine/\allowbreak{}runtime/\allowbreak{}device\_\allowbreak{}pool/\allowbreak{}device\_\allowbreak{}pool\_\allowbreak{}test}, 28 checks, 0 failed): the tower's four buffers for 1,100,000 voxels cost 10,485,760 bytes as one pool, the plan to the byte, and 14,680,064 as four allocations, read through tessera's counter (\textbf{measured}). Compression's two pools raise \texttt{root\_\allowbreak{}universal}'s peak from 154,304,512 to 158,498,816 bytes, two pages (\textbf{measured}). That each pool takes a page its small buffers would otherwise share is a reading and is not measured.
\item[{status}] proved: the field, and the RSNA knee as \texttt{.\allowbreak{}kcr}
\item[{next}] The floor: F2 per floor on 44b6\_\allowbreak{}0113de3b against the Rice coder (compression\_\allowbreak{}table.md). The 25 are re-ingested as \texttt{.\allowbreak{}kcr} (item 37 of the cell workbook's build\_\allowbreak{}plan.md). The prose and identifiers still say iapx (\texttt{engine\_\allowbreak{}iapx\_\allowbreak{}head}, \texttt{entry\_\allowbreak{}iapx\_\allowbreak{}encode}, the cell README) where the file is \texttt{.\allowbreak{}kcr} (\texttt{ENTRY\_\allowbreak{}CRYSTAL\_\allowbreak{}SUFFIX}, engine.cu). The export. The engine's opens need a long-path form (a set path past 260 characters fails). Flatten: per-request orders per sample (ruled), and a slab not sized as a fraction of memory.
\end{description}

\tablerecord{\textbf{M10. The record machine}}
\begin{description}
\item[{the algebra it holds to}] A record program is steps over limbs, run for every record in one sweep; truthy verdicts are fields (zero false, nonzero true; AND a product, OR a sum). A step may read a lookup table: a general one-variable function of one register's low bits. The machine's edges need not be linear (A13). The step count is the scheduler's n, held apart from the register file, the binding resource (A13).
\item[{does}] \texttt{keymath} imprints record programs, \texttt{key\_\allowbreak{}schedule} lays them out, \texttt{cycle} runs them on the device (with a host reference); every entry is error-wired. The operations are field, constant, product, sum, difference, ladder, absolute, compare, signed field, and table (\texttt{ENGINE\_\allowbreak{}RECORD\_\allowbreak{}TABLE}, op 10), threaded through the config, \texttt{keymath}, \texttt{key\_\allowbreak{}schedule} and \texttt{cycle}, on the device and the host port. A table maps the low \texttt{index\_\allowbreak{}bits} (at most 32) of one register to an \texttt{out\_\allowbreak{}bits} value, laid out as 2\textasciicircum{}index\_\allowbreak{}bits rows of ⌈out\_\allowbreak{}bits/32⌉ limbs. The exact integer's division is four operations (M1): quotient (\texttt{ENGINE\_\allowbreak{}RECORD\_\allowbreak{}QUOTIENT}, op 11), remainder (12), gcd (13) and exact quotient (14). The quotient rounds toward zero and the remainder carries the numerator's sign. The gcd is never negative. A zero divisor, or an exact quotient that leaves a remainder, refuses the lane. The widths are read at imprint: a quotient is as wide as its numerator, a remainder as the narrower operand and a gcd as the wider. On the device, Knuth's algorithm D divides, the gcd is Euclid's on it and the exact quotient is the multiply and mask. Only a program that divides carries the scratch, 5 · WIDE + 4 limbs a lane. The host's steps are the exact integer library itself. Xor (\texttt{ENGINE\_\allowbreak{}RECORD\_\allowbreak{}XOR}, op 15), and (16) and the two's complement wrap (17, to 4 bits or more), bitwise on each operand's two's complement sign-extended without end. The step count has no bound: each sign sits beside its register, and \texttt{ENGINE\_\allowbreak{}RECORD\_\allowbreak{}LIMBS\_\allowbreak{}MOST} = 256 limbs of registers is the only bound, and floors of steps stack into one program. Register reuse (\texttt{reuse} on the key-schedule and record requests, off by default) frees a register once its last reader has run, and a chain needs only its live registers; off, every step keeps its own register. A sweep whose frame grows the device's stack limit sets it back once it ends (\texttt{cycle\_\allowbreak{}stack\_\allowbreak{}return}), and that returns the local memory the limit reserves.
\item[{wants}] Limits measured, not written in (see the audit). The table's functions chosen without fitting (item 3 of \hyperref[chap:kolmogorov-arnold]{kolmogorov\_\allowbreak{}arnold.md}).
\item[{tried, and what it gave}] Device equal to host on every program run (the programs' proofs). The table (\texttt{test/\allowbreak{}engine/\allowbreak{}compiler/\allowbreak{}cycle/\allowbreak{}record\_\allowbreak{}table\_\allowbreak{}test\_\allowbreak{}compose.\allowbreak{}cu}, run by \texttt{test/\allowbreak{}engine/\allowbreak{}compiler/\allowbreak{}cycle/\allowbreak{}record\_\allowbreak{}table\_\allowbreak{}test.\allowbreak{}sh}): \textbf{proved}, 15 checks, 0 failed, every table filled by running an ordinary-ops program over the whole 16-bit alphabet, and the ops are the oracle. x² (product) and |x − 30000| + 100 (constant, difference, absolute, sum): the table equals the ops on 4,096 lanes, device, host and table alike. Composition, f(x) = 65535 − x and g(y) = |y − 30000|: a table read through a table equals the ops for |35535 − x|, a single table composed on the host (g read through f) equals the two, and the reversed order f(g(x)) differs on some lane: the chain keeps its order. Register reuse: a 21-step additive chain gives the same output from 3 limbs as from 21. Two refusals: an index wider than its source register, and a table index with no table behind it. The division (\texttt{test/\allowbreak{}engine/\allowbreak{}compiler/\allowbreak{}cycle/\allowbreak{}record\_\allowbreak{}divide\_\allowbreak{}test}): \textbf{proved}, 15 checks, 0 failed. The device equals the host word for word on 4,096 lanes of signed 160-bit numerators in the 64-limb register file and on 512 lanes of 2,048-bit numerators in the 256-limb file. On every lane numerator = quotient · divisor + remainder, the remainder is below the divisor and signed as the numerator, the gcd divides both, and the exact quotient returns the factor it was built from. A zero divisor and an inexact division refuse the lane, host and device. 3\textasciicircum{}40 divides exactly by 3\textasciicircum{}20 to 3\textasciicircum{}20. A quotient reading a later step is refused at imprint. The stack return (\textbf{measured} on the RTX 3070): a 7,600-byte frame holds 467,668,992 bytes until the limit is set back. \texttt{test/\allowbreak{}engine/\allowbreak{}compiler/\allowbreak{}cycle/\allowbreak{}record\_\allowbreak{}tower\_\allowbreak{}test} reads the limit at 1,024 bytes a thread before its blocks and after them (210 checks, 0 failed); the peak is unchanged, since the reservation stands during the sweep.
\item[{status}] built, error-wired; the table, register reuse and the division proved (A13); the programs compiled and measured, 12 to 17 times the interpreter's speed (item 9)
\item[{next}] \texttt{cycle\_\allowbreak{}record\_\allowbreak{}run\_\allowbreak{}host} (cycle.c) is not error-wired yet. The stack return is dropped once the compiled programs, the automata, are emitted (item 9). They are emitted for every program of 1,024 steps or fewer, the divisions included; the interpreter still runs the longer ones, and the stack return stays for it until they compile too.
\end{description}

\tablerecord{\textbf{M11. The golden ladder}}
\begin{description}
\item[{the algebra it holds to}] β(n): n's rung among the Fibonacci numbers that fit the word.
\item[{does}] \texttt{golden\_\allowbreak{}bands}, \texttt{band\_\allowbreak{}of}, 92 rungs (every Fibonacci number below 2⁶⁴: a word's width, not a scale).
\item[{wants}] none
\item[{tried, and what it gave}] none
\item[{status}] built
\item[{next}] none
\end{description}

\tablerecord{\textbf{M12. Errors and integrity}}
\begin{description}
\item[{the algebra it holds to}] The only acceptable error is a stateful one: request, resource (a failed CUDA call, allocation or read, with its status or errno) and logic (unforgivable; the result is withheld). First write wins; each request carries its record and the engine keeps its own. The record holds kind, module, site, status, execaddr, evacaddr, imagebase and up to 16 return frames, and the builds emit a \texttt{.\allowbreak{}pdb}, and execaddr − imagebase steps to the line.
\item[{does}] Wired: the engine's own entries, max\_\allowbreak{}tree, flatten, decimal\_\allowbreak{}double, unit\_\allowbreak{}sweep, cycle, keymath, key\_\allowbreak{}schedule, residual, grow, krep, compression, tower, entropy\_\allowbreak{}history, zip, npy, dicom, obsignatio (modules 0 to 17). Wired: cycle's host record run. \texttt{Cycle\allowbreak{}Record\allowbreak{}Host\allowbreak{}Request} (cycle.h) gains a last member, \texttt{Engine\allowbreak{}Error *error}, and \texttt{cycle\_\allowbreak{}record\_\allowbreak{}run\_\allowbreak{}host} (cycle.c, \texttt{CYCLE\_\allowbreak{}HELD}, module 5) refuses a NULL request or error record outright. It fills a request error for a layout or out that is NULL; zero steps, out limbs or members; a member's in, bodies or count missing; a step wider than \texttt{ANCHOR\_\allowbreak{}EXACT\_\allowbreak{}LIMBS}; a lane whose index runs past its member's records; a step the host arm refuses (\texttt{cycle\_\allowbreak{}host\_\allowbreak{}step}); or a result that does not fit its out limbs (\texttt{cycle\_\allowbreak{}host\_\allowbreak{}fits}). It fills a resource error when its working allocation fails. \texttt{engine\_\allowbreak{}record\_\allowbreak{}host} passes its request's error record through and maps \texttt{CYCLE\_\allowbreak{}REFUSED} to \texttt{ENGINE\_\allowbreak{}REFUSED}. Output: no print in any loop; each functional fills its result struct and one print format renders it on demand through \texttt{scriptura} (our own formatting and SWAR memory scans, its memory operations in M17); errors always print. The run log's default place is beside the program (\texttt{engine\_\allowbreak{}program\_\allowbreak{}directory}): \texttt{build/\allowbreak{}logs/\allowbreak{}track\_\allowbreak{}driver.\allowbreak{}log} for the published driver, whatever folder it runs from; \texttt{-\allowbreak{}-\allowbreak{}log} still names any other. A failed ingest names its stage: no place, source unread, crystal unwritten, or not rebuilt.
\item[{wants}] Every module wired: stack, schedule, golden\_\allowbreak{}bands, the other ingest readers and codecs, run\_\allowbreak{}cfg, run\_\allowbreak{}log, cfg\_\allowbreak{}json, relate\_\allowbreak{}frames, double\_\allowbreak{}fields. Three of these are not a mechanical wiring: golden\_\allowbreak{}bands and double\_\allowbreak{}fields have no refusal path at all, and schedule skips where the others would refuse. Wiring any of the three changes what it does, and each is Doug's call. Integrity is M13; device scheduling is M14.
\item[{tried, and what it gave}] Every wiring build checked against a program's run, 0 lanes differ in both modes (numbers in the cell workbook's on\_\allowbreak{}the\_\allowbreak{}engine.md). A series the archive does not hold is refused as no such member (the entry layer asks the archive; a folder search that finds nothing does not count as found); a failed ingest removes the directories it made (the maker counts what it created; it has no walk bound); a failed prove prints no set root; a structurally refused crystal prints as not read; every path buffer is sized by the platform's own limit (\texttt{ENGINE\_\allowbreak{}PATH\_\allowbreak{}ROOM}: 32,768 on Windows with a long-path-aware manifest, \texttt{PATH\_\allowbreak{}MAX} elsewhere), and seven sites refuse a long path and do not truncate it silently; the npy and nrrd headers have no caps; the host exact integer is fitted to the record path's widest step (256 limbs), and host and device accept the same steps. A 395-character crystal path ingests and proves to the same root. The host record run: eight record tests, each a tessera job: gaussian 11 checks, 0 failed; guide 12, 0 failed; table 16, 0 failed; divide 20, 0 failed; bitwise 25, 0 failed; boundary 49, 0 failed; coherence 15, 0 failed; order 18, 0 failed. The guide test's check: an index past its member's records refuses the sweep, and the error is a request error from module cycle. \texttt{record\_\allowbreak{}tower\_\allowbreak{}test}, a tessera job at below-normal priority: 210 checks, 0 failed.
\item[{status}] wiring in progress
\item[{next}] The next unwired module.
\end{description}

\tablerecord{\textbf{M13. The seal}}
\begin{description}
\item[{the algebra it holds to}] Every file, set and state is sealed by a dimensional Merkle DAG of keyed BLAKE3 nodes, one derived key per level, 32 bytes a node, every node stored. A change anywhere changes exactly its path of nodes up to the root, and a mismatch names its place (A10; the theory is \hyperref[chap:obsignatio-seal]{obsignatio\_\allowbreak{}seal.md}).
\item[{does}] \texttt{obsignatio} (module 17), our own BLAKE3, host and device sharing one compression core: \texttt{obsignatio\_\allowbreak{}signum} (any length, extended output), \texttt{obsignatio\_\allowbreak{}many} (equal messages at a stride, level-parallel), \texttt{obsignatio\_\allowbreak{}bits} (bit ranges of a limb stream), \texttt{obsignatio\_\allowbreak{}lanes} (rows → planes → volumes → lanes), and the level keys. The \texttt{.\allowbreak{}kcr} stores its six roots, every lane node and every chunk leaf. Ingest seals from the source lanes and proves the file against them. Prove and load check in order: stored bytes, then the inflated side, then the decoded rows, then the root. The set root seals every sample's root in the order named. The tower's CRC is gone. The seal sits after the head: both placements measure alike on the knee (the same bytes, prove median 76 ms each), and the seal exists before the stream is written. Ingest witnesses twice (keep both): the decoded voxels against the source on the device, then a second, independent read of the source pixel for pixel. A flip in the first read is then not sealed as the truth.
\item[{wants}] Stage B: the history, bodies, flattened, key and schedule files sealed the same way. The universal root over every file the engine touches plus the engine itself (its hash injected at compile time), and the trust anchor that makes it tamper-proof and not only accident-proof (a signature or a published root).
\item[{tried, and what it gave}] The suite (\texttt{obsignatio/\allowbreak{}test}, the official BLAKE3 vectors): \textbf{proved}, 2,182,450 cases, 0 failures (14 levels). That includes every length 0 to 102,400 host against device, every single-bit flip of every vector input, lane locality 9,234 of 9,234 and bit locality 57,216 of 57,216. Five planted bugs (a dropped byte, a wrong merge, a lost carry, no root flag, a stuck counter) each fail it. On a sealed knee crystal: 8 single-bit tamperings each refused and placed (head shape, a root, row t 0 z 8 y 0, a plane node, chunk 1,200, chunk 1,234 from a stream limb, the deflated side, a member name), the clean copy proved. Cost: 5.4\% of the crystal, 2.1\% of raw at 960 columns (16/X of raw from the rows).
\item[{status}] built and proved: the crystal and the set
\item[{next}] Stage B, then the universal root.
\end{description}

\tablerecord{\textbf{M14. The scheduler}}
\begin{description}
\item[{the algebra it holds to}] Jobs from separate processes share one device by memory: admit while the declaration fits the measured headroom, grow and warn when a job outgrows its reservation, keep each request signum's peak (A11; the theory is \hyperref[chap:tessera-scheduler]{tessera\_\allowbreak{}scheduler.md}).
\item[{does}] \texttt{tessera}, next column, arbitrates between processes: the driver and the device sims submit through it. schedule.cu still plans one process against free/3·2, a chosen share.
\item[{wants}] \texttt{tessera}: a daemon per host and device over a local socket (named pipe, Unix socket, systemd socket activation, a docker socket mount); per-pid measurement (NVML on Linux, the PDH GPU Process Memory counter under WDDM, where NVML reports none); dead sweeps and lost and found; over budget asked at admission against the signum's last peak, held, overridable; the budget and every time a request field; the driver submitting through it; service files; a test suite.
\item[{tried, and what it gave}] Measured on this machine: the PDH counter lists each process's dedicated device bytes by pid. NVML per-process memory is unavailable under WDDM by NVIDIA's documentation. The ledger (\texttt{tessera\_\allowbreak{}ledger}: the accounting, backfill with the head kept, one min-heap of deadlines): \textbf{proved}, 143,808 cases, 0 failed, from \texttt{engine/\allowbreak{}daemon/\allowbreak{}test}. The frame (\texttt{tessera\_\allowbreak{}frame.\allowbreak{}c}, 128 little-endian bytes): \textbf{proved}, 100,000 random round trips byte for byte, the layout at fixed places, and 131 single-bit flips of the head, each refused or held exactly as the rule says. The measure (\texttt{tessera\_\allowbreak{}measure.\allowbreak{}c}: NVML loaded at run time, the PDH counter matched on pid and LUID, a process held by a handle or pidfd): \textbf{measured} on the RTX 3070 under WDDM, 8 of 8 cases. This process read 24,576 bytes, then 412,254,208 after allocating 268,435,456 (the rest is its CUDA context); the device's in-use rose by the same 412,070,912; a pid with no device work read 0; a child held before its exit read live, then ended. NVML's total is 655,360 bytes above what the runtime reports (WDDM keeps them), and C is NVML's figure. The client (\texttt{tessera\_\allowbreak{}client.\allowbreak{}c}: connect, spawn the daemon if absent, submit, override, wait, precalc kept, release) and the daemon (\texttt{tessera\_\allowbreak{}daemon.\allowbreak{}c}, linked with the ledger, the measure, the paths, the frame, scriptura and obsignatio) compile clean at /W4. The job test (\texttt{tessera\_\allowbreak{}job\_\allowbreak{}test}, the client against the built daemon, which the client starts through \texttt{daemon\_\allowbreak{}path}): \textbf{proved} on the RTX 3070 under WDDM, 24 checks, 0 failed. Its first four paths, in order. A new signum declaring 64 MiB is admitted with 64 MiB granted; it takes 256 MiB, and its release reports it grew to 412,254,208 bytes, with a peak of 412,254,208, the process's whole dedicated memory, its CUDA context included. The same signum declaring four times that peak is asked, the ask naming the kept peak, and the override admits it on its declaration. Declaring exactly the kept peak is admitted at once. Declaring four times the peak with no answer is held for its 0.4 s holding time, then lost; its ticket names a lost and found directory, and the precalc kept releases it. The fifth part runs in a scratch \texttt{TESSERA\_\allowbreak{}STATE} inside the run's build directory. The lost ticket's seal holds, with “precalc kept” as its body's last line. The daemon ends once idle, and the history it leaves is whole records and the seal (80 bytes here: one record and the 32-byte seal). A byte flipped, a byte cut and the seal stripped are each refused: the submit fails, and no daemon starts. The history restored, the same signum at four times the peak is asked with the kept peak, which proves the history was read. The daemon seals the history and the tickets with \texttt{obsignatio\_\allowbreak{}seal} (keyed BLAKE3, on the host), and it makes no CUDA call and holds no device context (\hyperref[chap:tessera-scheduler]{tessera\_\allowbreak{}scheduler.md}). The daemon the last check starts ended itself within 5 s of the test. The history is loaded before any endpoint exists. A refused history ends the daemon before a client can reach it. A program's driver submits through it (its jobs and the sample they run on are in the cell workbook's on\_\allowbreak{}the\_\allowbreak{}engine.md). \textbf{Measured} there: two proves of one crystal in a row, each granted its declaration of 838,860,800 bytes, peaked at 3,958,566,912 and 3,962,761,216. A second run on the real state gave the same declarations, and the two prove peaks in the other order: identical requests measured 4,194,304 bytes (2\textasciicircum{}22) apart. The second prove was admitted on its declaration, below its kept peak, as A11 says. Until its sweeps grew it, its reservation stood 3,119,706,112 bytes below what its signum is known to need, and another job could have been admitted into that room. The rule reserves what a job asks for, and if it costs less, remembers that for next time. A job that takes more than it declared is grown and warned by the sweeps, as A11 says. The code reserves the larger of the declaration and the kept peak (\texttt{tessera\_\allowbreak{}ledger\_\allowbreak{}wants}), not yet ruled on. Its ledger suite passes the budget's 22 checks, 0 failed, and with it two proves each reserve 3,962,761,216 bytes, the kept peak. \textbf{The standing} (A11). Each job declares on top of what its process already holds. Its standing is read as it asks, and it is admitted on the bytes it has yet to find. The ledger suite checks this on a device of 1,000 bytes with 300 in use (\textbf{proved}, 8 standing cases, 0 failed). A job standing on 200 and declaring 650 wants 850, more than the 700 free, but has only 650 to find. It is admitted with 850 reserved and 50 left. Its peak of 900 is kept, and the same signum standing on 300 and declaring 650 is not held and wants 950. Declaring 950 is held. The job test ran 25 checks, 0 failed, standing on 24,576 bytes before any context and on the 412,254,208 it kept. Clients make their context before they ask: \texttt{sim\_\allowbreak{}job\_\allowbreak{}submit} does, and the driver does too. For its largest sample the driver declares two parts. The first is the lattice rounded to the page, where the part holds one (all but iapx-prove). The second is the tower's pool and compression's chunk pool, where the part lifts or lowers one (all but schedule and fingerprint). The stream and seal buffers are left to the kept peak. The driver's declaration has not run on a real sample. The sims submit too (\texttt{src/\allowbreak{}sims/\allowbreak{}cu/\allowbreak{}sim\_\allowbreak{}job.\allowbreak{}cu}): each device sim is one job, with a signum over its name and arguments, the bytes of its first allocation declared before it makes it, the same three times, and \texttt{TESSERA\_\allowbreak{}OVERRIDE=1} to override. \textbf{Measured} here in a scratch \texttt{TESSERA\_\allowbreak{}STATE}: all six device sims and Ω's engine run are admitted and released, each peak above its declaration (the figures are in M19), and the history is 368 bytes, seven records and the seal. Two peaks differ between runs by 2,097,152 bytes (2\textasciicircum{}21): nbody\_\allowbreak{}lattice reads 175,276,032 against 173,178,880, and fixed\_\allowbreak{}pattern 145,915,904 against 143,818,752. No program in the repo loads the engine DLL, whose only builder is \texttt{build\_\allowbreak{}engine.\allowbreak{}sh}, and no other caller is left to submit. \textbf{Not rerun here:} the Linux build and suite under WSL 2 (gcc 13.3.0, CUDA 13.3), with 0 warnings. Under WSL, NVML reads a process as 0 bytes before and after it allocates 256 MiB, and no pid can be measured from outside there. So on a paravirtual kernel the daemon's measure opens in reported mode and measures no pid (\texttt{tessera\_\allowbreak{}self\_\allowbreak{}paravirtual}). Each admitted job's own process reads its device bytes through dxcore (\texttt{tessera\_\allowbreak{}self\_\allowbreak{}measure}) and sends them every sweep, and the daemon applies each reading as it would a pid measurement. There the measure test gives 9 cases, 0 failed (0 bytes, then 427,819,008 after allocating 256 MiB), and the job test 24 checks, 0 failed. The systemd units are in \texttt{engine/\allowbreak{}daemon/\allowbreak{}service} (\texttt{tessera@.\allowbreak{}socket} and \texttt{tessera@.\allowbreak{}service}). Under WSL, socket activation starts the daemon on a client's connection, and the whole job test passes with systemd starting every daemon, 24 checks, 0 failed. The docker socket mount, also not rerun here: a static client with no CUDA (\texttt{engine/\allowbreak{}daemon/\allowbreak{}test/\allowbreak{}tessera\_\allowbreak{}socket\_\allowbreak{}probe.\allowbreak{}c}) runs in a container under Docker Engine 29.1.3 inside the WSL distro, the host's socket bind-mounted and named by \texttt{TESSERA\_\allowbreak{}RUNTIME}. Its submit reaches the socket-activated daemon on the host, which admits the job, and the job releases. The container had no device mounted and so no bytes to report, and the history held no record, just the seal (32 bytes). How to submit a job and start the daemon: engine/daemon/README.md.
\item[{status}] ledger, frame, measure and the daemon end to end proved on Windows, the history and the tickets sealed; a program's driver and the device sims submit every run, measured; a job stands on its process's bytes and is admitted on what it has yet to find, proved; the code reserves the larger of standing plus declaration and the kept peak, committed and not yet ruled on; Linux built under WSL with each job reporting its own bytes, socket activation and the docker mount reached, none rerun here
\item[{next}] Doug's ruling on reserving the kept peak. The driver's declaration run on a real sample. Native Linux, where NVML measures a pid (this machine has no native NVIDIA Linux driver). A container on WSL reporting its bytes, which needs dxcore mounted into it (\texttt{-\allowbreak{}-\allowbreak{}gpus}). A job from WSL through the Windows daemon. The unsealed history kept aside as \texttt{history.\allowbreak{}unsealed}.
\end{description}

\tablerecord{\textbf{M15. The sift}}
\begin{description}
\item[{the algebra it holds to}] A pattern of p points is found in a lattice by probing anchors, then verifying the whole pattern at the alignments that pass. What the anchors filter is set before the search by the corpus's collision entropy H2 and its coherence (a period P found against its own multiples): k anchors on one residue class collapse to one probe, and a periodic corpus cannot be filtered below 1/P by anchors. The dispatch between the in-order and free-order arms is one exact comparison, 100·total² against 85·distinct·Σcount², which needs 143 bits and so at least 8 limbs.
\item[{does}] \texttt{nbody/\allowbreak{}orior/\allowbreak{}} is orior's sift: the search, the steering, the portable scan and one scan arm per instruction set. There is one engine, orior's, and \texttt{cell\_\allowbreak{}tracking/\allowbreak{}maint/\allowbreak{}build\_\allowbreak{}stamp.\allowbreak{}sh} builds from \texttt{orior/\allowbreak{}src/\allowbreak{}engine}. No copy is left to compare. \texttt{orior\_\allowbreak{}anchors\_\allowbreak{}for} returns 1 where a period is found and 4 otherwise.
\item[{wants}] One source for the sift and the exact integer, taken from orior by a pinned git dependency instead of a copy. An anchor count between 1 and 4 on a partially coherent corpus. The readers (M9) feed it real lattices of 1 to 8 dims.
\item[{tried, and what it gave}] orior's benches, reported in engine/README.md: soundness, 3,421 rows and 465,546 planted occurrences, none refused; the cost per alignment flat over 4 KB to 4 MB; periodic16 off the histogram by exactly 4096 = 2\textasciicircum{}((k−1)·H2), found before the search by coherence; dispatch on flatness alone, 41 of 42 fastest. Built and run here on four configurations: MinGW gcc (host), MSVC with CUDA 13.3, WSL gcc 13 with CUDA, WSL clang 18 with CUDA, all at sm\_\allowbreak{}86 and 256 limbs. Every bench exits 0 on every configuration (bench\_\allowbreak{}lattice is not built under MSVC, which has no C99 \texttt{\_\allowbreak{}Complex}). \texttt{bench\_\allowbreak{}exact}, \texttt{bench\_\allowbreak{}lattice} (3,253 rows, every one \texttt{hold}), \texttt{bench\_\allowbreak{}scaling\_\allowbreak{}reads} and \texttt{bench\_\allowbreak{}coherence} are byte-identical across every configuration that builds them. \texttt{bench\_\allowbreak{}exact\_\allowbreak{}arms} at 4,096: every arm agrees with portable, AVX2 1.11× to 1.31×, CUDA 0.77× to 1.83×; the planted period reads 4,092 against arithmetic's 4,092, and the device itself answers 2,046 against portable's 2,046. \texttt{bench\_\allowbreak{}steer\_\allowbreak{}arms}: 0 checks failed on each (\textbf{proved}, per configuration).
\item[{status}] orior's, the one engine; built and run here on four configurations
\item[{next}] Doug's rulings on the port (canonical copy, n\_\allowbreak{}body or nbody, the error model).
\end{description}

\tablerecord{\textbf{M16. Render}}
\begin{description}
\item[{the algebra it holds to}] A step's state drawn as a sheet or a volume straight from the state, with no export between. The host arm in C and the device arm in CUDA make the same bytes.
\item[{does}] \texttt{render/\allowbreak{}} is orior's renderer: \texttt{anchor\_\allowbreak{}raster\_\allowbreak{}render} (a sheet) and \texttt{anchor\_\allowbreak{}volume\_\allowbreak{}render} (a volume), each preferring the device and falling back to the host. \texttt{view/\allowbreak{}vis\_\allowbreak{}png} is at \texttt{examples/\allowbreak{}00\_\allowbreak{}blob\_\allowbreak{}viz\_\allowbreak{}tools/\allowbreak{}view/\allowbreak{}vis\_\allowbreak{}png/\allowbreak{}} in orior. There is one engine, orior's (as M15), and \texttt{render/\allowbreak{}} is its own: no copy is left to compare. A WGSL raster arm is not built: \texttt{engine\_\allowbreak{}view/\allowbreak{}gpu.\allowbreak{}js} draws meshes and is not an \texttt{anchor\_\allowbreak{}raster} arm. Once a WGSL arm exists it is held byte for byte against the host and device arms, the way \texttt{render\_\allowbreak{}test.\allowbreak{}py} holds the Python, C host and C device arms.
\item[{wants}] One renderer for the engine's state, taken from orior by a pinned dependency.
\item[{tried, and what it gave}] orior's \texttt{bench\_\allowbreak{}raster}: twenty sheet and twenty volume configurations, host against device byte for byte. Here, 20 of 20 \texttt{ok} on each of the four configurations, the device arm present and graded against the host on the three CUDA builds, the host graded alone on MinGW (\textbf{proved}, per configuration).
\item[{status}] orior's, the one engine; built and run here
\item[{next}] As M15.
\end{description}

\tablerecord{\textbf{M17. Memory operations}}
\begin{description}
\item[{the algebra it holds to}] length, find, copy, move up (target above source, overlapping), compare (−1, 0, 1) and fill, each an exact byte operation; every arm answers exactly what a byte loop answers. Buffers are aligned and terminated by the caller (Doug's rule). An arm may therefore read an aligned block past the end and mask it.
\item[{does}] \texttt{scriptura}'s arms, one file an instruction set, chosen once at first use; below 32 bytes a call goes straight to portable. \texttt{portable}: 64-bit SWAR, overlapping head and tail words to 32 bytes, 32-byte blocks and an overlapping last block above. \texttt{avx2}: four blocks a branch (\texttt{vpminub} for length, OR for find, AND for compare), overlapping tails. \texttt{avx2-\allowbreak{}erms}: the same, with \texttt{rep movsb} and \texttt{rep stosb} for copy and fill from 16 KiB. \texttt{avx512-\allowbreak{}unrun} (masked loads and stores), \texttt{neon-\allowbreak{}unrun}, \texttt{sve-\allowbreak{}unrun} (predicated).
\item[{wants}] Every arm graded on its own hardware (AVX-512 and SVE have none here; NEON has the Pi 5). The small sizes at the C runtime's speed.
\item[{tried, and what it gave}] The grade (\texttt{test/\allowbreak{}engine/\allowbreak{}runtime/\allowbreak{}scriptura/\allowbreak{}scriptura\_\allowbreak{}arm\_\allowbreak{}test.\allowbreak{}c}, every operation over 483 lengths from 0 to 4,160, the terminator, match or difference at four places): portable and avx2 agree with a byte loop on MinGW gcc 16, MSVC /W4, WSL gcc 13 and WSL clang 18, with 0 warnings (\textbf{proved}, per configuration). NEON and SVE build and link on aarch64 gcc 13 and clang 18 and emit \texttt{cmeq}, \texttt{shrn}; \texttt{whilelo}, \texttt{ld1b}, \texttt{brkb}, \texttt{cntp}; AVX-512 emits \texttt{zmm} masked byte compares and moves. Against the C runtime (\texttt{bench/\allowbreak{}bench\_\allowbreak{}scriptura.\allowbreak{}c}, time-stamp ticks, best of 7, aligned buffers, an i7-5960X with ERMS), compare runs 0.31 to 0.77× memcmp on MinGW, and length and find 0.25 to 0.31× the MinGW runtime at 1 to 16 KiB and 0.93 to 1.25× glibc from 4 KiB. Copy and fill at 256 KiB and 1 MiB run 1.13 to 1.40× every runtime; string moves from 16 KiB bring them to 1.00×, and from 4 KiB they lose (1.28×): 16 KiB is the measured crossover on this part. Portable's short path is a leaf with no frame: overlapping words, the long path kept out of line (\textbf{measured}).
\item[{status}] built; the x86 arms proved, the aarch64 arms built and emitting
\item[{next}] The small sizes against glibc (16 to 64 bytes at 1.06× to 1.7×, one call through the arm table); the string crossover measured per part; the NEON arm run on the Pi.
\end{description}

\tablerecord{\textbf{M18. The period reading}}
\begin{description}
\item[{the algebra it holds to}] Per axis of a lattice, the agreement count at every lag 1 to ⌊n/2⌋ over one set of start positions, and every lag shares one denominator. A candidate is a local peak: the agreement at p and at 2p each stands above both neighbor lags, and the height is the smaller rise. The floor is a drawn null per axis instead of chance: for each axis, every line along it is shuffled on its own: that axis loses its order while every other axis keeps its structure, and each shuffle reads its strongest height. P is the fundamental, the smallest candidate whose height stands strictly above every one of them, or above a set-wide top the program passes. Every comparison is an exact cross-multiply (A12).
\item[{does}] \texttt{period\_\allowbreak{}read} and \texttt{period\_\allowbreak{}draw} (src/cu/engine/analysis/period, module 19; its own base module, it returns P and never sets an order, and the program decides which order P sets): a device lattice of u16 lanes, rank 1 to 8. The request names \texttt{draws}, the number of shuffles; \texttt{content}, the lattice's content signum, from which the shuffles are seeded; \texttt{null\_\allowbreak{}top}, one given top per axis for the set-wide band; and optionally \texttt{band} with \texttt{band\_\allowbreak{}room}, at least rank · draws. It takes draws or a given top, never both: no draws without a top, or a top with draws of its own, is a request error. The caller passes the signum, for example the sample's lane root from obsignatio, because no module reaches another. Each draw shuffles each axis once (\texttt{period\_\allowbreak{}line\_\allowbreak{}shuffle\_\allowbreak{}kernel}): every line along the axis is permuted alone by a keyed Fisher–Yates, keyed by the content, the draw, the axis, the line and the step; the four keys come from the signum's four words mixed with the counter draw · rank + axis. A draw's histogram must equal the original's; a mismatch is a logic error. Per axis it returns P (0 when none), the candidate (P itself, or the strongest peak when none clears), the agreement at the candidate and at twice it with the higher neighbor of each (\texttt{agreement\_\allowbreak{}beside\_\allowbreak{}candidate}, \texttt{agreement\_\allowbreak{}beside\_\allowbreak{}double}), the pairs per lag, the exact \texttt{margin} (the height over the pairs), and the band's \texttt{band\_\allowbreak{}count} (the draws that reached a peak), \texttt{band\_\allowbreak{}bottom} and \texttt{band\_\allowbreak{}top} (both the given top when one is passed). It also returns N, Σc² as \texttt{collisions} (reported, not a floor), the band's heights sorted into \texttt{band} when it is given, and optionally every lag's count. \texttt{period\_\allowbreak{}draw(request, draw, heights)} shuffles each axis once along its own lines, keyed by the content, the draw number and the axis, and returns each axis's strongest height (0 over Q where no peak is reached); the program keeps the top per axis across its series and passes the tops as \texttt{null\_\allowbreak{}top}. The counts come from two kernels (a histogram and one agreement row per lag); the choice is made on the host. A volume of 2\textasciicircum{}32 voxels or more is a request error, because N² must fit 64 bits. The engine builds with it; no entry calls it yet.
\item[{wants}] The knee's orders read from it (the knee program: s = b = P, an odd P a request error, P = 0 order 0), with one band for the set: \texttt{period\_\allowbreak{}draw} across its series, the top per axis passed as \texttt{null\_\allowbreak{}top}. Flatten taking per-sample orders from it.
\item[{tried, and what it gave}] The test (\texttt{test/\allowbreak{}engine/\allowbreak{}analysis/\allowbreak{}period/\allowbreak{}period\_\allowbreak{}test\_\allowbreak{}main.\allowbreak{}cu}, run by \texttt{test/\allowbreak{}engine/\allowbreak{}analysis/\allowbreak{}period/\allowbreak{}period\_\allowbreak{}test.\allowbreak{}sh}): \textbf{proved}, 700 checks, 0 failed. Every lag's device count equals a host loop. P, the candidate and its height equal an independent host reference of the same rule (the smallest peak above the top; the strongest peak when none clears). Every band is sorted, and its ends are its bottom and top. A second read of the same content gives an identical reading and band. \texttt{period\_\allowbreak{}draw} for draws 0 to 7 reproduces \texttt{period\_\allowbreak{}read}'s band exactly, and the band's top given as \texttt{null\_\allowbreak{}top} judges every axis as the band did. Nine refusals each give a request error from module 19: no lanes, rank 0, rank 9, an extent of 0, 2\textasciicircum{}32 voxels, agreement room short, no draws, band room short, a given top with draws of its own. Planted periods still read exactly and clear their bands under the per-axis null: 7 on one axis of 4,096, height 2,048/2,048 against a band top of 39/2,048; on 24 × 48 × 60, 5 in z (height 34,560/34,560, band top 109/34,560), none in y (no draw reached a peak), and 12 in x (height 34,560/34,560, band top 64/34,560). z is 24 deep, not 20, because the peak rule needs 2P + 1 ≤ ⌊n/2⌋ and 20 slices allow P ≤ 4. \textbf{Measured} at 8 draws: a smooth ramp with noise (x + y + 2z plus uniform 0 to 63), 16 volumes of 24 × 64 × 64, a period on 3 of 48 axes (6.3\%); on one such lattice, y's strongest peak, 15, has height 14/49,152 against a band top of 35/49,152 and is refused, and x reads 5, height 24/49,152 against a band top of 16/49,152. Uniform random u16, 64 volumes of 32 × 64 × 64: a period on 5 of 192 axes (2.6\%), below the 1/(draws + 1) = 1/9 a permutation test allows (A12); the caller sets the false-period rate with draws (99 draws, 1\%). \textbf{Measured} by \texttt{period\_\allowbreak{}power} (M19), a period 5 planted along x of 32 × 64 × 64 under the camera law, 599 checks, 0 failed: at 8 draws x reads 5 on 32 of 32 at 64, 128 and 1,024 e, and on 31 with 1 multiple at 256 e; at 19 draws, on 32 of 32 from 64 e up. With the plant at 1,024 e the unplanted z and y read a period on 5 of 64 at 8 draws (against 64/9) and on 2 of 64 at 19 draws (against 64/20) (A12). A per-axis null draw is costly, and the band is set-wide (A12). Under the per-axis null (\textbf{measured}, 34 × 960 × 960, RTX 3070): \texttt{period\_\allowbreak{}draw} 455.9 ms a call, and 450.1 ms a marginal draw in \texttt{period\_\allowbreak{}read}; one band for the set holds the more (A12). Three rules do not hold and are not used: the strongest peak reads a multiple of P instead of P; the margin against the non-multiple lags reads smoothness as a period; and the chance floor Σc²/N² cannot tell structure from a rearrangement, since Σc² is unchanged by one.
\item[{status}] built, proved against its reference; the local-peak rule, the fundamental above a per-axis drawn null or a set-wide top
\item[{next}] The knee's set-wide band from \texttt{period\_\allowbreak{}draw}; per-sample orders in flatten; the knee pulls it.
\end{description}

\tablerecord{\textbf{M19. The sims}}
\begin{description}
\item[{the algebra it holds to}] A sim builds a lattice whose answer is known and grades a measurement against it, on the device, in exact integers. The camera law (\texttt{sim\_\allowbreak{}camera.\allowbreak{}h}, the same in every sim): electrons S = the background + a ramp · x + an optional planted wave by phase + the brightness of every body whose ellipsoid holds the voxel, each body an integer center, velocity and semi-axes with a born and an ended frame and a parent. Shot noise is \texttt{sim\_\allowbreak{}poisson\_\allowbreak{}four\_\allowbreak{}cumulants}: each expected electron adds 0, 1, 2 or 4 with chances 9, 8, 6 and 1 in 24. The count's first four cumulants equal Poisson's exactly, and it parts from Poisson at the fifth. \texttt{noise\_\allowbreak{}terms} reads it at 16 electrons as κ3 15,925 and κ4 14,018 thousandths, against 16,000 and 14,000. The symmetric law \texttt{SIM\_\allowbreak{}SHOT\_\allowbreak{}SYMMETRIC} is Binomial(4S, ½) − S, with Poisson's mean and variance but symmetric, and it reads 74 and −9,282. Under this law, floor\_\allowbreak{}match, floor\_\allowbreak{}track, nbody\_\allowbreak{}lattice, noise\_\allowbreak{}floor, period\_\allowbreak{}power and root\_\allowbreak{}universal pass every check. knf\_\allowbreak{}identity passes its 23 checks too. It ran with \texttt{TESSERA\_\allowbreak{}OVERRIDE=1}, since it declares 404,752,392 bytes over its kept peak of 370,311,168, and the lattice's entangled entropy reads 44.170 times the free null's mean. Read noise is Binomial(4r², ½) − 2r². The value is the offset + the fixed pattern at the voxel + the gain · the electrons + the read noise. Every draw is a keyed counter draw (splitmix64 of the key mixed with the counter), and nothing is a float.
\item[{does}] \texttt{engine/\allowbreak{}sims/\allowbreak{}}: eighteen sims, each an exact-integer program (on the device, except \texttt{ask\_\allowbreak{}state}, \texttt{ka\_\allowbreak{}psi} and \texttt{goodstein}, which run on the host, and \texttt{pi\_\allowbreak{}tower}, whose turn runs on the host and whose BBP digits run on the engine's record machine) that \texttt{bash src/\allowbreak{}sims/\allowbreak{}run.\allowbreak{}sh <sim>} builds and runs, exiting 0 only when every check holds. \texttt{nbody\_\allowbreak{}lattice}: moving and dividing bodies with their truth (the mass and the three coordinate sums a frame a body), and \texttt{-\allowbreak{}-\allowbreak{}out <dir>} for \texttt{lattice.\allowbreak{}npy} (u16, T, Z, Y, X) and \texttt{truth.\allowbreak{}tsv}. \texttt{noise\_\allowbreak{}floor}: the Kolmogorov mock (C11), the linear complexity of every bit plane, the photon transfer curve (C3, C7) with the motion removed by the truth and by the data alone, and the neighbor coherence of the frame differences (C5, C6). \texttt{period\_\allowbreak{}power}: \texttt{period\_\allowbreak{}read} (M18) on a planted period from no plant to a strong one, at 8 and 19 draws. \texttt{fixed\_\allowbreak{}pattern} and \texttt{classify\_\allowbreak{}reject\_\allowbreak{}recover}: orior's ART-4-005 and ART-4-007 on the device, the phase dispersion ratio exact against a drawn null band. \texttt{root\_\allowbreak{}universal}: the tower with lookup edges between its floors, on camera-law volumes (M20). \texttt{ask\_\allowbreak{}state}: a qubit carried exactly as the answers of a complete ask, on the host in exact arithmetic (M21). \texttt{ka\_\allowbreak{}psi}: the superposition theorem's inner function ψ, Sprecher's and Köppen's, held exactly to depth 60 with the exponent never expanded, and the separation of its inner sums (A15; \hyperref[chap:kolmogorov-arnold]{kolmogorov\_\allowbreak{}arnold.md}). \texttt{chaitin\_\allowbreak{}omega}: Chaitin's Ω for Tromp's binary lambda calculus, where a closed term halts when it has a normal form, bracketed between two exact dyadics. Every closed term through L bits is run by normal-order reduction on the device, one term a thread (\texttt{omega\_\allowbreak{}kernel}), or on host threads with \texttt{cpu}. The device takes the host's decisions in the host's order under budgets of at most 256 steps and tokens, and hands every term that reaches a smaller budget or outgrows its room back to the host, which runs it under the full budgets. The host's run of every term through 30 bits must give the device's fates and busy beavers, record holders included. \texttt{-\allowbreak{}-\allowbreak{}ledger <path>} plans the run as jobs of 2\textasciicircum{}28 terms, each finished job's tally one flushed line, and a rerun runs only the jobs its ledger lacks. A normal form proves a halt, and Brent's watcher proves a loop when a term returns to a state it held. A term is proven to grow forever when a part of its spine, reduced only by head steps at or below its own position, comes back deeper on the spine (\texttt{omega\_\allowbreak{}grows\_\allowbreak{}forever}). A term still unsettled is proven to halt when it has a simple type (\texttt{omega\_\allowbreak{}simply\_\allowbreak{}typed}, Hindley's unification with an occurs check; a typed term normalizes, Tait 1967), and no loop or growth term may type. The rest stay open. The mass past L is counted with directed rounding instead of run. Given L alone, it runs on the engine's record machine (\texttt{omega\_\allowbreak{}engine}), and every term the engine settled through 24 bits is run again by \texttt{omega\_\allowbreak{}run}: a halt must give the same fate at the same step and the same normal form, a loop the same fate, and a term that grew, by the growth proof or parked as outgrown, is not run again (\texttt{chaitin\_\allowbreak{}omega.\allowbreak{}cu}). At the default L of 28 lengths 25 to 28 are not cross-checked, and no default run has finished (item 12). \texttt{knf\_\allowbreak{}identity}: the entropy history (M8) of \texttt{nbody\_\allowbreak{}lattice}'s law in a 64³ cube over 177 frames against spatial null permutations: the one-bit window rule, the 48 exact cube motions, the entangled entropy E against 19 free draws, the departure curves of two inverse nulls (inside b³ tiles, and whole tiles moved) at tiles 1 to 64 and over each body's box, the scene with no bodies, and alike-voxel controls (A13, \hyperref[chap:two-crystals]{two\_\allowbreak{}crystals.md}). \texttt{goodstein}: Goodstein's sequences held as hereditary trees and never expanded, with b\textasciicircum{}L − 1 at a large L held as one run of b − 1 over the exponents 0 to L − 1; with the base read as ω each tree is an ordinal below ε₀: the start 16 is ω\textasciicircum{}ω\textasciicircum{}ω and 65536 is ω↑↑4 (A13, \hyperref[chap:two-crystals]{two\_\allowbreak{}crystals.md}). \texttt{pi\_\allowbreak{}tower} (pi turning at the boundary, the idea behind tower recursion): the rotation by π − 3 on the circle, held as the integer turn nA mod 2\textasciicircum{}P with π bracketed by Machin's formula. Its closest returns are the floors, π's partial quotients. At 2\textasciicircum{}k cells, the step that fills every cell is read from the floors by the three-gap theorem, and each gap's question is a first hit answered by Euclid's descent. Whether the integer turn lands in the real turn's cell is itself a first hit, checked at every resolution. The resolutions come with the request: \texttt{run.\allowbreak{}sh pi\_\allowbreak{}tower -\allowbreak{}-\allowbreak{} n} reads 2\textasciicircum{}n, \texttt{-\allowbreak{}-\allowbreak{} from to} reads 2\textasciicircum{}from to 2\textasciicircum{}to, and no argument reads 2\textasciicircum{}1 to 2\textasciicircum{}100. P is 3n + 64 rounded up to a word and at least 384, and a width that cannot hold 3(P + 32) + 64 bits is refused with the limbs it needs. Its deepest turns run on the engine. The turn at depth n is α's bit n, and the BBP formula (Bailey, Borwein and Plouffe~\cite{src:Bailey-Borwein-and-Plouffe-1997}, Math. Comp. 66, 1997) reads it where it stands. Each term i ≤ d is a lane of the record machine: 16\textasciicircum{}(d − i) mod 8i + j by base-16 powering, each square taken mod 8i + j, then its fraction floored to W bits. A tail program takes the terms past d, and a pair program adds two lanes mod 2\textasciicircum{}W, and rounds of it reduce a sweep on the record machine. keymath sizes every register, the scheduler lays them with reuse, a sweep is one \texttt{cycle\_\allowbreak{}record\_\allowbreak{}run} over as many lanes as the device keeps resident, and tessera admits the job. The error, 4(N + T) + 1 units over N terms and T tail terms, decides the leading bits, and only those are printed. A resolution past 2\textasciicircum{}20, whole or as base\textasciicircum{}exponent, is held as (2, n): P and the exact width a host turn would need print exact, and the engine sums the depth-n terms sweep by sweep, printing the terms done and the rate. Neither tower, T or T⁻¹ (A13), is in it. \texttt{noise\_\allowbreak{}root}: the root noise of an object's box. Each shared term is planted alone at three strengths, and every yank is priced through the tower and compression's coder (E16). \texttt{noise\_\allowbreak{}terms}: each noise vector planted in the camera law and read back (\hyperref[chap:noise-vector-integration-table]{noise\_\allowbreak{}vector\_\allowbreak{}integration\_\allowbreak{}table.md}). \texttt{floor\_\allowbreak{}match}, \texttt{floor\_\allowbreak{}track}, \texttt{pi\_\allowbreak{}plane} and \texttt{omega\_\allowbreak{}computer} are not described here yet. \texttt{SIM\_\allowbreak{}EXACT\_\allowbreak{}LIMBS} sets \texttt{ANCHOR\_\allowbreak{}EXACT\_\allowbreak{}LIMBS} for every object a sim links. Every sim that uses the device is one tessera job (\texttt{sim\_\allowbreak{}job.\allowbreak{}cu}, M14): before its first allocation it declares those bytes, under a signum over its name and arguments, and its admission and release are two of its checks. The files are in engine/README.md.
\item[{wants}] The engine run on \texttt{nbody\_\allowbreak{}lattice}'s lattice and graded against its truth (not yet run). F4 in bits (C1) from a sim's planted σ²: not computed, because log₂ of 2πeσ² is irrational and the sims carry no floats.
\item[{tried, and what it gave}] \textbf{nbody\_\allowbreak{}lattice}, 9 checks, 0 failed (11 with \texttt{-\allowbreak{}-\allowbreak{}out}): 24 frames of 20 × 96 × 96, 16 bodies, 10 founders and 3 divisions at frames 8, 12 and 16. \textbf{Proved:} the device truth equals the host's walk of each body's box; the signal equals the host's on all 4,423,680 lanes, 0 differ; no lane clipped; a noiseless camera renders the signal exactly; a second render is byte-identical. \textbf{Measured}, the camera law's two moments at gain 1, read variance 3, offset 100, fixed pattern 0 to 8: the residuals sum to 3,271 against a spread of the square root of 636,863,096 (between 25,236 and 25,237), and their squares sum to 635,890,703 against the law's 636,863,096, a ratio of 0.99847 (the check allows 1\%). \textbf{noise\_\allowbreak{}floor}, 11 checks, 0 failed. The Kolmogorov mock (C11): 4,096 streams of 512 bits, each from a linear generator of degree 8 to 64 and a drawn seed; Berlekamp–Massey on the device, one thread a stream, puts the complexity within the degree on 4,096 of 4,096 and regenerates each whole stream from its first L bits on 4,096 of 4,096 (\textbf{proved}). The description, 2L bits a stream, is 286,744 of 2,097,152 bits (13.67\%). Keyed draws of the same count read 252 to 262 against n/2 = 256. The camera lattice's bit planes, 8,640 streams of 512 consecutive lanes a plane, random meaning within 16 of n/2 (\textbf{measured}): planes 0 to 8 random on 8,640 of 8,640 (mean L/n 0.5004), plane 9 on 850 (0.0638), plane 10 on 29 (0.0029), planes 11 to 15 on none. The photon transfer curve (C3, C7) on d = I(t + 1) − I(t), 4,239,360 pairs, planted g = 1 and r² = 3 (\textbf{measured}): with the motion removed by the truth (S unchanged, 4,226,574 pairs), d² on S has slope 2.0046 against 2g² = 2 (the check allows 2\%) and intercept 5.643 against 2r² = 6 (the check allows 6 either side). From the data alone, d² on I(t) + I(t + 1) over every pair has slope 7.0460 against g = 1, and the neighbor coherence (C5), (d − d′)² on the four lanes' sum over 4,195,200 pairs, has slope 4.5292 against 1: sharp-edged moving bodies bias the data-only gain, and the neighbor coherence does not remove the bias. That is the measured gap between removing the motion by the truth and removing it from the data. The neighbor correlation of d (C6), Σdd′/Σd², is −71,727/1,161,503,631 on truth-static pairs (the check: within 5σ of 0) and 403,193,327/2,288,244,270 (0.1762) over every pair. \textbf{period\_\allowbreak{}power}, 599 checks, 0 failed under the fundamental and the per-axis null: from 64 e up x reads 5 on every volume at 19 draws, and the unplanted axes' false periods stay within 5σ of 1/(draws + 1) (M18, A12; the counts in \hyperref[chap:ledger]{ledger.md}). \textbf{fixed\_\allowbreak{}pattern} (ART-4-005 ported), 10 checks, 0 failed: frame 8 × 8, 48 frames, swing 20, pattern 40, 8 draws, key 0xF17A. The detector reads period 64, the frame, at a dispersion ratio of 179.565 against its own shuffle's 0.972. The batch and Welford routes (reduced fractions, on the device) agree on 64 of 64 phases, the broken route splits from each on 61, and the residual equals the moving scene on 3,072 of 3,072 values, a reduction of exactly 100\% (\textbf{proved}). The shuffled stack reads period 11 at 2.161, and rejecting at the wrong period, 63, removes 0.2298\%. Against the band top of 5.241 over 8 shuffles, the matched pattern (179.565) is removed at 100\%, no pattern (1.477) is declined and returned untouched, and impulses (2.064) are declined. A static scene feature cannot be told from a fixed pattern and is removed with it: at depth 0, 5, 15 and 30 the reduction is 100.0000\%, 99.9263\%, 99.3370\% and 97.3480\%. It reads the raw integers, not the reference's byte view, and draws by key, not by Python's Mersenne Twister. The numbers are not the reference's; the claims are. \textbf{classify\_\allowbreak{}reject\_\allowbreak{}recover} (ART-4-007 ported), 7 checks, 0 failed. A fixed pattern over a varying subject is found (179.565 against 5.241), the subject exact on 3,072 of 3,072, a reduction of 100\%, 0 of 64 pixels flagged. Incoherent noise of variance 1,600 over a varying subject is declined, but its best ratio, 3.523 at period 3, clears the band top of 3.294: it is declined only because the period is not the frame, a 1-in-9 event under the null. Impulses on a repeat: 0 flagged, exact on 64 of 64. Incoherent noise of variance 144 on a repeat: 55 of 64 flagged, and where the count and the median agree, exact on only 4 of 9 (\textbf{measured}). Agreement between the routes does not make a strict majority, and a strict majority does not make the value exact: the mode and the median can agree on a wrong value. \textbf{root\_\allowbreak{}universal}, 1,268 checks, 0 failed (M20, A14). \textbf{ask\_\allowbreak{}state}, 33 checks, 0 failed (M21, A15). \textbf{ka\_\allowbreak{}psi}, 152 checks, 0 failed (A15): Sprecher's ψ descends, Köppen's is strictly increasing and continuous in both readings, and Braun and Griebel's separation Lemma 3.4 is false at k = 1 and restored by a corrected bound with a narrower ramp (\textbf{proved}). \textbf{chaitin\_\allowbreak{}omega}, 13 checks, 0 failed, at L = 30 on the host (2,048 steps and 2,048 tokens a term). Of 647,463 closed terms, 645,942 halt, 934 are proven to loop, 213 are proven to grow forever, 0 are proven to halt by a type, and 374 stay open (86 out of steps, 288 outgrew the space). 0.12254295… ≤ Ω ≤ 0.12599214…. The lower bound is the halted mass plus the normal forms past L. The upper bound is the halted mass, the open mass and the closed mass past L, and it equals their exact sum. At 200,000 steps and 16,384 tokens (73 s) the same 374 stay open (47 out of steps, 327 outgrew), and the bracket is the same. At L = 33 on the device (14 checks, 0 failed), of 3,910,730 closed terms 3,900,966 halt, 5,663 loop, 1,180 grow forever, 0 type and 2,921 stay open; 0.12312353… ≤ Ω ≤ 0.12599095…. The table equals the host's L = 33 run line for line. At L = 36 (14 checks, 0 failed), 0.12358892… ≤ Ω ≤ 0.12598988…; at 256 steps and tokens (13 checks, 0 failed, the strict busy beaver check dropped under 2,048) the bounds move by 2\textasciicircum{}−28.2 and 2\textasciicircum{}−25.5 against a gap of 2\textasciicircum{}−8.7. At L = 44 and 256 steps and tokens (13 checks, 0 failed), 3,392,860,908 terms, the first typed halt, and 0.12443857… ≤ Ω ≤ 0.12598796…. Every bracket holds 1/8, which puts Ω at 0.00… in binary: 2 bits \textbf{proved}, and the third is not settled. The checks include the closed-term counts against a parse of every code through 22 bits, no loop or growth term typing, the device against the host through 30 bits, and the busiest normal form at each length against BusyBeaverWiki's BB λ (OEIS A333479) at every length through the smaller of L and 33, up to 1812 at 33: at most the published value and equal where every term halted under any budgets, and equal everywhere at 2,048 steps and tokens or more. An open length holds only a lower bound: meeting the published maximum means the run reached the record holder. A known defect, and not fixed: the ledger keeps a line cut inside its last field once a rerun has ended it, and takes it over the whole line after it. At L = 31 that gave a wrong normal form record holder with every check passing (\textbf{measured}; \hyperref[chap:ledger]{ledger.md}). The fates and the bracket cannot be hit. L = 44 at 256 steps and tokens on the idle device (RTX 3070): 28,080 ms, 120.8 million terms a second (\textbf{measured}). \textbf{As tessera jobs} (M14), each device sim runs in a scratch \texttt{TESSERA\_\allowbreak{}STATE}, and the job's two checks are added to its count. Declared bytes, then the measured peak: period\_\allowbreak{}power 262,144 and 145,915,904 (601 checks, 0 failed); nbody\_\allowbreak{}lattice 26,542,080 and 175,276,032 (11, 0 failed); noise\_\allowbreak{}floor 9,142,272 and 187,858,944 (13, 0 failed); root\_\allowbreak{}universal 339,510 and 152,207,360 (1,270, 0 failed); fixed\_\allowbreak{}pattern 141,056 and 145,915,904 (12, 0 failed); classify\_\allowbreak{}reject\_\allowbreak{}recover 118,016 and 145,915,904 (9, 0 failed). \texttt{ask\_\allowbreak{}state} (33, 0 failed) and \texttt{ka\_\allowbreak{}psi} (152, 0 failed) run on the host and submit no job. \textbf{The Ω engine run} at L = 16 (16 checks, 0 failed, declared 30,256 bytes, peak 286,425,088): 226 closed terms, all halt, 0.11664483… ≤ Ω ≤ 0.12600284…, 2 bits \textbf{proved}, BB λ equal to the published value at every length through 16, and every term the engine settled gives, under \texttt{omega\_\allowbreak{}run}, the same fate at the same step and the same normal form. The host path at L = 16 (13 checks, 0 failed) gives the same table line for line. \textbf{knf\_\allowbreak{}identity}, 23 checks, 0 failed (a tessera job, admitted with 404,752,392 bytes reserved, peak 370,311,168). \textbf{Proved:} the host's walk equals the engine's history at all 262,144 voxels; of 8,192 single flipped bits 1,500 leave it unchanged, and the window rule names all 8,192; the 48 exact motions carry it, with E and the cloud unchanged, on 48 of 48; the whole is its tiles plus its seams at all 7 tile sizes, and a rigid move keeps every tile's inside on 133 of 133 draws; alike voxels, 4 of 64 volumes identified, within 5σ of 1/20. \textbf{Measured:} E is 44.626 times the free null's mean, above all 19 draws; the inside and between departure shares cross between tiles of 2 and 4 and sum to 0.590 at 4; each moving body's curve runs 0.112 to 0.196 at 2 up to 0.947 to 1.019 at 32; with no bodies E is −10.675 times the null's mean, not identified (the table in \hyperref[chap:two-crystals]{two\_\allowbreak{}crystals.md}). \textbf{goodstein}, 7 checks, 0 failed. \textbf{Proved:} the bump keeps the tree on 80,256 of 80,256 values below 20,000 at bases 2 to 6; 2↑↑k in hereditary base 2 is ω↑↑k for k = 1 to 4; the starts 1, 2 and 3 reach 0 in 1, 3 and 5 steps, each value the numeric sequence's; the starts 4, 16 and 65536 each run 1,000,000 steps to base 1,000,002 with the ordinal falling on every step, and every value that fits 64 bits equals the numeric sequence's (all 10⁶ steps for 4). \textbf{Measured:} 65536's top run holds 7,625,597,484,984 terms as one block, and all three starts end on the same tail, ω² + ω·8 + 572861 (an observation). \textbf{pi\_\allowbreak{}tower}, 8 checks, 0 failed. \textbf{Proved:} floor(π·2\textasciicircum{}384) is one integer at both ends of Machin's bracket (416 bits, 116 terms), and after the point π begins 0x243F6A8885A308D3; the bracket certifies 109 partial quotients, the first 21 the published [3; 7, 15, 1, 292, …, 1, 84]; the turn's 66 record closest returns to 2\textasciicircum{}112 steps are the convergent denominators q\_\allowbreak{}j in order; the integer turn lands in the real cell on every step searched at 2\textasciicircum{}1 to 2\textasciicircum{}100 cells; the fill step and last cell read from the floors equal a walk of every step at 2\textasciicircum{}1 to 2\textasciicircum{}24; at every resolution the last cell's first touch is the fill step. \textbf{Measured:} 2\textasciicircum{}64 cells fill at step 81,856,329,715,838,968,403 (8.186e19) and 2\textasciicircum{}100 at 1.593e30 (q\_\allowbreak{}57 − 1). The fill takes 1 to 5 times the cells except inside a large floor: up to 71.9 inside 292 (2\textasciicircum{}8), 20.9 inside 84 (2\textasciicircum{}34) and 23.8 inside 99 (2\textasciicircum{}61). The last cell filled is often just behind the start, at 1 − α, 1 − 2α, 1 − 3α (an observation). Then 17 checks, 0 failed, for the arc and balance. \textbf{Proved:} the helix over our disk (radius 1/(2π), rising 1/π a turn), in ℚ(π) at four quarter turns, has κ² = 4π⁶/(π²+1)², τ = 2π²/(π²+1) and τ/κ = 1/π, the tangent of its plane's tilt; its Darboux vector lies along our axis; the torque about the axis is zero and L\_\allowbreak{}z² = 1/(4(π²+1)). From the bracket to 12 places: κ = 5.705134342735, τ = 1.816000663299, tilt 17.656787151412°, L\_\allowbreak{}z = 0.151657235526. In the square billiard from the corner at slope π, unfolded, the segment in cell (i, j) carries L = (−1)\textasciicircum{}(i+j)((j + ½) − π(i + ½)) about the center. Over 2\textasciicircum{}24 columns walked from the integer turn with every floor certain, no wall hit is a corner, no L is zero, and the record near-corners are exactly the q\_\allowbreak{}j. \textbf{Measured:} 52,707,178 floor-wall hits to 16,777,216 side-wall hits; the lead changes hands 35,929,962 times in 69,484,394 segments and is never held longer than 2; the nearest tie, |L| = 1.910e-8, is at π ≈ 5419351/1725033, a convergent. Then 18 checks, 0 failed. At every resolution, H is the first step after the fill in cell 0 and C the first step after H in the fill's last cell, each one first hit, with the certainty check run to C. Both equal a walk at 2\textasciicircum{}1 to 2\textasciicircum{}24. \textbf{Measured:} H lands on a record return or a sum of them, and the wait H − T is either a few steps or about one floor's q. The fill is within 4 steps of the next record return at 64 of 100 resolutions. The second run C − H has the first run's period at 13 of 100, and is shorter at 86 by a floor combination (a record return at 80, c·q\_\allowbreak{}j at 6) and longer at 1 (2\textasciicircum{}6). Then 23 checks, 0 failed, for the residue, the golden helix and any 2\textasciicircum{}n. \textbf{Proved}, the residue: on every floor with q\_\allowbreak{}j ≤ 2\textasciicircum{}21 (j = 0 to 11), π's first q\_\allowbreak{}j steps have the whole parts of the permutation n ↦ n·p\_\allowbreak{}j mod q\_\allowbreak{}j, and at step q\_\allowbreak{}j the walk stands δ\_\allowbreak{}j = q\_\allowbreak{}jπ − P\_\allowbreak{}j off the whole, with q\_\allowbreak{}j|δ\_\allowbreak{}j| < 1 certified. The identities are 3/1, 22/7, 333/106, 355/113, 103993/33102, … 5419351/1725033, with the residues +1.415e-1, −8.851e-3, +8.821e-3, −3.014e-5, +1.912e-5, … −3.820e-8. \textbf{Proved}, the golden helix, with every residue bracketed from the turn on floors q\_\allowbreak{}j ≤ 2\textasciicircum{}112 (j = 0 to 65): the residues flip sign every floor; q\_\allowbreak{}j ≥ F\_\allowbreak{}\{j+1\} and |δ\_\allowbreak{}j| < 1/q\_\allowbreak{}\{j+1\}; and the shrink |δ\_\allowbreak{}j|/|δ\_\allowbreak{}\{j−1\}| lies between ½ and 1 exactly where a\_\allowbreak{}\{j+1\} = 1, each side of 1/φ decided at both ends. \textbf{Measured:} of 65 shrinks, 45 fall below 1/φ and 20 above (every one above on a floor of 1, while 8 floors of 1 fall below), and 28 lie between ½ and 1. The growth a floor, q\_\allowbreak{}65\textasciicircum{}\{1/65\}, is 3.210576 against φ = 1.618033. \textbf{Proved}, any 2\textasciicircum{}n: at 2\textasciicircum{}1000 (\texttt{SIM\_\allowbreak{}EXACT\_\allowbreak{}LIMBS=512}, 16,384 bits; P = 3,072, Machin at 3,104 bits and 864 terms, 888 floors certified), all 23 checks hold. The default width refuses 2\textasciicircum{}1000 by name, “the turn needs an exact width of 9376 bits: run with SIM\_\allowbreak{}EXACT\_\allowbreak{}LIMBS=512”. \textbf{Measured} at 2\textasciicircum{}1000: the fill is step 1.288e302, 12.02 times the cells, on floor 600 (a = 106), and the second run is not the first's period. Then 41 checks, 0 failed, for the engine. \textbf{Proved} on the engine's record machine: π's first 64 bits after the point, 0x243F6A8885A308D3; the hex digits at Bailey's published positions (D. H. Bailey, “The BBP Algorithm for Pi”, 2006, table 1 and its text), 26C65E52CB4593 at 10\textasciicircum{}6, 6C65E52CB459350050E4BB1 at 10\textasciicircum{}6 + 1, 17AF5863EFED8D at 10\textasciicircum{}7 and ECB840E21926EC at 10\textasciicircum{}8; and the engine's cell at 2\textasciicircum{}100, its low 64 bits, equals the exact turn's, 0x5a308d313198a2e0. \textbf{Measured} (RTX 3070, 70,656 lanes a sweep): 10\textasciicircum{}8 terms in 1,416 sweeps and 3.887 s. At 2\textasciicircum{}(2\textasciicircum{}30) cells (hex position 268,435,440), 268,435,441 terms in 3,800 sweeps and 10.638 s, 2.52e7 terms a second, 127 bits certified, the cell's low 64 bits 0x8293097a8232c37f. At 2\textasciicircum{}googol cells (n = 10\textasciicircum{}100, hex position 2.5e99), 46,844,928 of 2.5e99 terms in 663 sweeps and 64.0 s, 7.31e5 terms a second, about 1.08e86 years at that rate; the run was stopped, and no digit at that depth was printed. The job declared 3,413,540 bytes and peaked at 221,413,376: the declaration counts the sim's own buffers, not the record kernel's per-thread file. The runs are in \hyperref[chap:ledger]{ledger.md}. \textbf{noise\_\allowbreak{}root}, 72 checks, 0 failed (a tessera job; E16). \textbf{noise\_\allowbreak{}terms}, 109 checks, 0 failed (\hyperref[chap:noise-vector-integration-table]{noise\_\allowbreak{}vector\_\allowbreak{}integration\_\allowbreak{}table.md}). \textbf{omega\_\allowbreak{}computer}, 5 checks, 0 failed (runs/omega\_\allowbreak{}computer\_\allowbreak{}run.txt).
\item[{status}] built; every sim run recorded here passes every check, the device sims as tessera jobs, Ω runs on the engine, and π's BBP digits run on the engine
\item[{next}] The engine on \texttt{nbody\_\allowbreak{}lattice}'s lattice against its truth.
\end{description}

\tablerecord{\textbf{M20. The root universal}}
\begin{description}
\item[{the algebra it holds to}] The tower is the additive skeleton of the Kolmogorov–Arnold form: sums, collapsed in order onto the floor. The lookup table is its one-variable edge. Laid between the floors, the edges make the collapse an outer sum with nonlinear edges. Inside the reversible stream an edge must be a bijection: it permutes the low bits of a coefficient and passes the high bits. Any other function rides widened, (x, y) → (x, y ⊕ f(x)), the Bennett embedding: a dimensional expansion, and the bound is one of width, not of kind (A14).
\item[{does}] \texttt{tower\_\allowbreak{}lift} and \texttt{tower\_\allowbreak{}lower} take edges. \texttt{Tower\allowbreak{}Edge} \{floor, index\_\allowbreak{}bits, forward\} permutes the low index\_\allowbreak{}bits of every approximation coefficient on floor k through \texttt{forward}, a permutation of [0, 2\textasciicircum{}index\_\allowbreak{}bits). Floor k is the approximation just before level k's lifting; floor = floors is the collapsed floor. Lift lays E0, L0, E1, …, L(N−1), EN, and lower undoes EN⁻¹, L(N−1)⁻¹, …, L0⁻¹, E0⁻¹: the edges on one floor unwind as a stack. index\_\allowbreak{}bits runs 1 to 20. Every edge is checked before any device work (a permutation, the width in range, the floor at most floors) and refused otherwise. With no edges the tower is unchanged. The crystal path in \texttt{engine.\allowbreak{}cu} passes no edges.
\item[{wants}] The edges in the crystal and seal path, with the \texttt{.\allowbreak{}kcr} recording the edge program (Doug's call). Which functions to lay, without fitting (item 3 of \hyperref[chap:kolmogorov-arnold]{kolmogorov\_\allowbreak{}arnold.md}). A fitted edge's price.
\item[{tried, and what it gave}] \texttt{tower\_\allowbreak{}edge\_\allowbreak{}test}, 25 checks, 0 failed (\textbf{proved}; 20 for the edges and 5 for the Bennett edges, A14): one edge, stacked edges on one floor, and edges on every floor including the collapsed one each rebuild the exact lanes; an edge changes the crystal; a non-permutation, a width of 0, a width of 21 and floor 50 are each refused; a clean edge after the refusals still round-trips; a function that is not a bijection rides as a Bennett edge, (x, y) → (x, y ⊕ f(x)), at the price of width: |x| and x ≥ 100 as 16-bit edges round-trip on 64 of 64 lattices and read out exactly, and |x| laid bare is refused (A14). \texttt{root\_\allowbreak{}universal} (M19's camera law), 1,268 checks, 0 failed. \textbf{Proved}, the root stays exact: 48 volumes, 24 of 4 × 12 × 48 × 48 and 24 of 3 × 13 × 37 × 41, 6 floors each, carry 171 random edges (1 to 6 a volume, widths 1 to 12 bits, and one 20-bit edge per extent) through lift, code, wipe of the held coefficients, decode and lower: plain exact on 48 of 48, edged exact on 48 of 48, the crystal changed on 48 of 48. \textbf{Proved}, the fold keeps order: on floor 1, 8-bit edges a then b lay the same crystal as the single table b(a(i)) on 8 of 8; b then a lays a different crystal on 8 of 8; a on floor 1 and b on floor 2 lay a different crystal on 8 of 8: a floor between two edges blocks the fold. \textbf{Measured}, the price of an unfitted edge (one keyed random permutation a floor, coded bits a voxel, 8 volumes of 4 × 12 × 48 × 48): 6.802 with no edge; at widths 2, 4, 8 and 12, floor 0 (110,592 coefficients) 6.810, 6.908, 8.719 and 12.608, floor 3 (72) 6.802, 6.802, 6.804 and 6.817, the collapsed floor 6 (1) 6.802 at every width, every floor 6.808, 6.911, 8.742 and 12.714.
\item[{status}] built and proved in \texttt{tower}; not in the crystal path
\item[{next}] Doug's ruling on the \texttt{.\allowbreak{}kcr} recording the edge program; a fitted edge's price against the unfitted.
\end{description}

\tablerecord{\textbf{M21. The ask and the state}}
\begin{description}
\item[{the algebra it holds to}] The terms. An \textbf{ask} (a measurement, formally a POVM) is finitely many outcomes with positive semidefinite operators E\_\allowbreak{}i summing to the identity. A \textbf{state} ρ (the quantum functional) is a positive semidefinite matrix of trace 1; for a qubit ρ = (I + r·σ)/2, |r| ≤ 1. The \textbf{answer distribution} is p\_\allowbreak{}i = Tr(ρ E\_\allowbreak{}i): “the order falling out of an ask”, not an order imposed. One basis's ask returns the \textbf{magnitude}, the squared size of each amplitude; the \textbf{phase}, the relative angles between the amplitudes, it cannot see, and other asks can. A \textbf{complete ask} (informationally complete) is one whose answers determine the state, one to one; it needs at least d² outcomes. The \textbf{valid set} is the distributions some state gives, a proper subset of all of them. The \textbf{crossing rule} is the fixed linear map from a complete ask's answers to any other ask's; it is not classical total probability, and the difference is interference (A15; the protocol is \hyperref[chap:query-protocol-table]{query\_\allowbreak{}protocol\_\allowbreak{}table.md} and the plan engine\_\allowbreak{}plan.md).
\item[{does}] \texttt{engine/\allowbreak{}sims/\allowbreak{}ask\_\allowbreak{}state}, on the host, no GPU work: a qubit carried exactly as the answers of a complete ask. Rationals on the exact integer, reduced by a word gcd when both parts fit a word and carried exact, unreduced, when they do not. The SIC's numbers in Q(√3), a + b√3, multiplied with (√3)² = 3, each sign decided exactly by comparing a² with 3b². Two asks, both E\_\allowbreak{}i = (I + v\_\allowbreak{}i·σ)/4: a rational tetrahedral ask, v = (±1, ±1, ±1)/2 with an even number of minus signs, and the SIC, v = (±1, ±1, ±1)/√3.
\item[{wants}] Two statements, not built (A15): the sparse form of ψ's terms, and the countable numbers held exactly.
\item[{tried, and what it gave}] \texttt{ask\_\allowbreak{}state}, 33 checks, 0 failed. \textbf{Proved} for each ask: the operators sum to the identity and are positive, and the frame is c·I, and the ask is complete (the rational ask: |v|² = 3/4, not rank one, frame I, r = 4 Σ p\_\allowbreak{}i v\_\allowbreak{}i; the SIC: |v|² = 1, rank one, frame (4/3)I, r = 3 Σ p\_\allowbreak{}i v\_\allowbreak{}i). 25 states, 21 pure (14 on the equator from rational t, ((1 − t²)/(1 + t²), 2t/(1 + t²), 0), and rational points of the sphere) and 4 mixed, go state → answers → state exactly, 25 of 25; the pure lie on the boundary, |r|² = 1, 21 of 21, and the mixed strictly inside, 4 of 4. The crossings onto the +x, +y and +z asks equal (1 + r\_\allowbreak{}a)/2, 75 of 75. The phase falls out of the ask: the 14 equator states all answer ½, ½ to the 0/1 ask, yet the complete ask tells apart all 91 of their pairs. The crossing weights onto +x are 3/2, 3/2, −1/2, −1/2 (the rational ask) and ½ + ½√3, ½ + ½√3, ½ − ½√3, ½ − ½√3 (the SIC): a weight is negative: the crossing is not classical total probability. (1, 0, 0, 0) is refused, no state, and crosses onto +x at 3/2 and at ½ + ½√3, values that are no probabilities. The SIC's answers cross into the rational ask's exactly, 25 of 25. \textbf{Measured:} of the 455 distributions with denominator 12, 55 are states for the rational ask and 87 for the SIC; of the others, 104 and 72 still cross to probabilities on the x, y and z asks, and the valid set is stricter than “the three axis crossings are probabilities” (both counts recounted independently in exact fractions).
\item[{status}] built as a sim; proved exact on the host
\item[{next}] Not set (Doug's).
\end{description}

\tablerecord{\textbf{M22. Exact qubit states}}
\begin{description}
\item[{the algebra it holds to}] A state of n qubits is held with every amplitude an exact number, never a float. The field is Q(√2)[i], the numbers a + b√2 with a and b Gaussian rationals; the entries of X, Y, Z, S, H, CNOT, CZ and the controlled T phase all lie in it. Its rational functions in ω, \texttt{Q(√2)[i](ω)}, hold a phase left as a symbol. Two forms: the dense statevector, 2\textasciicircum{}n amplitudes, and the matrix-product state, a chain of site tensors whose bond after each two-site gate is cut to its exact rank by Gaussian elimination (M = C·F). The boundary-lens expectation ⟨ψ|O|ψ⟩ contracts the chain from both edges inward (A17).
\item[{does}] \texttt{examples/\allowbreak{}qasm/\allowbreak{}}, error module 20, host only: the exact qubit states in C. \texttt{qasm\_\allowbreak{}field\_\allowbreak{}number.\allowbreak{}c} and \texttt{qasm\_\allowbreak{}field\_\allowbreak{}rational.\allowbreak{}c}: rationals and Q(√2)[i] on the exact integer (128 limbs by default); past that width a value refuses and never wraps. \texttt{qasm\_\allowbreak{}dense.\allowbreak{}c}: the dense statevector with X, Y, Z, S, H, CNOT, CZ, the controlled phase and a general 2×2 or 4×4 matrix gate. \texttt{qasm\_\allowbreak{}chain\_\allowbreak{}apply.\allowbreak{}c} and \texttt{qasm\_\allowbreak{}chain\_\allowbreak{}matrix.\allowbreak{}c}: the matrix-product state, run unchanged over either field, and the boundary-lens expectation. \texttt{qasm\_\allowbreak{}symbolic\_\allowbreak{}function.\allowbreak{}c} and \texttt{qasm\_\allowbreak{}symbolic\_\allowbreak{}polynomial.\allowbreak{}c}: \texttt{Q(√2)[i](ω)} in rat\_\allowbreak{}make's canonical form, evaluated at a number. \texttt{qasm\_\allowbreak{}lens\_\allowbreak{}hash.\allowbreak{}c} and \texttt{qasm\_\allowbreak{}lens\_\allowbreak{}rank.\allowbreak{}c}: the reduced-round SHA-256 wedge-field invariants. \texttt{qasm\_\allowbreak{}exact.\allowbreak{}c}: exact fractions, angles and fixed point. The witness root is the engine's seal (\texttt{obsignatio\_\allowbreak{}seal}). The files are in engine/README.md.
\item[{wants}] qasm in the engine library build: \texttt{src/\allowbreak{}build\_\allowbreak{}engine.\allowbreak{}sh} (MODULES and INGEST) and \texttt{src/\allowbreak{}cu/\allowbreak{}CMake\allowbreak{}Lists.\allowbreak{}txt} name no qasm source; it is a module with its own test and not part of the engine library. orior's fixed-point qasm (\texttt{qasm\_\allowbreak{}bitstring.\allowbreak{}cu}, \texttt{build.\allowbreak{}sh} and \texttt{run.\allowbreak{}sh}) sits beside the exact host qasm, whose types are \texttt{Qasm\allowbreak{}Exact*}.
\item[{tried, and what it gave}] \texttt{examples/\allowbreak{}qasm/\allowbreak{}test/\allowbreak{}run.\allowbreak{}sh}, on the host instead of as a tessera job: \textbf{39 checks, 0 failed}. Two are regression checks: the bond at the widest site index reads 0 instead of past the tensors; and w\textasciicircum{}(bits−1) at w = 2 evaluates to 2\textasciicircum{}(bits−1), which an extra power 2\textasciicircum{}bits must not refuse. Every value below is checked by \texttt{qasm\_\allowbreak{}exact\_\allowbreak{}test} (\textbf{proved} for these circuits). Dense: Bell and GHZ3 amplitudes “1/2*sqrt2 + 0 i”, ⟨ψ|ψ⟩ exactly 1; Bell and controlled T, |11⟩ = “1/2 + 1/2 i”; five gates and then their exact inverse return |000⟩ to the bit. Chain: Bell, bond [2], 8 elements; GHZ-6, bonds [2,2,2,2,2], 40 elements; GHZ-4 and controlled T, bonds [2,2,2], 24 elements, ⟨1111|ψ⟩ = “1/2 + 1/2 i”; GHZ-100, widest bond 2, 792 field elements, norm exactly 1; the scrambler (10 qubits, depth 6), bonds [2,3,6,8,8,8,6,3,2], 552 elements, norm exactly 1; a five-qubit circuit and its inverse return |00000⟩ exactly; on 6 qubits the chain and the dense state agree on all 64 amplitudes for the GHZ chain, GHZ and controlled T, and the scrambler at depth 6. Symbolic, lens and refusals in A17. \textbf{Measured}, one run at BelowNormal priority: the dense section 350 µs, GHZ-100 15.9 ms, the scrambler 466 ms, the whole chain section 732 ms, the symbolic section 89 ms (its new check evaluates w\textasciicircum{}4095), the lens 8.6 ms. \textbf{The device code}, \texttt{examples/\allowbreak{}qasm/\allowbreak{}qasm\_\allowbreak{}self\_\allowbreak{}program.\allowbreak{}cu} and \texttt{qasm\_\allowbreak{}self\_\allowbreak{}run.\allowbreak{}cu}, \texttt{qasm\_\allowbreak{}self\_\allowbreak{}run} (\texttt{Qasm\allowbreak{}Self\allowbreak{}Request}, \texttt{Qasm\allowbreak{}Self\allowbreak{}Reading}). A circuit becomes record floors of integer steps, four parts an amplitude (the real rational, the real √2, the imaginary rational, the imaginary √2), each floor scaling D by the lcm G of the gate's entry denominators. It is imprinted and laid with register reuse, and one \texttt{cycle\_\allowbreak{}record\_\allowbreak{}run} launch runs every lane through the whole stack. Host and device records are compared word for word. The suites give qasm\_\allowbreak{}test 86 passed, 0 failed; qasm exact test 39 checks, 0 failed; qasm self test 44 checks, 0 failed; the CLI on bell.qasm exit 3, bitstring 00, and on bernstein\_\allowbreak{}vazirani.qasm exit 0, bitstring 101101001101; qasm run.sh exit 0. Every circuit's device records equal the host's word for word, and every host lane equals the dense exact port. Bell: 50 steps, 20 file limbs, 1 in limb, 3 out limbs. The 4-qubit scrambler (15 gates): 1856 steps, 98 file limbs, 4 in, 42 out. GHZ5 and controlled T: 804 steps, 134 file limbs, 8 in, 29 out. 2\textasciicircum{}200 in a 5-qubit lane is past the register file, and the lay refuses it before any run.
\item[{status}] built; exact on the host, checked by \texttt{qasm\_\allowbreak{}exact\_\allowbreak{}test}; the device code proved by \texttt{qasm\_\allowbreak{}self\_\allowbreak{}test}, 44 checks, 0 failed; qasm not in the engine library build
\item[{next}] orior is this repository instead of an external dependency: the sift is at \texttt{src/\allowbreak{}cu/\allowbreak{}engine/\allowbreak{}nbody/\allowbreak{}orior/\allowbreak{}}, and there is no submodule to track.
\end{description}

\tablerecord{\textbf{M23. The refinement loop}}
\begin{description}
\item[{the algebra it holds to}] Two systems work at the machine's lowest layer. The generator compiles a program from its current understanding of the problem, and hands over the binary or its behavior. The critic takes that program as a raw, noisy object, strips the instructions, paths and structure that carry no signal, and hands back the signal. The generator recompiles to match and hands it back, and the cycle repeats. On this machine (a reading, not a ruling), the programs are record programs, the compiler's automata (M10, item 9 below). A verdict is a field, zero false and nonzero true. Two programs agree when their verdicts are equal on every lane, which the machine checks exactly, as it checks the device against the host.
\item[{does}] Nothing of the loop is built. What it would stand on is built: record programs imprinted, laid and run (M10), each step's width read at imprint, and every program's device run compared with its host run word for word.
\item[{wants}] The loop, with the generator and the critic both on the record machine, and each round kept only when its verdicts equal the last round's on every lane. The engine's own cost, with one best answer: the least cost at the same verdicts. The loop's progress is read as a departure from a null, as the entropy history's identity (M8, M19) and the period (M18) are read against their draws. The outcomes: the program sheds operations toward the fewest cycles the problem needs, near its Kolmogorov complexity; an error or exploit put into the code is flagged by the other system as noise, isolated and patched; and forms no one would write by hand may emerge.
\item[{tried, and what it gave}] none
\item[{status}] wanted, not built
\item[{next}] The cost metric, Doug's choice: from what the imprint already reads (the step count, each step's width, the register file's limbs) or from counted device cycles. Then the critic's first form, and the null its curve is read against.
\end{description}

\tablerecord{\textbf{M24. The ask}}
\begin{description}
\item[{the algebra it holds to}] Every ask is \texttt{[ADDRESS] -\allowbreak{}> (QUALIFIER) -\allowbreak{}> [COST] -\allowbreak{}> BIT} (E17, A18): unbound it returns its cost, bound by β it returns one bit. Candidates pass a gate of exact cases, and the survivors are ranked by cost read against its own spread; the two are never added together. A chain's links are read by asks over the rows of a Hadamard matrix, and contention by the square of how many links an ask covers. The protocol, step by step, is \hyperref[chap:query-protocol-table]{query\_\allowbreak{}protocol\_\allowbreak{}table.md}.
\item[{does}] On the host: \texttt{query\_\allowbreak{}ask} (\texttt{compiler/\allowbreak{}bootstrap/\allowbreak{}query\_\allowbreak{}ask.\allowbreak{}\{h,c\}}) holds, reads equal or advancing, and bounds a run by the clock it finds (16 checks, 0 failed, Q1), and \texttt{query\_\allowbreak{}interface\_\allowbreak{}walk} asks addresses in a child the cell can lose (8 checks, 0 failed). The gate (\texttt{gate\_\allowbreak{}descent.\allowbreak{}c}) keeps exactly what \texttt{chain\_\allowbreak{}build} keeps on every relation the ladder holds (Q4). \texttt{ask\_\allowbreak{}order} solves every link at 3 to 63 links in exact integers (1,231 checks, 0 failed, Q5), and \texttt{ask\_\allowbreak{}links\_\allowbreak{}read} reads contention (Q7). One bit a check takes the fewest asks on every relation (Q15). On the part: \texttt{loop\_\allowbreak{}back\_\allowbreak{}if} is asked of sm\_\allowbreak{}86 (Q16, item 14 below): 8 of 2927 forms come back on every count, every one a \texttt{BRA}, and none is cheaper above its spread. A writing of one instruction is searched for on the part for every precept and every ladder relation, every arrangement of the \texttt{.\allowbreak{}kdm} is written from them, run and read back (Q16), and the cases the descent places stand on the part as on the host (Q4).
\item[{wants}] The run channel made of these asks (\texttt{run\_\allowbreak{}channel.\allowbreak{}h}, Open 2 in the plan): a container that puts a chain's covered links and reads the part's clock around them, so that NVRTC, nvJitLink and the CUDA runtime leave the loop. Every slot a \texttt{.\allowbreak{}krs} writes by hand asked the way \texttt{loop\_\allowbreak{}back\_\allowbreak{}if} is (Q16), and every program held against NVIDIA's compiler (Q17).
\item[{tried, and what it gave}] Slicing by adjacent cuts is refuted (Q6). One combined score and membership by agreement inside a floor are not so (Q12, Q13).
\item[{status}] \textbf{proved} (the ask on the host and in the cell, the gate, the order and its solve, the 8 forms), \textbf{measured} (the contention read, one bit a check, the cost of the 8 forms, the writings and arrangements on the part), \textbf{theory} (the run channel)
\item[{next}] The run channel (Open 2), then Q17's first step: the 16 lanes \texttt{sass\_\allowbreak{}lane\_\allowbreak{}needs} writes, run beside the cubins of their C route.
\end{description}

\section{\texorpdfstring{The machine's algebra}{The machine's algebra}}

Every line is exact. No line assumes a scale: orders, windows, widths and spacings are symbols the request fills.

\subsection{\texorpdfstring{A0. Exact numbers (M1): \texttt{decimal\_\allowbreak{}double}}{A0. Exact numbers (M1): decimal\_double}}

A decimal d = (−1)\textasciicircum{}neg · digits · 10\textasciicircum{}ex is 2\textasciicircum{}\{ex\} · 5\textasciicircum{}\{ex\} · digits. Its binary value's leading W bits are found by multiplying or dividing by the widest power of 5 that fits a limb, chunk by chunk, with the fives stripped from a fraction first: the result is the exact integer ⌊|d| · 2\textasciicircum{}\{−e2\}⌋ with e2 chosen so the integer has W bits, and \texttt{terminates} says whether the rest is zero. When the value needs more than W bits, \texttt{needed\_\allowbreak{}bits} says how many and \texttt{fits} = 0. Nothing is rounded.

\subsection{\texorpdfstring{A1. The residual (M2): \texttt{residual\_\allowbreak{}program}, \texttt{keymath}, \texttt{cycle}, \texttt{unit\_\allowbreak{}sweep}}{A1. The residual (M2): residual\_program, keymath, cycle, unit\_sweep}}

Per axis a, with [1 1] the unit step:

L\_\allowbreak{}n = [1 1]\textasciicircum{}n ∗ I,  B\_\allowbreak{}s ∗ B\_\allowbreak{}b = B\_\allowbreak{}\{s+b\},  R = 2\textasciicircum{}\{|b|\} L\_\allowbreak{}s − L\_\allowbreak{}\{s+b\}

\begin{itemize}
\item \textbf{The two terms are two rungs of one ladder.} B\_\allowbreak{}b applied to B\_\allowbreak{}s I is B\_\allowbreak{}\{s+b\} I, and R is the difference of rungs s and s + b of one binomial ladder of the same image. Every residual for every (s, b) is a difference of two rungs, and one sweep that keeps its rungs gives all of them. Orders per request, and a sweep over candidate periods, cost one pass (theory).
\item \textbf{The kernel sums to zero.} B\_\allowbreak{}s sums to 2\textasciicircum{}\{|s|\}, and 2\textasciicircum{}\{|b|\} B\_\allowbreak{}s and B\_\allowbreak{}\{s+b\} both sum to 2\textasciicircum{}\{|s|+|b|\}. A constant added to I leaves R unchanged: the medium's mean cancels exactly.
\item \textbf{The heat identity.} In centered steps H = [1 2 1] = [1 1]², with m = b/2 per axis:

4\textasciicircum{}m δ − H\textasciicircum{}m = Σ\_\allowbreak{}\{k=0\}\textasciicircum{}\{m−1\} 4\textasciicircum{}\{m−1−k\} H\textasciicircum{}k (4δ − H),  and 4δ − H = [−1 2 −1] = −Δ²

so, per axis, R = Σ\_\allowbreak{}\{k=0\}\textasciicircum{}\{m−1\} 4\textasciicircum{}\{m−1−k\} · (−Δ² S\_\allowbreak{}k), with S\_\allowbreak{}k = H\textasciicircum{}k S and S = H\textasciicircum{}\{s/2\} I. H/4 is one explicit step of the discrete heat equation, and R is 4\textasciicircum{}m times what m heat steps remove from the smoothed field: the residual is exactly a weighted sum of negative discrete Laplacians down the ladder. In three axes the same telescoping runs axis by axis: 2\textasciicircum{}\{G\}δ − B\_\allowbreak{}\{b\_\allowbreak{}z\}⊗B\_\allowbreak{}\{b\_\allowbreak{}y\}⊗B\_\allowbreak{}\{b\_\allowbreak{}x\} = (2\textasciicircum{}\{b\_\allowbreak{}z\}δ − B\_\allowbreak{}\{b\_\allowbreak{}z\})⊗2\textasciicircum{}\{b\_\allowbreak{}y+b\_\allowbreak{}x\}δ + B\_\allowbreak{}\{b\_\allowbreak{}z\}⊗(2\textasciicircum{}\{b\_\allowbreak{}y\}δ − B\_\allowbreak{}\{b\_\allowbreak{}y\})⊗2\textasciicircum{}\{b\_\allowbreak{}x\}δ + B\_\allowbreak{}\{b\_\allowbreak{}z\}⊗B\_\allowbreak{}\{b\_\allowbreak{}y\}⊗(2\textasciicircum{}\{b\_\allowbreak{}x\}δ − B\_\allowbreak{}\{b\_\allowbreak{}x\}) (theory: exact algebra, not built as a pass).
\item \textbf{Parity.} 2δ − [1 1] = [1 −1], a first difference, centered half a voxel off. An odd background order therefore puts the subtracted term half a voxel from the kept one, and an odd smooth order moves both by the same half voxel. That is ruling (a): refuse odd b, hold odd s with its offset.
\end{itemize}

The positive set is P(x) = [R(x) > 0].

\subsection{\texorpdfstring{A2. The component tree (M3): \texttt{max\_\allowbreak{}tree}}{A2. The component tree (M3): max\_tree}}

\begin{itemize}
\item c(x) is the dense rank of R(x) among the admitted voxels. Each tree node n has a level λ(n); ν(x) is x's own node, the node at level c(x) that holds x. The component at level r holding x is the ancestor a of ν(x) with λ(a) ≥ r > λ(parent(a)). Its voxels are exactly those x′ with ν(x′) in a's subtree.
\item Number the nodes in DFS order; a's subtree is one range [in(a), out(a)).
\item \textbf{Probes by LCA.} Two voxels share a component at level r exactly when r ≤ λ(LCA(ν(x), ν(x′))). Sort k probes by in(ν). The partition at every level at once is given by the k − 1 levels λ(LCA(p\_\allowbreak{}i, p\_\allowbreak{}\{i+1\})) of neighbors in that order: at level r, neighbors i, i + 1 are together exactly when r ≤ that level, and a probe whose own level is below r is absent. For u before v in DFS order, λ(LCA(u, v)) is the least level among the parents of the nodes in (u, v]: each of those nodes lies under the LCA, and the LCA's child toward v is one of them. It is not the parent of the lowest-level node in (u, v], since a node in u's own branch can sit lower than the LCA's child toward v. The slide's bisection over levels becomes k − 1 LCA queries (built: \texttt{max\_\allowbreak{}tree\_\allowbreak{}probe\_\allowbreak{}levels}, with \texttt{max\_\allowbreak{}tree\_\allowbreak{}probe\_\allowbreak{}partition} for one level from them; proved against a flood fill).
\item \textbf{The overlap at every level from one count.} For a lag v, the own-node pair table J(v)[n, n′] = |\{ x : ν\_\allowbreak{}t(x) = n, P\_\allowbreak{}\{t+1\}(x + v), ν\_\allowbreak{}\{t+1\}(x + v) = n′ \}|. Then for components A at level r of t and B at level r′ of t + 1:

\texttt{C\_\allowbreak{}\{r,r′\}[A, B](v)} = Σ\_\allowbreak{}\{n ∈ sub(A)\} Σ\_\allowbreak{}\{n′ ∈ sub(B)\} J(v)[n, n′]

a rectangle sum over the two DFS ranges, and one Fenwick sweep of J answers every asked pair of levels. The census (one partner, unpaired, mutual, forked) is the degrees of the support of C at the chosen levels. The tracker's overlap still recounts per sampled cut (built: \texttt{max\_\allowbreak{}tree\_\allowbreak{}pairs} and \texttt{max\_\allowbreak{}tree\_\allowbreak{}overlap\_\allowbreak{}sums}, one Fenwick sweep instead of a summed-area table; proved against direct counts).
\item \textbf{Built}, in \texttt{src/\allowbreak{}cu/\allowbreak{}engine/\allowbreak{}nbody/\allowbreak{}max\_\allowbreak{}tree/\allowbreak{}}: \texttt{max\_\allowbreak{}tree\_\allowbreak{}nodes} (the nodes in DFS order, each with its level as a dense rank, its parent, its subtree's end and the root, and ν(x) for every voxel), \texttt{max\_\allowbreak{}tree\_\allowbreak{}node\_\allowbreak{}at}, \texttt{max\_\allowbreak{}tree\_\allowbreak{}probe\_\allowbreak{}levels}, \texttt{max\_\allowbreak{}tree\_\allowbreak{}probe\_\allowbreak{}partition}, \texttt{max\_\allowbreak{}tree\_\allowbreak{}pairs} and \texttt{max\_\allowbreak{}tree\_\allowbreak{}overlap\_\allowbreak{}sums} (\texttt{max\_\allowbreak{}tree.\allowbreak{}h} :99, :105, :126, :130, :155 and :176). No caller outside the test.
\item \textbf{Proved:} \texttt{max\_\allowbreak{}tree\_\allowbreak{}levels\_\allowbreak{}test} (\texttt{test/\allowbreak{}engine/\allowbreak{}nbody/\allowbreak{}max\_\allowbreak{}tree/\allowbreak{}}, host only, with run.sh beside it), a tessera job: “max\_\allowbreak{}tree levels test: 110 volumes, 18071 checks, 0 failed”. Each piece is checked against an independent 6-connected flood fill of \{rank ≥ r\}: the node structure; the component at every level in both frames; the probe partition at every level, with duplicate and absent probes; C at every pair of levels (r, r′) at random lags in [−1, 1]³, against direct counts; and the refusals. The volumes are small and random, with values in [−2, 5] and [−40, 60], plus values across the 32-bit limb boundary.
\item \textbf{Prior art, cited from knowledge.} The component tree of an image's upper level sets, the max-tree: P. Salembier, A. Oliveras and L. Garrido, “Antiextensive connected operators for image and sequence processing”, IEEE Trans. Image Processing 7(4) (1998) 555–570; L. Najman and M. Couprie, “Building the component tree in quasi-linear time”, IEEE Trans. Image Processing 15(11) (2006) 3531–3539. The LCA and its reduction to a range minimum over the DFS order: D. Harel and R. E. Tarjan, “Fast algorithms for finding nearest common ancestors”, SIAM J. Computing 13(2) (1984) 338–355; M. A. Bender and M. Farach-Colton, “The LCA problem revisited”, LATIN 2000, LNCS 1776, 88–94. The Fenwick sweep: P. M. Fenwick, “A new data structure for cumulative frequency tables”, Software: Practice and Experience 24(3) (1994) 327–336.
\end{itemize}

\subsection{\texorpdfstring{A3. The moments (M4)}{A3. The moments (M4)}}

With own-node moments μ(n) = Σ\_\allowbreak{}\{ν(x)=n\} (1, x, x xᵀ), every node's moments are Σ\_\allowbreak{}\{n′ ∈ sub(n)\} μ(n′): one prefix over the DFS order gives every component at every level. The centered second moment is K = m M − S Sᵀ, carried with its mass, never divided.

\subsection{\texorpdfstring{A4. The correlation operator (M5)}{A4. The correlation operator (M5)}}

For frames t, t′, a shift v, a label a of t and a label b of t′ (or ∅, a positive voxel under no label):

\texttt{C\_\allowbreak{}\{t,t′\}[a, b](v)} = Σ\_\allowbreak{}x [L\_\allowbreak{}t(x) = a] · [P\_\allowbreak{}\{t′\}(x + v)] · [L\_\allowbreak{}\{t′\}(x + v) = b]

With a weight φ(x) ∈ \{1, x, x xᵀ\} at v = 0, t′ = t and b = a, the same sum gives the moments. Its reductions:

\begin{itemize}
\item the argmax over v of the sum over all labels (\texttt{shift\_\allowbreak{}agreement});
\item a local argmax per label from a start lag (\texttt{climb\_\allowbreak{}machine});
\item a transpose: \texttt{C\_\allowbreak{}\{t′,t\}[b, a](−v)} = \texttt{C\_\allowbreak{}\{t,t′\}[a, b](v)};
\item box sums over nested widths around a lag, from one summed-area table over v (the core: each width is the box on top of the one below);
\item the autocorrelation of a box, a tent, over one frame's own C (the contact form).
\end{itemize}

The labels may be any level's components: with A2, C at every pair of levels is a rectangle sum of J. The level is therefore one more axis of the same box, and all of the tree's frame-to-frame reading is C over (v, r, r′), counted once.

\subsection{\texorpdfstring{A5. The marginal (M6): \texttt{marginal}}{A5. The marginal (M6): marginal}}

For sources S and targets T with integer weights ω(s, t) and a null weight ω(s, ∅) per source, an arrangement α maps each source to a target or ∅, no two sources to one target:

Z = Σ\_\allowbreak{}α Π\_\allowbreak{}s ω(s, α(s)),  Z\_\allowbreak{}\{s→t\} the same over α(s) = t,  P(s → t) = Z\_\allowbreak{}\{s→t\} / Z

Z is the permanent of the |S| × (|T| + |S|) matrix [ω(s, t) | diag ω(s, ∅)]. \textbf{It factors over the connected components of the support graph} (s \textasciitilde{} t where ω(s, t) ≠ 0): the only coupling between sources is a shared target, and the sum over all arrangements is the product of the sums over components, and Z\_\allowbreak{}\{s→t\} / Z over s's component is the global ratio exactly. A local set smaller than the component is exact only when it is the component. The cost lives in the component's size; a permanent is hard in general, and the measurement to report is the component's size (theory).

\subsection{\texorpdfstring{A6. Matching (M7): \texttt{heaviest\_\allowbreak{}matching}}{A6. Matching (M7): heaviest\_matching}}

The one-to-one set of largest total weight on exact integer costs; no weight is scaled or compared as a fraction.

\subsection{\texorpdfstring{A7. The entropy history (M8)}{A7. The entropy history (M8)}}

For voxel x, bit j and frames 0 … F − 1, with f\_\allowbreak{}j(x) the number of transitions where bit j flips: f\_\allowbreak{}j(x) ≡ bit j of (I\_\allowbreak{}0(x) ⊕ I\_\allowbreak{}\{F−1\}(x)) (mod 2). A voxel where this fails means the history was not taken whole; it is the pass's own proof. The windows partition the transitions; their width is a symbol (see the audit).

\subsection{\texorpdfstring{A8. The files (M9)}{A8. The files (M9)}}

Integer lifting (predict, then update) takes each step's integer quotient by a power of two, and the inverse step subtracts the identical quotient, computed from the same integers. Each step is therefore a bijection on integers and nothing is lost: the tower rebuilds every voxel. The proof is not the algebra but the read-back: every sample is rebuilt from the file on the device and compared voxel for voxel, and again from a second read of the source pixel for pixel.

\subsection{\texorpdfstring{A9. The golden ladder (M11)}{A9. The golden ladder (M11)}}

β(n) = the number of Fibonacci rungs at or below n after the first, over the rungs that fit 64 bits.

\subsection{\texorpdfstring{A10. The seal (M13): \texttt{obsignatio}}{A10. The seal (M13): obsignatio}}

Each level ℓ keys its nodes: K\_\allowbreak{}ℓ = BLAKE3-derive-key-context(“obsignatio aeterna 2026-09-23 “ ‖ ℓ), H\_\allowbreak{}ℓ(m) = BLAKE3-keyed(K\_\allowbreak{}ℓ, m). For a sample of extent (T, Z, Y, X):

ρ(t, z, y) = H\_\allowbreak{}row(I[t, z, y, ·]),  π(t, z) = H\_\allowbreak{}plane(ρ(t, z, ·)),  υ(t) = H\_\allowbreak{}volume(π(t, ·)),  Λ = H\_\allowbreak{}lanes(υ(·)) κ(c) = H\_\allowbreak{}chunk(LE64(ℓ\_\allowbreak{}c) ‖ bits\_\allowbreak{}c),  Σ = H\_\allowbreak{}stream(offsets ‖ B ‖ κ(·)) σ = H\_\allowbreak{}side(H\_\allowbreak{}side stored(deflated) ‖ H\_\allowbreak{}side inflated(inflated)),  μ = H\_\allowbreak{}members(tables ‖ names) Ω = H\_\allowbreak{}sample(head ‖ Λ ‖ Σ ‖ σ ‖ μ),  Θ = H\_\allowbreak{}set(Ω\_\allowbreak{}1 ‖ … ‖ Ω\_\allowbreak{}n)

\begin{itemize}
\item \textbf{Locality.} A change at (t, z, y, x) changes exactly ρ(t, z, y), π(t, z), υ(t), Λ, Ω and Θ. A change in chunk c changes exactly κ(c), Σ, Ω and Θ. Any other node agrees except with probability 2\textasciicircum{}−256.
\item \textbf{Size is strength.} A k-bit check misses a random corruption with probability 2\textasciicircum{}−k, whatever the function. The DAG's gain over CRC-64 is structure: nonlinearity, locality and composition instead of a smaller word.
\item \textbf{The tree is BLAKE3's.} Pairing each level and carrying the odd node up is BLAKE3's left-heavy tree exactly (its left subtree is the largest power of two of chunks below the whole), and the level-parallel device fold equals the host's stack fold.
\item \textbf{Where it stops.} The tower lifts across all four axes, and a decoded row can be checked only after a whole decode. The stored bytes check alone.
\item \textbf{The shared shape.} The seal folds the lattice x, y, z, t, as the tower's floors lift it and the moments' prefix sums it (A3). The residual, the tree, C and the seal are each one lattice read once with a different combining operator. The seal's operator alone is not associative. Its tree shape is therefore fixed, and that fixed shape makes a mismatch locatable (observation).
\end{itemize}

\subsection{\texorpdfstring{A11. The scheduler (M14): \texttt{tessera}}{A11. The scheduler (M14): tessera}}

The device holds C bytes; U(τ) is in use. Each admitted job j has a declaration d\_\allowbreak{}j, a reservation r\_\allowbreak{}j (from d\_\allowbreak{}j, only growing) and a measured use u\_\allowbreak{}j(τ) by pid:

O(τ) = U(τ) − Σ\_\allowbreak{}j u\_\allowbreak{}j(τ),  H(τ) = C − O(τ) − Σ\_\allowbreak{}j max(r\_\allowbreak{}j, u\_\allowbreak{}j(τ))

\begin{itemize}
\item \textbf{Admission:} k is admitted when d\_\allowbreak{}k ≤ H(τ), and admission alone keeps H ≥ 0. The reservation rule, that what a job asks for is reserved and a lesser actual cost remembered for next time, is not the code's: the code reserves max(d\_\allowbreak{}k, p\_\allowbreak{}σ) instead (\texttt{tessera\_\allowbreak{}ledger\_\allowbreak{}wants}), which is not approved (M14).
\item \textbf{Standing} (the rule is s\_\allowbreak{}k plus d\_\allowbreak{}k). A job's process already holds s\_\allowbreak{}k bytes as it asks: its context and whatever it kept. The daemon reads them by pid, or from the process's own report under WSL. The job wants w\_\allowbreak{}k = max(s\_\allowbreak{}k + d\_\allowbreak{}k, p\_\allowbreak{}σ). U(τ) already counts s\_\allowbreak{}k. Admission, the head's shadow and the backfill weigh only n\_\allowbreak{}k = w\_\allowbreak{}k − s\_\allowbreak{}k, the bytes it has yet to find. An admitted job is reserved w\_\allowbreak{}k and counted as using s\_\allowbreak{}k until its first sweep. The held rule is d\_\allowbreak{}k alone against p\_\allowbreak{}σ.
\item \textbf{Growth:} u\_\allowbreak{}j(τ) > r\_\allowbreak{}j sets r\_\allowbreak{}j = u\_\allowbreak{}j(τ) and warns. Only growth and outside processes drive H negative, and both are measured: the tally equals C − H at every sweep.
\item \textbf{Memory of a request:} p\_\allowbreak{}σ = max over a run of u\_\allowbreak{}j(τ), kept under the request's signum σ. It is a measurement at the sweep's resolution, not a byte-exact constant: identical requests have read peaks 4,194,304 bytes apart (M14). The engine is deterministic in its request and its result, not in the bytes a sweep samples. A new declaration d > p\_\allowbreak{}σ is over budget: asked, held, overridable.
\item \textbf{Liveness:} a ticket is (pid, start time), and a reused pid is a different process. A closed socket or a gone process releases r\_\allowbreak{}j at once.
\item \textbf{Backfill with the head kept:} the history keeps (p\_\allowbreak{}σ, t\_\allowbreak{}σ). The head's shadow S is the least time with H + Σ\_\allowbreak{}\{e\_\allowbreak{}j ≤ S\} max(r\_\allowbreak{}j, u\_\allowbreak{}j) ≥ n\_\allowbreak{}h, and the spare is X = that sum − n\_\allowbreak{}h. A job k behind the head goes now if n\_\allowbreak{}k ≤ H and (now + t\_\allowbreak{}σk ≤ S or n\_\allowbreak{}k ≤ X). A signum with no history is never backfilled.
\item \textbf{Time:} every deadline sits in one min-heap keyed by time to change; the root is the next event, handled in order.
\end{itemize}

\subsection{\texorpdfstring{A12. The period reading (M18): \texttt{period\_\allowbreak{}read}}{A12. The period reading (M18): period\_read}}

On an axis of extent n with the other axes holding M = N/n voxels, take L = ⌊n/2⌋ lags and the start positions whose coordinate on the axis is below n − L. Every lag ℓ ≤ L then compares the same Q = (n − L)·M pairs:

a(ℓ) = \#\{x : I(x) = I(x + ℓ·e)\},  r(ℓ) = a(ℓ)/Q

\begin{itemize}
\item \textbf{The local peak.} A candidate p must stand above both its neighbor lags, and so must 2p:

a(p) > max(a(p − 1), a(p + 1))  and  a(2p) > max(a(2p − 1), a(2p + 1))

Its height is the smaller rise, h(p) = min(a(p) − max(a(p − 1), a(p + 1)), a(2p) − max(a(2p − 1), a(2p + 1))), an integer count over the axis's Q pairs, and its margin is the rational h(p)/Q. The candidates run over 2 ≤ p with 2p + 1 ≤ L, which gives P ≤ ⌊(L − 1)/2⌋: 8 on an axis of 34 slices (L = 17), 5 on 24, 4 on 20.
\item \textbf{The floor: a drawn null per axis.} The request names d draws. For axis a, draw k permutes each line along a on its own by π\_\allowbreak{}k, and the shuffled lattice I∘π\_\allowbreak{}k holds exactly the same values on every line, its histogram is the original's, and every other axis keeps its structure; only a's own order is destroyed. Each draw reads a's strongest height, h\_\allowbreak{}k = max\_\allowbreak{}p h\_\allowbreak{}\{I∘π\_\allowbreak{}k\}(p); the draws that reach a peak make a's band, sorted exactly. A read of rank r makes r·d shuffles, one per axis per draw.
\item \textbf{The choice: the fundamental.} P is the smallest candidate whose height clears the band's top:

P = min \{ p : h(p)/Q > max\_\allowbreak{}k h\_\allowbreak{}k/Q \}

with each comparison a 128-bit cross-multiply. A period P also peaks at 2P and 3P with heights matching P's, and the strongest peak picked a multiple; the smallest peak above the top picks P. An empty band, where no draw reached a peak, or a given top of 0 lets any peak through, and the smallest peak is P. When none clears, P = 0 and the reported candidate is the strongest peak (the largest height, the smallest p on a tie). A lag that makes no peak at p and 2p is never a candidate. Nothing is rounded, and no constant is written in.
\item \textbf{The fundamental leaves the rate alone.} Some peak clears the top exactly when the strongest peak does. An axis therefore reads a period under the fundamental exactly when it did under the strongest peak; the rule changes which period is read, not whether one is (theory).
\item \textbf{The set-wide band.} The knee's ruling is 73,158 draws a series, enough to expect fewer than one false period over the set's 24,386 series × 3 axes = 73,158 axes. \textbf{Measured} (RTX 3070, on 34 × 960 × 960 = 31,334,400 voxels, a smooth ramp x + y + 2z plus uniform noise 0 to 63): the per-axis null makes one shuffle and one count per axis. \texttt{period\_\allowbreak{}draw} takes 455.9 ms a call, the mean of 8 after a warm-up; a call is one count of the original, then a shuffle and a count on each of the 3 axes. \texttt{period\_\allowbreak{}read} takes 1,934.2 ms at 4 draws and 9,136.5 ms at 20: a marginal draw costs 450.1 ms and the fixed cost is 133.7 ms (each as the run prints it from the unrounded times). From the measured draw instead of a measured set run: about 9.1 hours a series and about 223,000 GPU-hours for the 24,386 series. So one band serves the whole set. Its limits: one run on one lattice; host wall-clock with the syncs; and \texttt{period\_\allowbreak{}draw} allocates its buffers on every call, and part of the 455.9 ms may be allocation (not separated). The timing program is a scratch program outside the repo's tests. \texttt{period\_\allowbreak{}draw} shuffles each axis once along its own lines, keyed by the content, the draw number and the axis, and returns each axis's strongest height (0 over Q where no peak is reached). The program calls it across its series, keeps the top per axis, and passes those tops to \texttt{period\_\allowbreak{}read} as \texttt{null\_\allowbreak{}top} with no draws of its own; a given top with draws is a request error. P is the smallest peak above the given top, and a top of 0 lets any peak through. The tops compare exactly across shapes, because the cross-multiply carries each Q. Whether one band is fair across shapes is Doug's call.
\item \textbf{The false-period rate.} With no structure along axis a, each line's order along a is one more arrangement of its values, whatever the other axes carry, and if the draws behave as uniform shuffles the d + 1 heights are exchangeable (theory). The lattice's own strongest height then stands strictly above all d others with probability at most 1/(d + 1), exactly that when no two tie. The caller sets the rate with d: 1/9 at 8 draws, 1\% at 99. \textbf{Measured,} 5 of 192 random axes at 8 draws (2.6\%, M18). With a period 5 planted along x at 1,024 e, the unplanted z and y read a period on 5 of 64 at 8 draws (against 64/9) and 2 of 64 at 19 draws (against 64/20) (\texttt{period\_\allowbreak{}power}, M19): the per-axis null holds the bound when another axis carries a strong period.
\item \textbf{Measured: the fundamental is read} (\texttt{period\_\allowbreak{}power}, M19). A period 5 is planted along x of 32 × 64 × 64 under the camera law, at amplitudes 0 to 1,024 e, 32 volumes each. At 8 draws x reads 5, a multiple (10 or 15), or another lag: 10, 6, 0 at 16 e; 26, 3, 0 at 32 e; 32, 0, 0 at 64 and 128 e; 31, 1, 0 at 256 e; 32, 0, 0 at 1,024 e. At 19 draws it reads 5 on 32 of 32 from 64 e up.
\item \textbf{The strongest peak does not hold and is not used: it reads a multiple.} The candidate is the largest height, the reading lands on a multiple of P, and the multiples' peaks are as high as P's, and the noise picks among them.
\item \textbf{The global shuffle does not hold and is not used: it reads a period where another axis carries one.} One Feistel shuffle of the whole index per draw serves every axis, and an unplanted axis reads a period: the global shuffle destroys the planted axis's structure too, the null's agreement falls below the live lattice's at the unplanted axis's lags, and the exchangeability the false-period rate rests on fails there.
\item \textbf{The margin rule does not hold and is not used.} For a candidate p with 2p ≤ L, the lags that are not multiples of p number B, with A = Σ a(ℓ) over them (the total less the multiples), and the margin m(p) = (a(p) + a(2p))/(2Q) − A/(B·Q) = (B·(a(p) + a(2p)) − 2A) / (2·B·Q) is the largest positive margin the candidate. A smooth lattice agrees more at every short lag than at the long ones, and the rule reads smoothness as a period: a smooth ramp with noise reads P = 2 or 3 on every axis.
\item \textbf{Why chance could not do this.} The old floor was m ≥ Σc²/N². Σc² is unchanged when the values are rearranged, and a floor built from it cannot tell structure from a rearrangement of the same values, and uniform random data passed 63 of 192 axes (33\%). The drawn null keeps the histogram and Σc² exactly and changes only the arrangement, and a period is an arrangement. Σc² is still reported, as \texttt{collisions}.
\item \textbf{The shuffle.} A keyed Fisher–Yates~\cite{src:Fisher} on each line along the axis, one thread a line: for step s from n − 1 down to 1, the element at s swaps with the one at h mod (s + 1), h a 32-bit hash of the line index and s under four keys (\texttt{PERIOD\_\allowbreak{}ROUNDS}). Each key is the content signum's words mixed with c · 4 + round, for the counter c = draw · rank + axis. The seed is the lattice's content, and the same content draws the same band. Superseded: a Feistel network on the whole index, 4 rounds, with cycle-walking to stay below N (the global shuffle).
\end{itemize}

\subsection{\texorpdfstring{A13. The record machine (M10): \texttt{keymath}, \texttt{key\_\allowbreak{}schedule}, \texttt{cycle}}{A13. The record machine (M10): keymath, key\_schedule, cycle}}

A record program is a list of steps s\_\allowbreak{}1, …, s\_\allowbreak{}n over registers of limbs, each step reading earlier registers and writing its own, run for every record in one sweep.

\begin{itemize}
\item \textbf{The steps are n} (the scheduler's steps, with the floors stacked).

\begin{itemize}
\item The step count has no bound. \texttt{ENGINE\_\allowbreak{}RECORD\_\allowbreak{}STEPS\_\allowbreak{}MAX} is gone: it sized a per-thread array holding one sign per step, and each sign now sits beside its register in the file.
\item The binding resource is the file: \texttt{ENGINE\_\allowbreak{}RECORD\_\allowbreak{}LIMBS\_\allowbreak{}MOST} = 256 limbs of live registers.
\item A stack of floors runs as one program in one launch. Each floor is a round that reads the floor below and leaves its state, often wrapped, for the next.
\item The layout frees each register in a single pass, bucketed by its last reader. It frees the same registers at the same steps as the per-step scan it replaced.
\item The algebra of the stack, what it compresses and what it leaves, with the stacked-against-chained measurement (13 to 22 times, records equal): \hyperref[chap:vertical-time-compression]{vertical\_\allowbreak{}time\_\allowbreak{}compression.md}.
\item \textbf{Proved:} 700 floors of a round over four 32-bit words (xor, and, sum, xor with a constant, a 32-bit wrap), 4,204 steps, run in an 8-limb file. On 1,024 lanes the device equals the host word for word, and both equal the same rounds on the CPU's own 64-bit two's complement at every tapped floor (\texttt{test/\allowbreak{}engine/\allowbreak{}compiler/\allowbreak{}cycle/\allowbreak{}record\_\allowbreak{}bitwise\_\allowbreak{}test}).
\end{itemize}
\item \textbf{Xor, and, and the two's complement wrap.}

\begin{itemize}
\item \texttt{ENGINE\_\allowbreak{}RECORD\_\allowbreak{}XOR} and \texttt{\_\allowbreak{}AND} work bitwise on each operand's two's complement, sign-extended without end, one bit wider than the wider operand. The imprint narrows that where it can show a register is never negative. An and with such a register is no wider than it, and an xor of two such is no wider than the wider. The result's sign is read from the operands' signs.
\item \texttt{ENGINE\_\allowbreak{}RECORD\_\allowbreak{}WRAP} takes left modulo 2\textasciicircum{}w, read back signed, with w ≥ 4 in \texttt{right}.
\item \textbf{Proved:}

\begin{itemize}
\item Every output bit is rebuilt from the raw fields one bit at a time, and each sign still holds 64 bits past its written width. This holds on 4,096 lanes in the 64-limb file and 512 in the 256-limb file.
\item Hand-worked answers match.
\item The narrowed widths are exactly the rule's: 32, 20, 20, 32, 41, 41, 32 and 32 over unsigned 32- and 20-bit fields and a signed 40-bit field. All 4,096 lanes equal the oracle past every narrowed width.
\item A wrap below 4 bits is refused at imprint (\texttt{test/\allowbreak{}engine/\allowbreak{}compiler/\allowbreak{}cycle/\allowbreak{}record\_\allowbreak{}bitwise\_\allowbreak{}test}, 24 checks, 0 failed).
\end{itemize}
\end{itemize}
\item \textbf{Register reuse} (opt-in). Register r is live from the step that writes it to its last reader. With reuse on, a liveness free-list returns r to the file once its last reader has run, and a chain needs only the registers live at once instead of one per step. \textbf{Proved} on a 21-step additive chain: the same output from 3 limbs as from 21. Off by default: a program that does not ask keeps one register per step: each register placed at the running total of the limbs before it, the same bump, and a program with no tables keeps the default layout, by construction.
\item \textbf{The lookup table.} A table step reads the low b bits of one register's magnitude, b ≤ 32, and returns T[|x| mod 2\textasciicircum{}b], an \texttt{out\_\allowbreak{}bits}-wide value, from 2\textasciicircum{}b rows of ⌈out\_\allowbreak{}bits/32⌉ limbs. The file holds a magnitude and a sign: a negative x reads the row of |x|, and x and −x read the same row. A register that can be negative is made never negative before it indexes a table. It is a general one-variable function on the lane's alphabet: 2\textasciicircum{}16 entries for a u16 lane, built once by pushing every value through. \textbf{Proved:} filled by running the ordinary operations over the whole 16-bit alphabet, the table equals them on 4,096 lanes, on the device and the host (M10).
\item \textbf{Composition keeps the order.} Reading g through f gives the table (g∘f)[x] = g[f[x]], one table for the chain: the floor is 2-D, but the operations need not be, for a chain of them is collapsed in order onto it. Composition regroups, (h∘g)∘f = h∘(g∘f), but does not reorder: g∘f ≠ f∘g. \textbf{Proved} with f(x) = 65535 − x and g(y) = |y − 30000|. Regrouping: g read through f equals the ops for |35535 − x|, and one table composed on the host equals the two. No reordering: f∘g differs on some lane.
\item \textbf{The two towers} (a second tower over the first one's boundary, inverted, and the two collapse together).

\begin{itemize}
\item Floor division is a record floor: ⌊v / 2\textasciicircum{}k⌋ = EXACT\_\allowbreak{}QUOTIENT(v − AND(v, 2\textasciicircum{}k − 1), 2\textasciicircum{}k), rounding toward −∞ for every integer v, as \texttt{tower\_\allowbreak{}floor\_\allowbreak{}shift} does.
\item With it, the crystal's 5/3 lifting T (A14) and its inverse are sums, differences, constants and that division: stacks of record floors. An operation program F over the crystal and T⁻¹ compose into one stack F ∘ T⁻¹ by the regrouping law.
\item \textbf{Proved:} one 5/3 level over 8 signed 16-bit samples and its inverse, 129 steps in one program and a 12-limb file. On 4,096 lanes the device equals the host word for word, the forward floor equals tower.cu's formulas, and the inverse returns every sample exactly (\texttt{test/\allowbreak{}engine/\allowbreak{}compiler/\allowbreak{}cycle/\allowbreak{}record\_\allowbreak{}bitwise\_\allowbreak{}test}).
\item Derived in \hyperref[chap:vertical-time-compression]{vertical\_\allowbreak{}time\_\allowbreak{}compression.md}:

\begin{itemize}
\item The file holds live registers only. A field counts against it only while its register is live; the atom stays in the record.
\item keymath folds no constant, drops no dead step and merges no duplicate step.
\item A record is addressable to 2\textasciicircum{}32 bits for each member; past that, device memory bounds it.
\item A lifting cone spans (2\textasciicircum{}\{L+2\} − 3)\textasciicircum{}D samples, yet one level over a 4-D cone of 625 samples runs in a 15-limb file.
\item keymath's widths grow every register 1 bit a level, the first level 2, where the low-pass coefficients' values grow 0.585 bits and the high-pass coefficients' 1 (the linear forms of A16; E4). The constant-divisor narrowing is \textbf{proved} (\texttt{test/\allowbreak{}engine/\allowbreak{}compiler/\allowbreak{}cycle/\allowbreak{}record\_\allowbreak{}divide\_\allowbreak{}test}, 19 checks, 0 failed).
\item T is a bijection; an operation G between the towers runs as T⁻¹GT, and the levels above G cancel past it.
\item The heap (the bits a floor's values occupy) mirrors exactly across the crystal. The ring (the imprint's widths) does not, until a wrap at the mirror makes it symmetric to one bit a register. The crystal with the ring's framing is a code length bounding K(x) from above, and by counting no lens shortens most lanes.
\end{itemize}
\item Open: the neighbor gather in D axes (the halo in the record, or one axis per sweep), F ∘ T⁻¹'s error in D axes, which operations commute with T, the heap and the wrap at the mirror as committed tests (both measured in scratch runs; \textbf{proved} by \texttt{test/\allowbreak{}engine/\allowbreak{}compiler/\allowbreak{}cycle/\allowbreak{}record\_\allowbreak{}boundary\_\allowbreak{}test}), the oval (Doug's call), and the whole crystal as one stack (\hyperref[chap:vertical-time-compression]{vertical\_\allowbreak{}time\_\allowbreak{}compression.md}).
\item \textbf{Where the two towers stand in the engine.}

\begin{itemize}
\item In the engine, T and T⁻¹ are \texttt{tower.\allowbreak{}cu}'s own kernels (\texttt{tower\_\allowbreak{}forward\_\allowbreak{}kernel} and \texttt{tower\_\allowbreak{}inverse\_\allowbreak{}kernel}, behind \texttt{tower\_\allowbreak{}lift} and \texttt{tower\_\allowbreak{}lower}). The crystal path (\texttt{engine.\allowbreak{}cu}) and \texttt{root\_\allowbreak{}universal} call them.
\item The tower module also emits them as record floors. \texttt{tower\_\allowbreak{}record\_\allowbreak{}lift} and \texttt{tower\_\allowbreak{}record\_\allowbreak{}lower} pass a ruleset (E4) through a block's extent into a caller's program, one block of the lattice to a lane. Each position's register goes in and comes out in t, z, y, x order, and T leaves the crystal where \texttt{tower\_\allowbreak{}lift} leaves its coefficients. With no ruleset named, the ruleset is the kernels' 5/3. The steps go in at the caller's count; with no room given, the call only counts them, which sizes the program. A register that is not earlier, a room too small, an empty extent, and a step with a zero weight, no taps or a sign other than ±1 are each refused, and the program is left as it was. A program composes F ∘ T⁻¹ by laying its own steps after T⁻¹'s registers.
\item \textbf{Proved} (\texttt{test/\allowbreak{}engine/\allowbreak{}compiler/\allowbreak{}cycle/\allowbreak{}record\_\allowbreak{}tower\_\allowbreak{}test}, 209 checks, 0 failed, a job on tessera). There are seven blocks: 1 × 1 × 1 × 64, 1 × 4 × 4 × 4, 2 × 3 × 4 × 4, 1 × 3 × 5 × 7, 3 × 1 × 2 × 9, 2 × 2 × 2 × 2 and 1 × 1 × 1 × 1. Each carries 256 volumes: all 0, all 65,535, the two alternating, a ramp, and drawn. On every block:

\begin{itemize}
\item with no ruleset named, and with the 5/3 named as its two steps, T equals \texttt{tower\_\allowbreak{}lift}'s crystal at every coefficient on 256 of 256 volumes, and T⁻¹ run on it returns the source on 256 of 256;
\item a host oracle's 5/3 equals \texttt{tower\_\allowbreak{}lift} and returns the source on 256 of 256; against that oracle, the S transform's T and the four-tap predict's T each equal it on 256 of 256, and their T⁻¹ returns the source on 256 of 256;
\item for all four rulesets, T ∘ T⁻¹ as one program returns 256 of 256 crystals drawn anywhere in the fields' widths;
\item F ∘ T⁻¹ as one program, F the block's sum and its energy along x (the sum of the squared differences of x neighbors), equals F on \texttt{tower\_\allowbreak{}lower}'s samples on 256 of 256;
\item every program runs on the device and the host, and the records agree word for word.
\end{itemize}
\item \textbf{Measured}, the steps and the file (limbs of live registers, with reuse), the 5/3's first:

{\scriptsize
\setlength{\tabcolsep}{4pt}
\begin{longtable}{>{\RaggedRight\arraybackslash}p{\dimexpr0.1360\linewidth-2\tabcolsep\relax}>{\RaggedRight\arraybackslash}p{\dimexpr0.1352\linewidth-2\tabcolsep\relax}>{\RaggedRight\arraybackslash}p{\dimexpr0.1789\linewidth-2\tabcolsep\relax}>{\RaggedRight\arraybackslash}p{\dimexpr0.1026\linewidth-2\tabcolsep\relax}>{\RaggedRight\arraybackslash}p{\dimexpr0.1026\linewidth-2\tabcolsep\relax}>{\RaggedRight\arraybackslash}p{\dimexpr0.1675\linewidth-2\tabcolsep\relax}>{\RaggedRight\arraybackslash}p{\dimexpr0.1773\linewidth-2\tabcolsep\relax}}
\toprule
block & T steps, file & widest coefficient & T⁻¹ file & T ∘ T⁻¹ file & S transform T steps, file & four-tap T steps, file, T⁻¹ file \\
\midrule
\endhead
1 × 1 × 1 × 64 & 697, 70 & 23 bits & 71 & 71 & 317, 67 & 1,015, 73, 80 \\
1 × 4 × 4 × 4 & 1,192, 71 & 23 bits & 75 & 139 & 542, 68 & 1,735, 74, 80 \\
2 × 3 × 4 × 4 & 2,264, 103 & 24 bits & 119 & 215 & 1,070, 100 & 3,207, 106, 142 \\
1 × 3 × 5 × 7 & 2,216, 112 & 25 bits & 136 & 234 & 1,108, 109 & 3,009, 115, 193 \\
3 × 1 × 2 × 9 & 1,046, 61 & 24 bits & 61 & 111 & 510, 58 & 1,449, 64, 74 \\
2 × 2 × 2 × 2 & 356, 23 & 21 bits & 23 & 43 & 162, 20 & 519, 26, 26 \\
\bottomrule
\end{longtable}
}

T ∘ T⁻¹ on 1 × 3 × 5 × 7, 105 samples, holds 234 of the file's 256 limbs: a block is one lane with its samples as fields, and the file bounds the block.
\item The record programs the engine ran before (the cell program's, Ω's reduct in \texttt{chaitin\_\allowbreak{}omega}, π's BBP terms in \texttt{pi\_\allowbreak{}tower}) read their own fields and use neither tower.
\item \textbf{Built:} the towers as kernels, and as record floors from any ruleset. \textbf{Not built:} F ∘ T⁻¹ in the crystal path (\texttt{engine.\allowbreak{}cu} calls the kernels); the gather across blocks, since a lane holds one block; edges as record floors; and the kernels taking a ruleset.
\item Posits (theory, neither built nor tested): the scheduler is bypassed and the whole tower beamformed using its own inverse energy potential; and a two wave-vector magnitude interference pattern, with reinforced alignment in the ruleset and the exact no rounding as the destructive interference.
\end{itemize}
\end{itemize}
\item \textbf{The two crystals} (the finite towers and their limit, with an infinite delta between them).

\begin{itemize}
\item A register is a 2-adic integer, and a wrap to w is the projection ℤ₂ → ℤ/2\textasciicircum{}w. Sum, difference, product, xor and and commute with every projection. An exact quotient by an odd c is the product by c⁻¹ in ℤ₂.
\item \textbf{Proved} (\texttt{test/\allowbreak{}engine/\allowbreak{}compiler/\allowbreak{}cycle/\allowbreak{}record\_\allowbreak{}coherence\_\allowbreak{}test}, 14 checks, 0 failed): 393,216 of 393,216 lane-widths agree three ways over 16 random programs, 73,728 of 73,728 for the odd divisors, and a quotient and a comparison break the agreement. Mod 3\textasciicircum{}5, 3\textasciicircum{}10, 5\textasciicircum{}4 and 7\textasciicircum{}3, 262,144 of 262,144 lane-moduli agree three ways over 16 ring-only programs, the CRT join with the 8-bit window equals the exact run on all of them, and xor breaks mod 3.
\item \textbf{Proved} (\texttt{test/\allowbreak{}engine/\allowbreak{}compiler/\allowbreak{}cycle/\allowbreak{}record\_\allowbreak{}boundary\_\allowbreak{}test}, 48 checks, 0 failed): the crystal of 64 samples at 4 levels as a boundary. T's reach is 3ℓ at the level-ℓ low-pass coefficients and 3ℓ − 2 at its high-pass coefficients and T⁻¹'s is L + 2, each met on the device; arbitrary crystals return through T⁻¹ then T; constants and 2\textasciicircum{}\{3L\}ℤ\textasciicircum{}n moves pass through T, negation and doubling do not; the heap mirrors, the ring is ring\_\allowbreak{}0 + 6(n − n/2\textasciicircum{}ℓ); the quotient toward zero keeps order and a sum within one; T and T⁻¹ keep Haar measure, counted; det M = ±1 exactly, its sign +1 in the run. T against 8 null shuffles a lane identifies each structured class past the null's 1/(d + 1) bound and noise within it, and the crystal is a one-to-one ID: a draw changes it exactly when it moves a value, and every single flipped bit changes the image.
\item Derived in \hyperref[chap:two-crystals]{two\_\allowbreak{}crystals.md}: what does not factor through the projection, the lifting as a homeomorphism of ℤ₂\textasciicircum{}n whose level-L low-pass coefficients read the samples to w + 3L bits, the windows' union ℤ and the projections' limit ℤ₂ with the solenoid between them, what passes to a limit, the top projection's limit ℝ, the odd crystals and the places of ℚ, the count as coset counting, and the posits bounded (the anchors' complexity open; the fingerprint of T against null permutations, proved; the knf's identity by spatial null permutation, built; the pairwise agreement between sections, open).
\item \textbf{Proved} (\texttt{knf\_\allowbreak{}identity}, 23 checks, 0 failed): the knf of the nbody lattice (64³, 177 frames) against spatial null permutations. The host's walk equals the engine's knf at every voxel, and a single flipped bit leaves it unchanged exactly where the window rule says (1,500 of 8,192): the knf is many-to-one. The 48 exact cube motions carry the knf and keep its entangled entropy E and its cloud. The whole is its tiles plus its seams at every tile size. Alike-voxel controls are identified within 5σ of 1/20. \textbf{Measured:} E is 44.6 times the free null's mean, the inside and between departure curves cross between tiles of 2 and 4 and sum to 0.590 at 4, and each body's box has its own curve. Written in \hyperref[chap:two-crystals]{two\_\allowbreak{}crystals.md}.
\end{itemize}
\item \textbf{The full Kolmogorov–Arnold form.} With linear steps only, the machine held the theorem's skeleton with linear edges; the table gives it general one-variable edges, which is item 1 of \hyperref[chap:kolmogorov-arnold]{kolmogorov\_\allowbreak{}arnold.md} and the LUT of \hyperref[chap:noise-sieve-tower]{noise\_\allowbreak{}sieve\_\allowbreak{}tower.md} §2. Which functions to tabulate, without fitting, is still open (item 3 there).
\end{itemize}

\subsection{\texorpdfstring{A14. The root universal (M20): \texttt{tower}}{A14. The root universal (M20): tower}}

The tower and the table are separate objects; laid together they are one. Let W\_\allowbreak{}k be the approximation on floor k, L\_\allowbreak{}k level k's lifting, and E\_\allowbreak{}k the edges laid on floor k.

\begin{itemize}
\item \textbf{The edge is a bijection.} E(w) = (w with its low b bits x replaced by π(x)), π a permutation of [0, 2\textasciicircum{}b), the high bits passed through. So E is a bijection on the coefficient word and E⁻¹ uses π⁻¹ (proved by construction; the round trips below are the check).
\item \textbf{The order.} Lift is EN ∘ L(N−1) ∘ … ∘ E1 ∘ L0 ∘ E0, and lower is its inverse, E0⁻¹ ∘ L0⁻¹ ∘ … ∘ EN⁻¹. The edges on one floor are a stack, undone last laid first. \textbf{Proved:} every edge program in \texttt{tower\_\allowbreak{}edge\_\allowbreak{}test} and \texttt{root\_\allowbreak{}universal} rebuilds the exact lanes, 171 random edges on 48 volumes through the whole crystal path.
\item \textbf{The fold.} Two edges on one floor are one edge: π\_\allowbreak{}b ∘ π\_\allowbreak{}a, laid as one table, lays the same crystal. It keeps the order, since π\_\allowbreak{}a ∘ π\_\allowbreak{}b lays another, and a lifting level between two edges blocks the fold. \textbf{Proved} on 8 volumes each.
\item \textbf{What it is} (theory). The tower is the Kolmogorov–Arnold form's additive skeleton, the table its one-variable edge, and with edges between the floors the collapse is an outer sum with nonlinear edges.
\item \textbf{Its two bounds} (theory; kept). It is universal over the finite coefficient alphabet: every function of the indexed bits is one table, exactly. It is not the continuous-function theorem of Kolmogorov and Arnold. And inside the reversible stream every edge is a permutation, and the second bound is one of width, not of kind (below).
\item \textbf{The dimensional expansion} (reaching a boundary means a dimensional expansion realizes the transform). A function f that is not a bijection (|x|, a comparison) rides the stream through the Bennett embedding, (x, y) → (x, y ⊕ f(x)): a permutation of the pair for any f, since y comes back as the carrier ⊕ f(x) (proved by construction). The price is width: a permutation takes b\_\allowbreak{}x bits, any f takes b\_\allowbreak{}x + b\_\allowbreak{}y, within the 20-bit limit. No engine change; the tables are built on the host. \textbf{Proved} by \texttt{tower\_\allowbreak{}edge\_\allowbreak{}test}, 25 checks, 0 failed. |x| of an 8-bit two's complement field and x ≥ 100, each a 16-bit Bennett edge (8 bits of x, 8 of carrier) on the collapsed floor of 1 × 8 × 8 × 8, alternate over 64 lattices: they round-trip on 64 of 64; f(x) reads out of the carrier (x kept, the carrier moved by exactly f(x), the high bits passed) on 64 of 64; the edge touches only its one coefficient on 64 of 64. The two edges on floors 1 and 2, where lifting mixes their outputs, round-trip. |x| laid bare on 16 bits, not embedded, is refused, since x and −x collide. The quantum forms of the same expansion are standard theorems (theory here): Naimark (an ask is a projective measurement in a larger space), Stinespring (a transform is a unitary in a larger space, then a discard) and purification (a mixed state is part of a pure state in dimension d²); the ask is A15.
\item \textbf{The price} (measured, M20). An unfitted edge is exact but not free. On floor 0, where every coefficient is, a random 12-bit permutation raises the coded bits from 6.802 to 12.608 a voxel; on the collapsed floor, one coefficient, the cost does not show at three decimals at any width. The reading, which is not measured: a random permutation of the low bits sends small magnitudes to large ones, and the magnitude coder pays for each: the cost scales with the coefficients the edge touches. Whether a fitted or data-named edge lowers the bits, not only raises them, is not measured.
\end{itemize}

\subsection{\texorpdfstring{A15. The ask and the state (M21): \texttt{ask\_\allowbreak{}state}}{A15. The ask and the state (M21): ask\_state}}

The terms are M21's. The idea is standard: QBism (Fuchs and Schack) and the SIC-POVMs.

\begin{itemize}
\item \textbf{The ask returns the state.} For a qubit and E\_\allowbreak{}i = (I + v\_\allowbreak{}i·σ)/4, the answers are p\_\allowbreak{}i = (1 + r·v\_\allowbreak{}i)/4. With Σ v\_\allowbreak{}i = 0 and Σ v\_\allowbreak{}i v\_\allowbreak{}iᵀ = c·I, Σ p\_\allowbreak{}i v\_\allowbreak{}i = (c/4) r, which gives r = (4/c) Σ p\_\allowbreak{}i v\_\allowbreak{}i and the state is read back from the answers alone (theory; \textbf{proved} on 25 states for each ask, c = 1 and c = 4/3). A complete ask needs at least d² outcomes, here 4 (theory, standard).
\item \textbf{The phase falls out of the ask.} The equator states differ only in phase: each answers ½, ½ to the 0/1 ask, which sees magnitudes only. The complete ask tells all 91 pairs of the 14 apart (\textbf{proved}). So the phase is not lost; one basis's ask cannot see it, and a complete ask holds it.
\item \textbf{The crossing rule.} The chance of + on axis a is (1 + r\_\allowbreak{}a)/2 = Σ p\_\allowbreak{}i (½ + (2/c) v\_\allowbreak{}\{i,a\}), a fixed linear map from the complete ask's answers (theory; \textbf{proved}, 75 of 75 for each ask). Its weights onto +x are 3/2, 3/2, −1/2, −1/2 for the rational ask and ½ ± ½√3 for the SIC. Classical total probability never carries a negative weight, and the difference is interference.
\item \textbf{The valid set.} A distribution is a state exactly when |r|² ≤ 1. (1, 0, 0, 0) gives |r|² = 12 for the rational ask and 9 for the SIC, and it is refused, and it crosses onto +x at 3/2 and ½ + ½√3, no probabilities (\textbf{proved}). On the 455 distributions with denominator 12 the valid set is 55 (rational) and 87 (SIC), and 104 and 72 more pass the three axis crossings without being states (\textbf{measured}): the valid set is stricter than the axis crossings.
\item \textbf{Exact numbers.} Every value is exact: rationals on the exact integer, and the SIC's √3 carried by its relation, (a + b√3)(c + d√3) = (ac + 3bd) + (ad + bc)√3, with the sign of a + b√3 read from a² against 3b², never equal unless both are 0, since √3 is irrational (proved by construction). The SIC's answers cross into the rational ask's exactly, 25 of 25 (\textbf{proved}).
\item \textbf{The dimensional expansion} (theory, standard theorems). Naimark: an ask is a projective measurement in a larger space. Stinespring: a transform is a unitary in a larger space, then a discard. Purification: a mixed state is part of a pure state in dimension d². The engine's form is the Bennett embedding of A14, proved.
\item \textbf{The statements} (the first built in \texttt{ka\_\allowbreak{}psi}, the rest theory). The exp is symbolic after the second nth exp: ψ's terms γ\textasciicircum{}−β(r), with β(r) = (nʳ − 1)/(n − 1), are carried as a digit and an exact integer exponent, a sparse form whose cost is polynomial in the scale instead of the expanded digits. \texttt{ka\_\allowbreak{}psi} carries ψ this way to β(60) = 1,152,921,504,606,846,975 (built; 152 checks, 0 failed). ψ is the inner function of Sprecher's version of the superposition theorem, as J. Braun and M. Griebel state it (“On a constructive proof of Kolmogorov's superposition theorem”, Constructive Approximation 30(3) (2009) 653–675, doi:10.1007/s00365-009-9054-2, the record checked at Crossref; the preprint read, pages 1 to 5). Their Theorem 2.1 takes integers n ≥ 2, m ≥ 2n, γ ≥ m + 2 and a = [γ(γ − 1)]⁻¹, and writes f(x) = Σ\_\allowbreak{}\{q=0\}\textasciicircum{}\{m\} Φ\_\allowbreak{}q ∘ ξ(x\_\allowbreak{}q), with ξ(x\_\allowbreak{}q) = Σ\_\allowbreak{}\{p=1\}\textasciicircum{}\{n\} α\_\allowbreak{}p ψ(x\_\allowbreak{}p + qa), α\_\allowbreak{}1 = 1 and α\_\allowbreak{}p = Σ\_\allowbreak{}\{r≥1\} γ\textasciicircum{}−(p−1)β(r). ψ is defined on the terminating base-γ rationals D\_\allowbreak{}k = \{ Σ\_\allowbreak{}\{r≤k\} i\_\allowbreak{}r γ\textasciicircum{}−r \} (their 2.2), and in Köppen's form (their 2.7) level k adds its digit at the weight γ\textasciicircum{}−β(k). So γ\textasciicircum{}−β(r) is the γ-adic weight of level r. What ψ is, strictly increasing and continuous in Köppen's form and not so in Sprecher's, is proved exactly by the sim \texttt{ka\_\allowbreak{}psi} (M19), to depth 60 with the exponent never expanded, in “Kolmogorov's inner function, exactly” of \hyperref[chap:kolmogorov-arnold]{kolmogorov\_\allowbreak{}arnold.md}. Its “Separation” part finds the paper's Lemma 3.4 false at k = 1 and restores Lemmas 3.7 and 3.8 at every k with a corrected bound and a ramp narrowed by γ². The outer functions' contraction is not built. A two-sqrt product represented as two ints and their words of math operations is countable: the numbers named by finitely many integers and a finite word of operations (the algebraic, the computable, π as its program) are countable, and the exact storage holds each of them. And π compresses to 1 kb: π as a spigot program.
\end{itemize}

\subsection{\texorpdfstring{A16. The Gaussian step (M10, M19): \texttt{test/\allowbreak{}engine/\allowbreak{}compiler/\allowbreak{}cycle/\allowbreak{}record\_\allowbreak{}gaussian\_\allowbreak{}test}}{A16. The Gaussian step (M10, M19): test/engine/compiler/cycle/record\_gaussian\_test}}

BBP's base is a power of a Gaussian prime (the shape of the algorithm dictates the geometric ruleset to enforce, and where the geometry looks absurd the algorithm is adjusted until it is fluid). Each rule carries one label: \textbf{cited} (a published or standard theorem), \textbf{derived} (exact algebra worked here and checked in floating point by a scratch script; theory, not built), \textbf{proved} (the run named) or \textbf{open} (not known).

\begin{itemize}
\item \textbf{The ring and its prime} (cited, standard). ℤ[i] = \{a + bi : a, b ∈ ℤ\} is Euclidean under the norm N(a + bi) = a² + b². 2 = −i(1 + i)², and so 1 + i is the only prime over 2, ramified, with N(1 + i) = 2. The completion there is ℤ₂[i], with 1 + i its uniformizer: v(1 + i) = 1 and v(2) = 2.
\item \textbf{Eight steps are sixteen} (derived; proved below). (1 + i)² = 2i, which gives (1 + i)\textasciicircum{}8 = (2i)\textasciicircum{}4 = 16. The powers for k = 1 to 8 are 1 + i, 2i, −2 + 2i, −4, −4 − 4i, −8i, 8 − 8i and 16.
\item \textbf{The step turns and scales} (derived). 1 + i = √2·e\textasciicircum{}\{iπ/4\}: z ↦ z(1 + i) turns z by π/4 and scales it by √2, and N doubles exactly. Eight steps are a whole turn and a scale of 16 = 2\textasciicircum{}4. So BBP's digit shift, ×16, one hex digit, is eight Gaussian steps.
\item \textbf{The record floor} (derived; proved below). On z = a + bi the step is (a, b) ↦ (a − b, a + b): one DIFFERENCE and one SUM, the matrix [[1, −1], [1, 1]], determinant 2.

\begin{itemize}
\item Its image is the pairs of equal parity: the ideal (1 + i), of index 2. Each floor therefore drops one bit, and after k floors the image is (1 + i)\textasciicircum{}k ℤ[i], of index 2\textasciicircum{}k. v rises by 1 a floor and N doubles: the bit leaves the archimedean side for the (1 + i)-adic one.
\item The inverse floor is ((c + d)/2, (d − c)/2): two EXACT\_\allowbreak{}QUOTIENTs by 2, exact on the image. A pair of mixed parity is exactly a lane the machine refuses (derived; proved below).
\item SUM and DIFFERENCE commute with every wrap (E1, proved by \texttt{test/\allowbreak{}engine/\allowbreak{}compiler/\allowbreak{}cycle/\allowbreak{}record\_\allowbreak{}coherence\_\allowbreak{}test}), and the floor runs exactly in any window. In a window of w bits each floor keeps an index-2 part of the pairs, and 2w floors send every pair to 0, since (1 + i)\textasciicircum{}\{2w\} = (2i)\textasciicircum{}w ≡ 0 mod 2\textasciicircum{}w (derived).
\end{itemize}
\item \textbf{The product formula over ℚ(i)} (the formula cited, standard; the values derived). Take |x|\_\allowbreak{}ℂ = x·x̄ and |x|\_\allowbreak{}\{1+i\} = 2\textasciicircum{}\{−v(x)\}; every other place gives 1 on these values. Then |16|\_\allowbreak{}ℂ·|16|\_\allowbreak{}\{1+i\} = 256·2\textasciicircum{}\{−8\} = 1. A floor multiplies |·|\_\allowbreak{}ℂ by 2 and |·|\_\allowbreak{}\{1+i\} by ½, and the product holds at 1 on every floor: the floor moves a bit from one place to the other and loses none.
\item \textbf{π from the step} (the series cited; the rest derived).

\begin{itemize}
\item BBP (Bailey, Borwein and Plouffe~\cite{src:Bailey-Borwein-and-Plouffe-1997}, Math. Comp. 66, 1997): π = Σ\_\allowbreak{}k 16\textasciicircum{}\{−k\}(4/(8k + 1) − 2/(8k + 4) − 1/(8k + 5) − 1/(8k + 6)).
\item Expanding 1/(1 − x⁸) term by term gives Σ\_\allowbreak{}k 16\textasciicircum{}\{−k\}/(8k + j) = 2\textasciicircum{}\{j/2\} ∫\_\allowbreak{}0\textasciicircum{}\{1/√2\} x\textasciicircum{}\{j−1\}/(1 − x⁸) dx. So π = ∫\_\allowbreak{}0\textasciicircum{}\{1/√2\} (4√2 − 8x³ − 4√2x⁴ − 8x⁵)/(1 − x⁸) dx, the integral their proof runs through. The upper limit is 1/√2 = |1/(1 + i)|, and (1/√2)\textasciicircum{}8 = 1/16.
\item \textbf{Adjusted: four of the eight roots of unity are poles instead of all eight.} 1 − x⁸ vanishes at the eighth roots of unity, e\textasciicircum{}\{ikπ/4\} = ((1 + i)/√2)\textasciicircum{}k. The ruleset as first stated put a pole at each. The numerator is −8(1 + x²)(x − 1/√2)(x² + √2x + 1), which cancels ±i and e\textasciicircum{}\{±3iπ/4\}. What remains is 8(1/√2 − x)/((1 − x²)(x² − √2x + 1)), with poles at ±1 and e\textasciicircum{}\{±iπ/4\} = (1 ± i)/√2 only, and it vanishes at the upper limit.
\item With y = √2x, the form BBP's proof reaches: π = ∫\_\allowbreak{}0\textasciicircum{}1 16(y − 1)/((y² − 2)(y² − 2y + 2)) dy. y² − 2y + 2 = (y − ρ)(y − ρ̄) with ρ = 1 + i. The poles are the Gaussian prime, its conjugate, and ±√2 = ±|ρ|. By partial fractions the integrand is 2/(y − √2) + 2/(y + √2) − 2ρ/(y − ρ) − 2ρ̄/(y − ρ̄). The residue at the prime is −2ρ, and the four residues sum to 0.
\item Over [0, 1] the poles at ±√2 give −2 ln 2 and the real part at ρ gives +2 ln 2, and they cancel. What remains is 4·arctan 1 = 4·(π/4), and π/4 is the angle [0, 1] subtends at ρ, arg(−i) − arg(−1 − i). So π is four times one Gaussian step's turn, and BBP's logarithms, the step's scale (ln N(ρ) = ln 2), cancel exactly: the series keeps the turn and drops the scale.
\end{itemize}
\item \textbf{The width, made fluid} (the keymath rule; built, and proved at every floor).

\begin{itemize}
\item The superseded widening takes a SUM or DIFFERENCE to max + 1, with floor k at w + k bits, 25 to 32 at w = 24, and is not used.
\item Floor k's values are (a·x\_\allowbreak{}k − b·y\_\allowbreak{}k, a·y\_\allowbreak{}k + b·x\_\allowbreak{}k) with (1 + i)\textasciicircum{}k = x\_\allowbreak{}k + i·y\_\allowbreak{}k, and |x\_\allowbreak{}k| + |y\_\allowbreak{}k| = 2\textasciicircum{}\{⌈k/2⌉\}. So they grow half a bit a floor. The script's corner figures (25, 26, 27, 27, 28, 28, 28 and 28) are two's complement widths, sign included (checked by hand at floors 2, 3, 7 and 8). The register holds the magnitude with its sign beside it, and there floor k needs 24 + ⌈k/2⌉.
\item The cause: the old rule bounded each register alone. An odd floor's values fill a turned square, |c| + |d| ≤ 2\textasciicircum{}\{⌈k/2⌉\}·2\textasciicircum{}\{w−1\}, and a per-register box keeps its bounding square, which the next floor turns into twice the reach. Over 8n floors the imprint grew 8n bits and the values 4n.
\item The rule built (\texttt{keymath.\allowbreak{}cu}): every register is a linear form over atoms, integer coefficients plus a constant. A sum or difference adds its operands' forms, a product by a constant scales the other's, a constant is its value, and a wrap that passes its register through keeps that register's form. Every other register is an atom. The width is the fewer of the operation's own rule and the form's bound, |x| ≤ |c| + Σ |c\_\allowbreak{}i|(2\textasciicircum{}\{b\_\allowbreak{}i\} − 1). A coefficient past 2\textasciicircum{}62 makes the register an atom. On the Gaussian floor the coefficients are the parts of (1 + i)\textasciicircum{}k: the bound is 2\textasciicircum{}\{⌈k/2⌉\}·2\textasciicircum{}w and the width is w + ⌈k/2⌉.
\item Read through Mathai and Thiang (arXiv:1505.05250): T-duality turns the bulk-boundary map ∂ into a restriction ι*. Here the width, a boundary quantity, is read by restricting the bulk form to its atoms' widths, not carried register by register (an analogy, not a theorem about keymath).
\item The wrap at w + ⌈k/2⌉ + 1 bits a floor is not needed.
\end{itemize}
\item \textbf{Proved} (\texttt{test/\allowbreak{}engine/\allowbreak{}compiler/\allowbreak{}cycle/\allowbreak{}record\_\allowbreak{}gaussian\_\allowbreak{}test}, 16 checks, 0 failed, a job on tessera). The program is two signed 24-bit fields and eight floors of one DIFFERENCE and one SUM, 18 steps with every floor an output. It ran on 4,096 lanes: every pair of the five edge values (0, −1, 1, −2\textasciicircum{}23 and 2\textasciicircum{}23 − 1), and keyed draws for the rest. Its checks:

\begin{itemize}
\item the program imprints, lays and loads, and floor k is 24 + ⌈k/2⌉ bits wide: 25, 25, 26, 26, 27, 27, 28 and 28;
\item some lane fills every floor's width, |value| ≥ 2\textasciicircum{}\{bits − 1\}, and no width is loose;
\item the device equals the host word for word;
\item every floor equals z(1 + i)\textasciicircum{}k, taken on the CPU from the powers above;
\item floor 4 is −4z and floor 8 is 16z on every lane;
\item each floor's two registers share their parity;
\item floor 8's values lie in [−2\textasciicircum{}27, 2\textasciicircum{}27), 28 bits, and reach −2\textasciicircum{}27;
\item eight inverse floors take 16z back through z(1 + i)\textasciicircum{}\{8−k\} to z on every lane, in the host's records and the device's alike;
\item the chosen pairs w(1 + i)\textasciicircum{}k, for k from 0 to 9 with six drawn w of valuation 0, and then 0, have valuation k on the CPU;
\item through 1 to 8 inverse floors the host and the device agree on every run, j inverse floors refuse a pair exactly where (1 + i)\textasciicircum{}j does not divide it, and where they run each level is z/(1 + i)\textasciicircum{}k, the Gaussian division on the CPU. The run printed 272 runs and 216 refused. Derived: 61 pairs through 8 floor counts are 488 runs; a pair of valuation k runs the j ≤ min(k, 8) of them, 6·(0 + 1 + … + 8 + 8) + 8 = 272, and 488 − 272 = 216;
\item tessera admits the job and it releases.
\end{itemize}
\item \textbf{The rule across the other record tests} (each a job on tessera, run one at a time): guide 12, table 16, divide 20, bitwise 25, coherence 15 and order 18 checks, 0 failed. The boundary ring law is ring\_\allowbreak{}ℓ = ring\_\allowbreak{}0 + n + 2n(1 − 2\textasciicircum{}\{−ℓ\}) in E4's widths, and boundary is 49 checks, 0 failed.
\item \textbf{Where it stands in the engine.} The width rule is in \texttt{keymath.\allowbreak{}cu}, and every record program takes it. The Gaussian step itself is built only in the test. \texttt{pi\_\allowbreak{}tower} runs 16\textasciicircum{}\{d−i\} mod m a nibble at a time (E13) and does not use the step.
\item \textbf{The bridge} (open).

\begin{itemize}
\item The digit at d is \{16\textasciicircum{}d π\} = \{4S₁ − 2S₄ − S₅ − S₆\}, where S\_\allowbreak{}j = Σ\_\allowbreak{}\{i≤d\} (16\textasciicircum{}\{d−i\} mod m\_\allowbreak{}i)/m\_\allowbreak{}i plus a tail, with m\_\allowbreak{}i = 8i + j.
\item Each term's power reduces by Carmichael's λ (cited, standard): on m\_\allowbreak{}i's odd part, the exponent is taken mod λ; on its power of 2, the term is 0 once 4(d − i) reaches that power (derived). Still, d + 1 terms remain.
\item The Gaussian form does not shorten them. 16 is a rational integer, and (1 + i)\textasciicircum{}\{8(d−i)\} mod m\_\allowbreak{}i in ℤ[i]/(m\_\allowbreak{}i) is 16\textasciicircum{}\{d−i\} mod m\_\allowbreak{}i (derived).
\item What would close it is a sum over i taken at once: a modular form or Hecke eigenvalue whose coefficient is Σ\_\allowbreak{}i r\_\allowbreak{}i/m\_\allowbreak{}i mod 1. The Hecke recursion, T\_\allowbreak{}\{p\textasciicircum{}\{n+1\}\} = T\_\allowbreak{}p·T\_\allowbreak{}\{p\textasciicircum{}n\} − p\textasciicircum{}\{k−1\}·T\_\allowbreak{}\{p\textasciicircum{}\{n−1\}\} on level-1 forms of weight k (cited, standard), makes a coefficient at 2\textasciicircum{}\{4d\} a power of a 2 × 2 matrix, about log d steps (derived).
\item No form whose coefficients give BBP's sum is known. No method is known to us that finds π's hex digit at d in fewer than about d steps (open).
\end{itemize}
\end{itemize}

\subsection{\texorpdfstring{A17. Exact qubit states (M22): \texttt{qasm}}{A17. Exact qubit states (M22): qasm}}

The module is the exact qubit states in C at \texttt{examples/\allowbreak{}qasm/\allowbreak{}}, and each reading below is checked by the qasm tests, \texttt{examples/\allowbreak{}qasm/\allowbreak{}test/\allowbreak{}run.\allowbreak{}sh} (39 checks, 0 failed, on the host). Two of the checks are regression checks: the bond at the widest site index reads 0, and w\textasciicircum{}(bits−1) at w = 2 evaluates to 2\textasciicircum{}(bits−1).

\begin{itemize}
\item \textbf{The field.} Q(√2)[i] is the numbers a + b√2 with a and b Gaussian rationals. H's 1/√2 is √2/2, and the T phase e\textasciicircum{}\{iπ/4\} is (1 + i)·√2/2. A circuit over X, Y, Z, S, H, CNOT, CZ and the controlled T phase keeps every amplitude in the field. Each rational sits on the exact integer (\texttt{arithmetic/\allowbreak{}no\_\allowbreak{}rounding}, 128 limbs by default), and a value past that width is refused, never wrapped (built).
\item \textbf{The chain.} Each qubit is a site tensor. After a two-site gate the pair's matrix M is factored as M = C·F by Gaussian elimination, and the bond between them is M's exact rank; nothing is truncated. The same code runs over Q(√2)[i] and over \texttt{Q(√2)[i](ω)}. The readings are in M22 (\textbf{proved} for those circuits).
\item \textbf{The symbolic phase.} \texttt{Q(√2)[i](ω)} holds rational functions in ω, each in rat\_\allowbreak{}make's canonical form, and evaluates them at a number. Readings (\textbf{proved}):

\begin{itemize}
\item Delta-Bell: bond [2]; ⟨00⟩ = “1/2sqrt2”, ⟨11⟩ = “(1/2sqrt2)w”, ⟨ψ|ψ⟩ = 1.
\item ⟨X0X1⟩ = “(1/2 + (1/2)w\textasciicircum{}2) / (w)”, ⟨Z0Z1⟩ = 1, ⟨X0⟩ = 0.
\item Delta-GHZ3: bonds [2,2], and ⟨X0X1X2⟩ is the same function.
\item |+++⟩ with a phase on wire 0 keeps bonds [1,1].
\item At ω = e\textasciicircum{}\{iπ/4\}, ⟨X0X1⟩ equals the numeric controlled-T reading, “1/2sqrt2”, and the norm is exactly 1.
\end{itemize}
\item \textbf{The boundary lens.} ⟨ψ|O|ψ⟩ is contracted from both edges of the chain inward. \texttt{qasm\_\allowbreak{}lens\_\allowbreak{}hash.\allowbreak{}c} and \texttt{qasm\_\allowbreak{}lens\_\allowbreak{}rank.\allowbreak{}c} compute the reduced-round SHA-256 wedge-field invariants: rank, complement, degree-2 lift rank, curvature vanish counts, fibers and clock closure. At mid 13, p 13, q 8, 6 bits and seed 2024, rounds 16 to 30 each read raw rank 126, complement 130, lift extra 0, vanish 30 and 14 of 400, and fibers 64 and 64; the clock closed and the root was stable (\textbf{proved} at those arguments).
\item \textbf{The witness.} The root is the engine's seal (\texttt{obsignatio\_\allowbreak{}seal}). The checks are that a seal repeats, that one rational moved by 1/10\textasciicircum{}9 changes the dense state's seal, and that one flipped generator bit changes the lens's (\textbf{proved}).
\item \textbf{Refusals.} Each is a request error from qasm: a gate past the state, a division by zero, a two-site gate on the last site, inverting the zero function, and an aperture past 20 bits (\textbf{proved}).
\item \textbf{The device code} (M22) is built, on the engine's fully vectorized self run: record floors of integer steps, four parts an amplitude, one \texttt{cycle\_\allowbreak{}record\_\allowbreak{}run} launch over every lane, host and device records equal word for word. A tessera job: qasm\_\allowbreak{}test 86 passed, 0 failed; qasm exact test 39 checks, 0 failed; qasm self test 44 checks, 0 failed; the CLI on bell.qasm exit 3, bitstring 00, and on bernstein\_\allowbreak{}vazirani.qasm exit 0, bitstring 101101001101.
\item \textbf{Not built into the engine library.} qasm in the engine library build: \texttt{src/\allowbreak{}build\_\allowbreak{}engine.\allowbreak{}sh} (MODULES and INGEST) and \texttt{src/\allowbreak{}cu/\allowbreak{}CMake\allowbreak{}Lists.\allowbreak{}txt} name no qasm source. The QASM ingester is \texttt{qasm\_\allowbreak{}read.\allowbreak{}c} and \texttt{qasm\_\allowbreak{}lexer.\allowbreak{}c}.
\end{itemize}

\subsection{\texorpdfstring{A18. The ask (M24): \texttt{query\_\allowbreak{}ask}, \texttt{gate\_\allowbreak{}descent}, \texttt{ask\_\allowbreak{}order}, \texttt{interface\_\allowbreak{}sass\_\allowbreak{}probe}}{A18. The ask (M24): query\_ask, gate\_descent, ask\_order, interface\_sass\_probe}}

The algebra of each step and the run behind it are P1 to P8 of \hyperref[chap:query-protocol-table]{query\_\allowbreak{}protocol\_\allowbreak{}table.md}. What is settled:

\begin{itemize}
\item \textbf{The ask holds on the host and in the cell} (\textbf{proved}, Q1). \texttt{query\_\allowbreak{}ask} finds memory the test owns holding and leaves it as it was, reads a written word as that word and no other, reads a word nothing writes as not advancing, and reads a bound with no clock as past the bound. \texttt{query\_\allowbreak{}interface\_\allowbreak{}walk} keeps an ending as the answer of the address that caused it.
\item \textbf{The gate is a conjunction} (\textbf{proved}, Q4). The descent's survivors equal every candidate asked every case, and both equal what \texttt{chain\_\allowbreak{}build} keeps with its sweep: add 1,202 of 1,227, take 1,047 of 1,083, up 411 of 698, down 408 of 559, same 0 of 8, 0 lost.
\item \textbf{The order and its solve} (\textbf{proved}, Q5). \texttt{4·Sᵀ − 2·J} times \texttt{S} is \texttt{(n + 1)} times the identity entry by entry, every ask covers \texttt{(n + 1)/\allowbreak{}2} links and any two share \texttt{(n + 1)/\allowbreak{}4}, and integer link costs solve back exactly at 1 to 8 passes. Against one ask a link the gain is √(n + 1)/2 (\textbf{measured} 1.09, 2.07 and 8.16 at 3, 15 and 255 links).
\item \textbf{Contention} (\textbf{measured}, Q7): found on 3.0\% with none, 83.6\% with up to 80 a pair and 100\% with up to 200, over 15 links and 45 sweep asks.
\item \textbf{One ask, many answers} (\textbf{measured}, Q15). One bit a check takes the fewest asks on every relation. Past a failing share of (3 − √5)/2 asking one check at a time costs no more than any scheme that asks a set as one bit (cited from knowledge: P. Ungar, “The cutoff point for group testing”, Comm. Pure Appl. Math. 13 (1960) 49–54).
\item \textbf{A form asked of the part} (\textbf{proved}, Q16). On sm\_\allowbreak{}86, 8 of the 2927 forms in \texttt{loop\_\allowbreak{}back\_\allowbreak{}if}'s place answer every count, every one a \texttt{BRA}, and \texttt{sass\_\allowbreak{}loop\_\allowbreak{}walk} agrees with all 8. The next above the least costs 0.0525 ns a turn more against a spread of 0.9260 (\textbf{measured}): no form is cheaper, and \texttt{sass.\allowbreak{}krs} keeps its \texttt{BRA}.
\end{itemize}

\section{\texorpdfstring{The scale audit: every number in the machine that fixes a scale}{The scale audit: every number in the machine that fixes a scale}}

A number in this table is a debt unless it is a word's width or a format's specification. Each is to come from the request or the data, or to be measured and reported, never tuned to a program.

{\footnotesize
\setlength{\tabcolsep}{6pt}
\begin{longtable}{>{\RaggedRight\arraybackslash}p{\dimexpr0.2589\linewidth-2\tabcolsep\relax}>{\RaggedRight\arraybackslash}p{\dimexpr0.2429\linewidth-2\tabcolsep\relax}>{\RaggedRight\arraybackslash}p{\dimexpr0.2141\linewidth-2\tabcolsep\relax}>{\RaggedRight\arraybackslash}p{\dimexpr0.2841\linewidth-2\tabcolsep\relax}}
\toprule
number & where & what it fixes & why it is a debt \\
\midrule
\endhead
the even-order guard & \texttt{residual.\allowbreak{}cu}, \texttt{keymath.\allowbreak{}cu}, \texttt{unit\_\allowbreak{}sweep.\allowbreak{}cu} & refuses every odd order & stricter than the algebra (A1); ruling (a) replaces it \\
\texttt{ENGINE\_\allowbreak{}RESIDUAL\_\allowbreak{}LIMBS} 9 & engine\_\allowbreak{}config.h & the residual's width: 16 input bits plus the orders' growth & sized for one program's orders; the width should follow the request's orders (each step adds a bit). The planes path derives it (input bits + Σ orders + 1) and returns it (M2); the 16-bit path and the key still use 9 \\
\texttt{UNIT\_\allowbreak{}SWEEP\_\allowbreak{}INPUT\_\allowbreak{}BITS} 16, unsigned 8 and 16 bit input only & unit\_\allowbreak{}sweep.cu, \texttt{engine\_\allowbreak{}source\_\allowbreak{}read} & the input's depth & a source deeper or signed needs a representation, not a refusal (the RSNA ruling: signed stored + 2\textasciicircum{}15 into the u16 lane, Doug to confirm). The unit sweep takes planes of any declared width (M2); its volume path and \texttt{engine\_\allowbreak{}source\_\allowbreak{}read} stay at 16 \\
\texttt{ENGINE\_\allowbreak{}AXES} 3 & engine\_\allowbreak{}config.h & three spatial axes & a word of the problem, but the residual, the tree and C are separable per axis \\
\texttt{ENGINE\_\allowbreak{}HISTORY\_\allowbreak{}WINDOW} 11, \texttt{ENGINE\_\allowbreak{}HISTORY\_\allowbreak{}WINDOWS\_\allowbreak{}MAX} 16 & engine\_\allowbreak{}config.h & the entropy window and how many fit & a window is a scale; it belongs in the request \\
\texttt{ENGINE\_\allowbreak{}COEFFICIENT\_\allowbreak{}LIMIT} 2\textasciicircum{}30 & engine\_\allowbreak{}config.h, read by tower & the largest lifted coefficient & a word's width; an overflow is a request error, not a silent clip \\
\texttt{COMPRESSION\_\allowbreak{}BLOCK} 64, 64 blocks a chunk, \texttt{COMPRESSION\_\allowbreak{}K\_\allowbreak{}BITS} 5, \texttt{COMPRESSION\_\allowbreak{}ESCAPE} 24 & compression.cu & the Rice coder's block and chunk & chosen; measure them against the data, or let the file carry them \\
\texttt{CRC\_\allowbreak{}SEGMENT\_\allowbreak{}PIXELS} 64 & crc.h & the CRC's parallel segment & no longer read by the tower (the seal replaced its CRC); crc.h goes with stage B \\
free / 3 · 2 & schedule.cu & the share of free device memory the schedule plans against & chosen (flagged in the RSNA work); ingest, prove and flatten consult no free memory at all; tessera arbitrates the device between processes (M14) \\
extents ≤ 2\textasciicircum{}20 and lanes ≤ 2\textasciicircum{}40 & krep history read, \texttt{krep\_\allowbreak{}lanes}, \texttt{entry\_\allowbreak{}lanes} & the largest sample & chosen ceilings; the words hold more \\
\texttt{MARGINAL\_\allowbreak{}ARRANGEMENTS\_\allowbreak{}MOST} 2\textasciicircum{}16, \texttt{MARGINAL\_\allowbreak{}TARGETS\_\allowbreak{}MOST} 64, \texttt{MARGINAL\_\allowbreak{}LIMBS} 16 & marginal & the largest local set & OrganoidTracker's bound and not measured here; A5 says report the component's size instead \\
\texttt{SHIFT\_\allowbreak{}AGREEMENT\_\allowbreak{}LONGEST\_\allowbreak{}AXIS} 2\textasciicircum{}23 & shift\_\allowbreak{}agreement & the NTT's longest padded axis & the prime's transform length; a word's width, but a larger view needs a second prime, not a refusal \\
\texttt{ENGINE\_\allowbreak{}RECORD\_\allowbreak{}LIMBS\_\allowbreak{}MOST} 256, \texttt{CYCLE\_\allowbreak{}RECORD\_\allowbreak{}BLOCKS\_\allowbreak{}MOST} 4096, \texttt{ENGINE\_\allowbreak{}RECORD\_\allowbreak{}MEMBERS\_\allowbreak{}MAX} 3 & engine\_\allowbreak{}config.h, cycle.h & the record machine's reach; 256 limbs of live registers is the binding resource (A13) & not measured \\
\texttt{ENGINE\_\allowbreak{}RECORD\_\allowbreak{}WRAP\_\allowbreak{}BITS\_\allowbreak{}LEAST} 4 & engine\_\allowbreak{}config.h & the narrowest two's complement wrap, a nibble (A13) & a floor the request sets instead of a resource; there is no step cap, since each sign sits beside its register \\
\texttt{ENGINE\_\allowbreak{}RECORD\_\allowbreak{}TABLE\_\allowbreak{}INDEX\_\allowbreak{}BITS\_\allowbreak{}MOST} 32 & engine\_\allowbreak{}config.h & a table's index, the low bits of one register (A13) & a word's width: the index fits one limb \\
\texttt{TOWER\_\allowbreak{}EDGE\_\allowbreak{}INDEX\_\allowbreak{}BITS\_\allowbreak{}MOST} 20 & tower.h & a tower edge's widest index, 2\textasciicircum{}20 entries of 4 bytes, 4 MiB a table (A14) & not measured \\
\texttt{BODY\_\allowbreak{}OVERLAP\_\allowbreak{}CHUNK} 256, \texttt{CLIMB\_\allowbreak{}MACHINE\_\allowbreak{}BUNDLE} 32, \texttt{MAX\_\allowbreak{}TREE\_\allowbreak{}BUCKETS} 256, \texttt{UNIT\_\allowbreak{}SWEEP\_\allowbreak{}SHARED\_\allowbreak{}BYTES\_\allowbreak{}MOST} 48 KiB & kernels & device grains & the device's, not the data's; measure per device \\
\texttt{DEVICE\_\allowbreak{}POOL\_\allowbreak{}SLICE\_\allowbreak{}ALIGN} 256 bytes, \texttt{DEVICE\_\allowbreak{}POOL\_\allowbreak{}PAGE\_\allowbreak{}BYTES} 2 MiB & device\_\allowbreak{}pool.h & where a pool's slices sit, and the grain a pool is rounded up to (M9) & the device's allocation grain, not the data's: on this machine four allocations cost 14,680,064 bytes where one pool of the same buffers cost 10,485,760; measure per device \\
\texttt{STACK\_\allowbreak{}UNCACHED\_\allowbreak{}BLOCK} 64 MiB, \texttt{NPY\_\allowbreak{}FORTRAN\_\allowbreak{}BLOCK} 16 MiB & stack, npy & IO grains & size by experiment to the disk's throughput (build plan 22); there is no header-room cap: the npy header is bounded by its own length field and the file, and the nrrd header is read in doubling pieces until its blank line \\
flatten's slab, half the free device memory & flatten & frames per residual pass & measured 107.0 s against 9.0 s one frame at a time; a fraction of memory is not a scale of the problem \\
\texttt{SCRIPTURA\_\allowbreak{}DIRECT\_\allowbreak{}BYTES} 32 & scriptura.c & below it a memory call skips the arm table for portable & the portable short path's reach (four words), not the data's; measured only on this part \\
\texttt{SCRIPTURA\_\allowbreak{}AVX2\_\allowbreak{}STRING\_\allowbreak{}BYTES} 16 KiB & scriptura\_\allowbreak{}arm\_\allowbreak{}avx2.c & where copy and fill switch to \texttt{rep movsb}/\texttt{stosb} & the part's, not the data's: the measured crossover on a Haswell-E with ERMS (4 KiB lost 1.28×, 16 KiB and up at 1.00× the runtime); other parts move it, and each part is measured on its own \\
\texttt{PERIOD\_\allowbreak{}THREADS} 256, \texttt{PERIOD\_\allowbreak{}GRID\_\allowbreak{}ROWS\_\allowbreak{}MOST} 65,535 & period.cu & a block's width and the grid's row limit & the device's, not the data's: the columns are read from the device (multiprocessors × resident blocks), and the lags beyond the row limit are walked by stride \\
\texttt{PERIOD\_\allowbreak{}ROUNDS} 4 & period.cu & the keys of the line shuffle's hash in each draw of the null (A12; the Feistel shuffle's rounds before it) & a mixing depth, not a scale of the data \\
fewer than 2\textasciicircum{}32 voxels a reading & period.cu & the largest lattice a period is read on & a word's width: N² and Σc² must fit 64 bits; a larger lattice needs wider counts, not a refusal \\
\texttt{ANCHOR\_\allowbreak{}EXACT\_\allowbreak{}LIMBS} 256 in the cell\_\allowbreak{}tracking build (the default is 128) & \texttt{cell\_\allowbreak{}tracking/\allowbreak{}build\_\allowbreak{}engine.\allowbreak{}sh} and \texttt{build\_\allowbreak{}driver.\allowbreak{}sh}, fitted from \texttt{ENGINE\_\allowbreak{}RECORD\_\allowbreak{}LIMBS\_\allowbreak{}MOST} by a define, with no change to the source & 8,192 bits, 1,024 digits & a word's width, derived and not chosen; \texttt{decimal\_\allowbreak{}double} reports \texttt{needed\_\allowbreak{}bits} past it \\
\bottomrule
\end{longtable}
}

There is no npy zip directory cap (64 MiB).

\section{\texorpdfstring{What the engine offers programs, and what it asks}{What the engine offers programs, and what it asks}}

{\small
\setlength{\tabcolsep}{6pt}
\begin{longtable}{>{\RaggedRight\arraybackslash}p{\dimexpr0.7484\linewidth-2\tabcolsep\relax}>{\RaggedRight\arraybackslash}p{\dimexpr0.2516\linewidth-2\tabcolsep\relax}}
\toprule
the engine offers & through \\
\midrule
\endhead
exact numbers from decimal strings, with \texttt{needed\_\allowbreak{}bits} & M1 \\
the residual at any orders and any input width the request names, its width returned & M2 \\
the component tree at every level: cuts, probe partitions, the bipartite census over probe links & M3 \\
moments per component & M4 \\
one correlation C and its reductions & M5 \\
the exact marginal and one-to-one matching on the caller's integer weights & M6, M7 \\
the entropy history and its cloud & M8 \\
proved files, and every source format read exactly & M9 \\
a seal on every crystal and set, and the place of any mismatch & M13 \\
one device shared between processes by memory, each job measured by pid & M14 \\
the sift: a pattern found by anchors whose count the corpus's coherence sets & M15 \\
a step's state drawn as a sheet or a volume, host and device alike & M16 \\
the record machine for the program's own verdicts & M10 \\
reversible lookup edges between the tower's floors, the root universal & M20 \\
the root noise of a box, the shared term whose yank saves the crystal the most bits, and the box back from the yank exactly & E16 \\
a stateful error on every failure & M12 \\
length, find, copy, move, compare and fill at each instruction set's speed & M17 \\
\bottomrule
\end{longtable}
}

What the engine asks of a program: every scale (orders, windows, widths, spacings, lags, probes) in the request, and nothing about what the program's bodies are. Programs on it: cell tracking (cell\_\allowbreak{}tracking\_\allowbreak{}table.md) and RSNA knee MRI.

\section{\texorpdfstring{What to work toward, in order}{What to work toward, in order}}

From the build plan (\hyperref[chap:build-plan]{build\_\allowbreak{}plan.md}):

\begin{enumerate}
\item \textbf{The seal} (M13): stage B (history, bodies, flattened, key, schedule), then the universal root and its trust anchor. The RSNA converter is done (M9); the export follows.
\item \textbf{Errors in every module} (M12), and its two fail-closed defects.
\item \textbf{The scale audit}, row by row: the entropy window into the request, ruling (a) at the three guards, the residual's width from the orders, the marginal's limits reported instead of set.
\item \textbf{Orders from the per-axis coherence reading} (M2), now built as M18 (\texttt{period\_\allowbreak{}read}, A12), and orders come from the data (order = period); the knee pulls it.
\item \textbf{The compression floor, by noise functionals} (M8, M9, M5) is its own theory: its plan and order of work are the compression workbook's, compression\_\allowbreak{}table.md.
\item \textbf{The algebra as passes}: the ladder-keeping sweep (A1), the LCA probe partition and the DFS-range overlap (A2), moments at every level (A3), C as the one pass (A4), the marginal over the whole component (A5).
\item \textbf{The port into orior} (close the gaps before anything moves). These gaps stand:

\begin{itemize}
\item (a) The layout is pinned to a commit hash before they re-check it.
\item (b) Our \texttt{arithmetic/\allowbreak{}no\_\allowbreak{}rounding} arms, \texttt{render} and \texttt{nbody/\allowbreak{}orior} carry their full doc text and are byte-identical to their files (their code is identical). Our \texttt{exact\_\allowbreak{}integer.\allowbreak{}\{c,h\}} is ahead of theirs: it has the open width, the multiplication ladder, the division and Newton's division (M1), and theirs multiplies by long multiplication alone and has no division. This tree is the source their engine is to be copied from (orior). The root \texttt{bench/\allowbreak{}} still differs from theirs in every file.
\item (c) \texttt{build\_\allowbreak{}engine.\allowbreak{}sh} and \texttt{build\_\allowbreak{}driver.\allowbreak{}sh} compile the exact integer from this tree's \texttt{arithmetic/\allowbreak{}no\_\allowbreak{}rounding} by default, and \texttt{ANCHOR\_\allowbreak{}EXACT\_\allowbreak{}ROOT} only overrides, and moving their \texttt{no\_\allowbreak{}rounding} does not break our build.
\item (d) \texttt{engine/\allowbreak{}CMake\allowbreak{}Lists.\allowbreak{}txt}: every path is this tree's, the C-only warnings are kept off the CUDA sources, the CUDA arms join where a CUDA compiler is found, and \texttt{bench\_\allowbreak{}lattice} is built only where the compiler has C99 \texttt{\_\allowbreak{}Complex}. It builds and runs every bench on four configurations (M15).
\item (e) Their C builds host-only (CMake, MinGW, a Pi 5, WSL), but obsignatio, tessera and most modules are \texttt{.\allowbreak{}cu} only. Each module needs a host-only path, or CUDA optional per module.
\item (f) Their \texttt{src/\allowbreak{}engine/\allowbreak{}n\_\allowbreak{}body} holds older copies of modules we have (and some we retired). One definition, and the old copies retired in the same change.
\item (g) The sift and the exact integer are outside M12's stateful errors: wrap at the boundary or wire them.
\item (h) The \texttt{representation} package, the Python model of games, particles, sound, text and proteins, sits in the engine at \texttt{src/\allowbreak{}engine/\allowbreak{}python/\allowbreak{}representation/\allowbreak{}} (game/, particle/, atom/, constants/, structure/, sound/, exact.py). Proofs under \texttt{evidence/\allowbreak{}proofs/\allowbreak{}} and examples under \texttt{examples/\allowbreak{}} (any\_\allowbreak{}corpus, art, 0\_\allowbreak{}experimental, chemistry, crystallography) import it. The engine is not clear of domain code while \texttt{representation/\allowbreak{}game/\allowbreak{}chess.\allowbreak{}py} sits under it.
\item (i) \texttt{src/\allowbreak{}cu/\allowbreak{}CMake\allowbreak{}Lists.\allowbreak{}txt} sets \texttt{ORIOR\_\allowbreak{}ROOT} to the parent of \texttt{src/\allowbreak{}engine}, which is \texttt{src/\allowbreak{}}, one level short of the repository's root. The four tests it guards lie under the root's \texttt{test/\allowbreak{}engine/\allowbreak{}}, and its \texttt{if(EXISTS .\allowbreak{}.\allowbreak{}.\allowbreak{})} blocks skip \texttt{test\_\allowbreak{}steer}, \texttt{test\_\allowbreak{}adversarial}, \texttt{test\_\allowbreak{}arm\_\allowbreak{}agreement} and \texttt{test\_\allowbreak{}o2\_\allowbreak{}spawn} without a word. The bench targets name \texttt{.\allowbreak{}.\allowbreak{}/\allowbreak{}bench/\allowbreak{}}, a path under \texttt{src/\allowbreak{}}, and the tree's \texttt{bench/\allowbreak{}} sits at the root. The fix is the root two levels up from \texttt{src/\allowbreak{}engine}. Measured with the fix: a configure, build and grade as a tessera job, 890.067 s, exit 0, “[+] all graders passed”, with \texttt{test\_\allowbreak{}steer}, \texttt{test\_\allowbreak{}adversarial} and \texttt{test\_\allowbreak{}arm\_\allowbreak{}agreement} among the graders. \texttt{maint/\allowbreak{}engine/\allowbreak{}build\_\allowbreak{}engine.\allowbreak{}sh} neither builds nor runs \texttt{test\_\allowbreak{}o2\_\allowbreak{}spawn}: its targets and its graders leave it out, and the fix only makes it a target. Built by name as a tessera job, it passes: “0 case(s) failed”, exit 0, 2.286 s.

\begin{itemize}
\item \textbf{The descent}, in \texttt{test/\allowbreak{}engine/\allowbreak{}nbody/\allowbreak{}orior/\allowbreak{}test\_\allowbreak{}o2\_\allowbreak{}spawn.\allowbreak{}c}. No log₂ N floor applies to isolating one of N alignments: the case isolates its 2,029 in 9 and prints so. It is given the needle, reads every survivor against the needle's byte at each candidate offset (\texttt{steer\_\allowbreak{}truthy\_\allowbreak{}after}, \texttt{src/\allowbreak{}cu/\allowbreak{}engine/\allowbreak{}nbody/\allowbreak{}orior/\allowbreak{}orior\_\allowbreak{}steer.\allowbreak{}c}), and places the offset leaving the fewest standing, a tie going to the earlier candidate. Its probes are not answers to yes or no questions about an unknown target. They form a certificate, a set of offsets at which every other alignment differs from the needle, and no general bound above one holds for its size. The comments make no bound claim. \textbf{Prior art, cited from knowledge.} A smallest certificate is a smallest set cover: each offset covers the alignments that differ from the needle there. Placing the offset that prunes most is the greedy set cover algorithm, which finds a cover within a factor H(N − 1) ≤ 1 + ln(N − 1) of the smallest cover drawn from the same offsets (Johnson 1974; Lovász 1975; Chvátal 1979). The cost case counts the survivors entering each placed level: the pruning field is 3,925 over 9 levels, against 18,261 for a re-scan, and the barely pruning field 18,514 over 20 levels, against 40,580. 3,925 is 1.93 times the 2,029-alignment universe (Derived), matching the case's own note “near twice the universe”. Choosing each level's probe also reads its survivors once for each candidate offset, up to 20. The composition and the re-scan are counted per alignment on the same footing, and neither count is of bytes read.
\end{itemize}
\item Their proposed order: new modules land as siblings under their \texttt{src/\allowbreak{}engine/\allowbreak{}c}; the nbody modules then replace their \texttt{n\_\allowbreak{}body} copies; their own base/nbody/sims split comes last. Since gap (c) closed, that split no longer has to land with a change to our build.
\end{itemize}
\item \textbf{The two towers in the engine.} T and T⁻¹ emitted as record floors by the engine itself, not only by the tests, and a program then runs F ∘ T⁻¹ on the crystal as one stack, checked against \texttt{tower.\allowbreak{}cu}'s kernels lane for lane (A13; \hyperref[chap:vertical-time-compression]{vertical\_\allowbreak{}time\_\allowbreak{}compression.md}). \textbf{Built and proved}: \texttt{tower\_\allowbreak{}record\_\allowbreak{}lift} and \texttt{tower\_\allowbreak{}record\_\allowbreak{}lower} focus each floor from a ruleset (E4), and \texttt{test/\allowbreak{}engine/\allowbreak{}compiler/\allowbreak{}cycle/\allowbreak{}record\_\allowbreak{}tower\_\allowbreak{}test} has 209 checks, 0 failed (A13), 210 since it reads the stack limit (M10). Still open: F ∘ T⁻¹ in the crystal path, the gather across blocks, edges as record floors, and the kernels taking a ruleset.
\item \textbf{The record machine as a compiler} (the kernel is a compiler). Its compiled programs are named automata: finite, and of fixed width, and an automaton can be part of a UTM program's toolbox. A UTM program then sits one level up, with storage that grows between runs and automata from the toolbox as its steps (theory). Once the automata are emitted, the stack return (M10) is dropped. \textbf{Built and measured} (\texttt{src/\allowbreak{}cu/\allowbreak{}engine/\allowbreak{}analysis/\allowbreak{}cycle/\allowbreak{}}, the record program compiled):

\begin{itemize}
\item \textbf{The automaton.} At load, a program's steps become CUDA source. Each step's value is a local array of its own exact limbs, and each operation is a template over those limbs, the interpreter's own arithmetic with the widths fixed. NVRTC, loaded at run time (\texttt{nvrtc64\_\allowbreak{}130\_\allowbreak{}0.\allowbreak{}dll}, \texttt{libnvrtc.\allowbreak{}so.\allowbreak{}13}), compiles it to a cubin for the device's architecture. The cubin is loaded as a library and launched in the interpreter's place, and the build links nothing new.
\item \textbf{Kept.} A process finds a program again by its whole source. Across processes, a cache file holds the source and its cubin (\texttt{\$CYCLE\_\allowbreak{}CACHE}, else \texttt{\%LOCALAPPDATA\%\textbackslash{}cycle} or \texttt{\textasciitilde{}/\allowbreak{}.\allowbreak{}cache/\allowbreak{}cycle}), and it is used only where its source is this source, byte for byte. A compile took 72 to 1,604 ms, and a cache read 0.4 to 1.4 ms.
\item \textbf{The interpreter stays}, as the oracle and the fallback. A program of more than 1,024 steps, or one NVRTC refuses, runs on it. Past 4,096 unrolled limb operations a program's loops are left rolled. The cap came from the runs: the bitwise test's 4,204-step stack held NVRTC past 1.2 GB, and the order test's 12,293-step program past 13 GB, before each was stopped.
\item \textbf{Switches.} \texttt{CYCLE\_\allowbreak{}RECORD\_\allowbreak{}INTERPRET=1} keeps every program on the interpreter. \texttt{CYCLE\_\allowbreak{}RECORD\_\allowbreak{}CHECK=1} runs both on every launch and refuses the launch where their records or refusals differ (\texttt{ENGINE\_\allowbreak{}ERROR\_\allowbreak{}LOGIC}). \texttt{CYCLE\_\allowbreak{}RECORD\_\allowbreak{}REPORT=1} says on stderr where each program came from and how long it ran. \texttt{cycle\_\allowbreak{}record\_\allowbreak{}compiled} says which ran.
\item \textbf{Proved.} The ten record tests, each run compiled, then checked, then on the interpreter: 30 runs, 0 failed (speed 12, bitwise 25, boundary 49, coherence 15, table 16, divide 20, order 18, tower 210, gaussian 11, guide 12 checks). The check compared both machines' records and refusals on 58 launches, and every one was the same.
\item \textbf{Measured} (\texttt{test/\allowbreak{}engine/\allowbreak{}compiler/\allowbreak{}cycle/\allowbreak{}record\_\allowbreak{}speed\_\allowbreak{}test}, 4,194,304 lanes, the best of 5 runs timed on the host around \texttt{cycle\_\allowbreak{}record\_\allowbreak{}run}, RTX 3070):

{\footnotesize
\setlength{\tabcolsep}{6pt}
\begin{longtable}{>{\RaggedRight\arraybackslash}p{\dimexpr0.3799\linewidth-2\tabcolsep\relax}>{\RaggedRight\arraybackslash}p{\dimexpr0.2442\linewidth-2\tabcolsep\relax}>{\RaggedRight\arraybackslash}p{\dimexpr0.2059\linewidth-2\tabcolsep\relax}>{\RaggedRight\arraybackslash}p{\dimexpr0.1700\linewidth-2\tabcolsep\relax}}
\toprule
program & interpreted & compiled & times faster \\
\midrule
\endhead
chain: the Gaussian step's 8 floors, 18 steps & 3.818 ms & 0.272 ms & 14.0 \\
powers: x to x\textasciicircum{}8 and their sum, 15 steps, products to 8 limbs & 33.307 ms & 1.920 ms & 17.3 \\
mix: 6 rounds of xor, product, sum and wrap, then the ladder, 32 steps & 6.010 ms & 0.500 ms & 12.0 \\
\bottomrule
\end{longtable}
}
\item \textbf{The divisions compiled.} Quotient, remainder, gcd and exact quotient are the interpreter's own Knuth D, Euclid on it, and multiply and mask. Each holds its working values at its operands' own limbs, where the interpreter holds them at the register file's. The five dividing tests, each run compiled, then checked, then on the interpreter: 15 runs, 0 failed (divide 20, coherence 15, bitwise 25, boundary 49, tower 210 checks). The check compared 101 launches, refusals included, and every one was the same. The kernels' own time, summed over the checked launches: the tower's 37, 11.518 ms compiled against 401.126 ms interpreted; the boundary's 14, 17.155 ms against 471.493 ms. The divide test's 4,096 lanes of 160-bit numerators took 0.926 ms compiled against 9.553 ms, and its 512 lanes of 2,048-bit numerators 27.634 ms against 99.652 ms and 16.607 ms against 22.055 ms. The cost is the compile. The tower's dividing programs of 761 steps took 66.9 s and 69.5 s, and its 1,012-step program 122.9 s for 328,744 bytes of cubin, paid once a program and then read from the cache.
\item \textbf{Where a division runs.} Inlined into every step, the divisions made the compile the cost. Called, one shared body a width, they halved the compile and slowed the tower's kernels. Each way was built by nvcc for the record tests' 52 dividing programs into a cache folder of its own, and the tests ran checked against each folder, three ways back to back and twice. The interpreter's own time on every run was the control: about 400 ms for the tower, 490 ms for the boundary, 160 ms for the divide test and 213 ms for the speed test. Every launch was the same on both machines. The kernels' time summed over each test's checked launches, first run and second:

{\scriptsize
\setlength{\tabcolsep}{4pt}
\begin{longtable}{>{\RaggedRight\arraybackslash}p{\dimexpr0.1748\linewidth-2\tabcolsep\relax}>{\RaggedRight\arraybackslash}p{\dimexpr0.1520\linewidth-2\tabcolsep\relax}>{\RaggedRight\arraybackslash}p{\dimexpr0.1678\linewidth-2\tabcolsep\relax}>{\RaggedRight\arraybackslash}p{\dimexpr0.1748\linewidth-2\tabcolsep\relax}>{\RaggedRight\arraybackslash}p{\dimexpr0.1505\linewidth-2\tabcolsep\relax}>{\RaggedRight\arraybackslash}p{\dimexpr0.1802\linewidth-2\tabcolsep\relax}}
\toprule
the divisions & nvcc, 52 programs & tower, 37 launches & boundary, 14 & divide, 8 & speed, 15 (no division) \\
\midrule
\endhead
inlined & 1,541 s & 10.8, 11.3 ms & 17.0, 17.0 ms & 45.0, 45.8 ms & 12.9, 13.2 ms \\
called & 864 s & 61.7, 64.3 ms & 124.8, 111.9 ms & 36.7, 37.3 ms & 12.8, 12.7 ms \\
called, operands copied first & not summed & 103.1, 99.4 ms & 106.6, 101.7 ms & 39.5, 40.0 ms & 12.7, 12.3 ms \\
\bottomrule
\end{longtable}
}

2,897 of the 2,904 division sites in those programs divide by one limb, and a call saves every live register around it. Copying the operands out kept the steps' own arrays in registers and was slower still. nvcc's fast compile (\texttt{-\allowbreak{}-\allowbreak{}Ofast-\allowbreak{}compile=mid}) built the 1,012-step program in 21.7 s, and its tower launches took 241.0 ms in one run (the interpreter's 409.7 ms).
\item \textbf{Built: a divisor of one limb inlined.} It runs the interpreter's own one-limb path, a 64-bit division a numerator limb from the top, known at compile time (\texttt{cycle\_\allowbreak{}divide\_\allowbreak{}limb}, and \texttt{cycle\_\allowbreak{}exact\_\allowbreak{}quotient\_\allowbreak{}limb}, which refuses a rest or a quotient past its limbs). Any wider divisor stays the one call (\texttt{cycle\_\allowbreak{}divide\_\allowbreak{}long}, \texttt{cycle\_\allowbreak{}exact\_\allowbreak{}quotient\_\allowbreak{}inverse}, \texttt{cycle\_\allowbreak{}gcd}). The ten record tests, each run compiled, then checked, then on the interpreter: 30 runs, 0 failed, and the check compared 128 launches, every one the same. Each test's NVRTC compiles and its checked launches summed, against the runs of the inlined and the called builds, with the interpreter's time on the same run:

{\footnotesize
\setlength{\tabcolsep}{6pt}
\begin{longtable}{>{\RaggedRight\arraybackslash}p{\dimexpr0.1742\linewidth-2\tabcolsep\relax}>{\RaggedRight\arraybackslash}p{\dimexpr0.3279\linewidth-2\tabcolsep\relax}>{\RaggedRight\arraybackslash}p{\dimexpr0.2502\linewidth-2\tabcolsep\relax}>{\RaggedRight\arraybackslash}p{\dimexpr0.2477\linewidth-2\tabcolsep\relax}}
\toprule
build & tower: compiles, kernels (interpreter) & boundary & divide \\
\midrule
\endhead
inlined & 706.9 s, 11.518 ms (401.126 ms) & 71.1 s, 17.155 ms (471.493 ms) & 20.6 s, 45.710 ms (161.301 ms) \\
called & 346.7 s, 49.179 ms (311.450 ms) & 35.2 s, 105.951 ms (513.240 ms) & 23.9 s, 39.030 ms (162.370 ms) \\
one limb inlined & 462.3 s, 14.782 ms (402.223 ms) & 44.1 s, 17.719 ms (462.862 ms) & 23.8 s, 40.696 ms (156.173 ms) \\
\bottomrule
\end{longtable}
}
\item \textbf{The compile grows with about the square of the steps.} The 1,012-step program cut to its first 128, 256, 512, 768 and 1,012 steps, each step's value kept live by one xor: cicc, NVRTC's front end and optimizer, took 2.2, 4.3, 18.1, 40.2 and 97.2 s, and 1.7, 3.2, 10.2, 20.5 and 52.4 s with its optimizer off. On the 381-step program ptxas took 0.6 s against cicc's 7.5 s. The cost follows the kernel's length, with or without the optimizer, and the loop below answers it where no flag does.
\item \textbf{No step cap.} A program of any length compiles. The interpreter runs a program only where one of its steps is not held by the compiler, where NVRTC or nvJitLink cannot be loaded, or under \texttt{CYCLE\_\allowbreak{}RECORD\_\allowbreak{}INTERPRET=1}.
\item \textbf{The operator block} (every program's operators split into one file that all programs access, for a faster and smaller build). Every record operation is one function in one module: the interpreter's own arithmetic, with its widths taken as arguments. NVRTC compiles the module once for the device as relocatable code, and the cache keeps it (1.4 s, 164,864 bytes of cubin). A program's source is its steps alone, each one call into the block with its places and widths as constants. NVRTC compiles that as relocatable code, and nvJitLink links it against the block (loaded at run time, \texttt{nv\allowbreak{}Jit\allowbreak{}Link\_\allowbreak{}130\_\allowbreak{}0.\allowbreak{}dll}, \texttt{libnv\allowbreak{}Jit\allowbreak{}Link.\allowbreak{}so.\allowbreak{}13}). The register file and the signs lie in the program's local frame, and each operator reads its operands and writes its register there by address. The divisions and the ladder share one scratch, laid at the program's widest division.
\item \textbf{The compile, with and without the block} (one program each):

{\footnotesize
\setlength{\tabcolsep}{6pt}
\begin{longtable}{>{\RaggedRight\arraybackslash}p{\dimexpr0.2406\linewidth-2\tabcolsep\relax}>{\RaggedRight\arraybackslash}p{\dimexpr0.2447\linewidth-2\tabcolsep\relax}>{\RaggedRight\arraybackslash}p{\dimexpr0.1666\linewidth-2\tabcolsep\relax}>{\RaggedRight\arraybackslash}p{\dimexpr0.1290\linewidth-2\tabcolsep\relax}>{\RaggedRight\arraybackslash}p{\dimexpr0.2191\linewidth-2\tabcolsep\relax}}
\toprule
build & program & compile & link & cubin \\
\midrule
\endhead
stage 1, the file in registers & tower, 761 steps & 37.9 s & none & 133,288 bytes \\
stage 1, the file in registers & tower, 1,012 steps & 96.4 s & none & 231,976 bytes \\
stage 1, loops rolled & tower, 2,094 steps & 778.7 s & none & 1,450,024 bytes \\
the operator block & coherence, 164 steps & 0.17 to 0.22 s & 0.7 to 1.3 ms & 91,360 to 96,480 bytes \\
the operator block & boundary, 968 steps & 1.29 s & 2.7 ms & 366,560 bytes \\
the operator block & boundary, 1,992 steps & 5.11 s & 4.1 ms & 731,744 bytes \\
the operator block & bitwise, 4,204 steps & 9.67 s & 5.4 ms & 1,201,888 bytes \\
the operator block & order, 12,293 steps & 217.5 s & 27.1 ms & 4,592,480 bytes \\
the operator block & order, 16,379 steps & 209.5 s & 24.9 ms & 4,747,616 bytes \\
\bottomrule
\end{longtable}
}

The bitwise stack of 4,204 steps and the order test's 12,293 steps both compile. Past about 4,000 steps the compile grows faster than the steps, and why is not yet measured.
\item \textbf{Proved on the block.} Each test run compiled, then checked against the interpreter, then on the interpreter. Every checked launch had the same records and refusals on both machines. The kernels' time summed over each test's checked launches, beside the interpreter's on the same run:

{\footnotesize
\setlength{\tabcolsep}{6pt}
\begin{longtable}{>{\RaggedRight\arraybackslash}p{\dimexpr0.2011\linewidth-2\tabcolsep\relax}>{\RaggedRight\arraybackslash}p{\dimexpr0.1660\linewidth-2\tabcolsep\relax}>{\RaggedRight\arraybackslash}p{\dimexpr0.1869\linewidth-2\tabcolsep\relax}>{\RaggedRight\arraybackslash}p{\dimexpr0.2045\linewidth-2\tabcolsep\relax}>{\RaggedRight\arraybackslash}p{\dimexpr0.2415\linewidth-2\tabcolsep\relax}}
\toprule
test & checks & checked launches & compiled & interpreted \\
\midrule
\endhead
bitwise & 25 & 6 & 8.566 ms & 21.748 ms \\
boundary & 49 & 16 & 94.854 ms & 438.492 ms \\
coherence & 15 & 37 & 14.459 ms & 143.879 ms \\
divide & 20 & 8 & 42.628 ms & 169.002 ms \\
gaussian & 11 & 1 & 0.264 ms & 2.138 ms \\
guide & 12 & 1 & 0.330 ms & 2.357 ms \\
order & 18 & 2 & 19.827 ms & 87.672 ms \\
\bottomrule
\end{longtable}
}

The cost is the call. The register file lies in memory and the operators read it by address. The boundary's kernels ran 4.6 times the interpreter's speed, against 26 times when the file was in registers (17.719 ms against 462.862 ms, the one-limb-inlined build above). The speed, table and tower tests, and the block built as LTO-IR and linked with link-time optimization (\texttt{CYCLE\_\allowbreak{}RECORD\_\allowbreak{}LTO=1}), are running.
\end{itemize}
\item \textbf{Resident programs.} Automata that run on the device on their own. Each reports its status to its own block and loops, talks to other programs through their blocks, and can be chained to another program or pointed at one. Built from these, a solution is small autonomous kernels composed on the GPU, with the host out of the loop.

\begin{itemize}
\item \textbf{The block} (\texttt{Engine\allowbreak{}Program\allowbreak{}Block}, \texttt{engine\_\allowbreak{}config.\allowbreak{}h}). Its core: the offset (the execaddr, where the program resumes), the step (the evacaddr, where it left), the span (the register map), and the time to live, the watchdog, the launch time, the run time and the exec time. Beside those it holds:

\begin{itemize}
\item the program's signum, the run's generation, and the launch that owns it;
\item the scheduler's command (run, yield, stop) and the grant: registers a thread, threads, device bytes (the local frame times the threads) and shared bytes (\texttt{engine\_\allowbreak{}config.\allowbreak{}h});
\item the state, the exits, and the error's module and site;
\item the check-ins, and the progress: lanes and refusals;
\item the wiring: the blocks it reads and the blocks that read it, each link's words produced and consumed, the result, and the parent;
\item a CRC-64 that seals the block whenever no launch holds it.
\end{itemize}
\item \textbf{Yield.} A launch runs its time to live and leaves, and its block holds where it stands. The next launch resumes from there, and no launch outlives Windows' 2 s watchdog.
\item \textbf{The exits.} True, false, malformed (no lattice was built), and answered: an answer table on the device, keyed by the signum of the program and its inputs, holds each answer once. An exit short-circuits every program reading the block.
\item \textbf{Two machines pointed at each other.} Running a machine that halts only on yes against one that halts only on no is dovetailing. It decides exactly the questions whose yes side and no side are both semi-decidable (Post~\cite{src:Post-1944}, 1944). A network of such processes wired through channels is a Kahn process network.
\item \textbf{What it must hold} (the adversarial pass):

\begin{itemize}
\item Programs that wait on each other must be resident together. One cooperative launch guarantees that, and separate launches do not.
\item One fault ends every program in its CUDA context.
\item Programs in separate processes cannot share device blocks under WDDM.
\item A block's data is fenced before its state is written.
\item The answer table's key is the program and its inputs, the full signum, confirmed on a hit.
\item Admission reads the compiled registers.
\item Wiring is data in the block and is never compiled in.
\item A link needs a bound or a handshake.
\end{itemize}
\item \textbf{Stages.}

\begin{enumerate}
\item The block, the register file in registers, and yield and resume, proved on the eleven record tests.
\item Resident programs as one cooperative launch, chained through their blocks.
\item The exits and the answer table.
\item Tessera admits with the compiled grant and reads the blocks from host memory.
\end{enumerate}
\item \textbf{Stage 1, built} (\texttt{src/\allowbreak{}cu/\allowbreak{}engine/\allowbreak{}analysis/\allowbreak{}cycle/\allowbreak{}}). The block, yield and resume, and the register file in registers: every place a constant in the source, and file\_\allowbreak{}limbs registers hold the file. The thread blocks take lanes a round at a time from one counter and leave past the time to live (500 ms unless \texttt{CYCLE\_\allowbreak{}RECORD\_\allowbreak{}TTL} names one). The block's command can also send them out. The thread block that opens a launch always runs one round, and every launch moves the program on. The last one out writes the offset, the times and the state, and the host launches again while the state is yielded. On the tower test the programs hold 128 to 255 registers a thread with local frames of 0 to 752 bytes. Each runs in one launch and two check-ins. The 761-step program holds 183 registers, a 0-byte frame, and 0.442 ms for 256 lanes. Stage 1's compiles take 10.2 to 778.7 s a program, which the operator block (item 9) answers.
\item \textbf{The block in device memory.} Check-ins stay on the device. The host lays the block, seals it and sends it before the first launch. It reads the block back once a launch has ended, checks the seal before the next launch, and seals and sends it again. A fault is left sealed on the host's copy.
\item \textbf{Still owed for stage 1:} a forced yield, \texttt{CYCLE\_\allowbreak{}RECORD\_\allowbreak{}TTL=1} on the eleven record tests, matching the interpreter. The eleventh, the lane test, comes with \texttt{ENGINE\_\allowbreak{}RECORD\_\allowbreak{}LANE}.
\end{itemize}
\item \textbf{The emitter over any ruleset.} A record program's lane written directly as NVIDIA PTX (\texttt{src/\allowbreak{}cu/\allowbreak{}transpiler/\allowbreak{}codegen/\allowbreak{}}, being built), with PTX's own rules held as data: its carry chains, its register classes and its calling convention, the construction set. Once this works it is split into functionals and support expanded to intel, amd, arm, riscv, vhdl and the rest, to prove it a universal compiler where it knows the hardware ruleset.

\begin{itemize}
\item \textbf{(a) Functionals.} Each operation's emitter is a functional over a target ruleset L. PTX is the first L. \textbf{The bootstrap test} (future work, not built): the emitter written as a record program and compiled by the emitter to PTX, then that PTX run on the same record program. The test passes when the two emitted texts match byte for byte, the fixed point a three-stage compiler bootstrap compares (GCC's bootstrap, from knowledge). It needs two things first: the operator block emitted as PTX by the same construction set, and the emitter expressible as a record program.

\begin{itemize}
\item \textbf{The bootstrap, built.} The device writes a lane's text with the text program, itself a record program run a lane a byte, and \texttt{codegen\_\allowbreak{}test} includes the bootstrap: the text program's own lane written by the device with the text program. The production route does the same: with \texttt{CODEGEN\_\allowbreak{}DEVICE=1} the device writes the assembly printer's own lane too, and while it does, the assembly printer's own program runs on the interpreter, the one program not compiled (\texttt{s\_\allowbreak{}cycle\_\allowbreak{}codegen\_\allowbreak{}depth} in \texttt{src/\allowbreak{}cu/\allowbreak{}engine/\allowbreak{}analysis/\allowbreak{}cycle/\allowbreak{}cycle\_\allowbreak{}compile\_\allowbreak{}route.\allowbreak{}cu}). A text the device writes and matches once in a process is not written again. The production route has not run.
\item \textbf{The emitter on the device, a reading} (a reading; built below). The emitter can recompile itself to run on the device. Emission is parallel by step: lane l reads step l from the step table and writes its text, with the templates as table data from the \texttt{.\allowbreak{}krs}, and a scan of the text lengths places each step's text. The analysis across steps, first and last puts, lifetimes and the live estimate, is scans and minima over lanes. On NVIDIA the device can write PTX, but assembly to SASS is the driver's, on the host, and a kernel cannot branch into code it wrote (from knowledge): self-hosting there leaves the host assembling and loading, and the host becomes a loader only once a SASS \texttt{.\allowbreak{}krs} exists, derived by the probes of stages 4 and 5. MaxAs and CuAssembler recover SASS encodings by probing (cited from knowledge; the sources are unread). A CPU process can write code and run it after remapping its pages executable, as a JIT does, where the platform allows it, and that makes stage 6's second target the likely first to host itself fully.

\begin{itemize}
\item \textbf{Built} (everything on the device). The emitter's decisions, keymath's imprint and key\_\allowbreak{}schedule's lay are each one source the host and the device compile (\texttt{codegen\_\allowbreak{}core.\allowbreak{}h}, \texttt{keymath\_\allowbreak{}core.\allowbreak{}h}, \texttt{key\_\allowbreak{}schedule\_\allowbreak{}core.\allowbreak{}h}). The device lays a program's step table, decides each step's forms a thread a step, cuts the lane into states in one thread, and writes the lane's text with the text program (\texttt{asm\_\allowbreak{}printer.\allowbreak{}h}, \texttt{codegen\_\allowbreak{}device.\allowbreak{}h}). \texttt{CODEGEN\_\allowbreak{}DEVICE=1} has the record machine's lanes written by the device and held to the host's text byte for byte. \textbf{Measured}, each device run a tessera job on Windows (RTX 3070): \texttt{codegen\_\allowbreak{}test} 121 checks, 0 failed, the bootstrap included, the emitter's text in PTX 170,922, C 412,182 and VHDL 554,754 bytes; \texttt{CODEGEN\_\allowbreak{}DEVICE=1} \texttt{record\_\allowbreak{}bitwise\_\allowbreak{}test} 25 checks, 0 failed, 8 programs written by the device byte for byte the host's, the largest 4,204 steps and 2,051,022 bytes. Not exercised: \texttt{engine.\allowbreak{}cu}'s \texttt{engine\_\allowbreak{}record\_\allowbreak{}lay\_\allowbreak{}device}.
\end{itemize}
\item \textbf{Tessera emitted for the device, a want} (a want; its ledger built on the device, below). Tessera is to run on the device and be emitted as soon as the compiler can. Tessera is host C (\texttt{tessera\_\allowbreak{}run.\allowbreak{}c} and the daemon's \texttt{tessera\_\allowbreak{}*.\allowbreak{}c} files). Once the emitter can express it, tessera becomes a program the emitter compiles for the device, and its decisions run where the programs run. This ties item 10's stage 4, “Tessera admits with the compiled grant and reads the blocks from host memory”, to the emitter on the device above. The part that can move is the ledger's arithmetic and the decisions over the blocks: each block's command (run, yield, stop) and its grant (registers a thread, threads, device bytes and shared bytes, the registers read from the compile), read and written by a resident program the way item 10's programs talk through their blocks. Three parts stay on the host, a reading. Admission watches operating-system processes, their tree, their signals and their peaks in processors and bytes, and the device sees none of it. A program enters the device through the host's loader (above). One fault ends every program in its CUDA context (item 10, what it must hold), and a scheduler in that context dies with the program it watches. Tessera is meant to be part of the compiler itself. \textbf{Tessera on every target} (built, below): the daemon, \texttt{tessera\_\allowbreak{}run} and the seal build with no CUDA toolchain, the Pi first, and runs there become tessera jobs.

\begin{itemize}
\item \textbf{Tessera on every target, built.} Where \texttt{\_\allowbreak{}\_\allowbreak{}CUDACC\_\allowbreak{}\_\allowbreak{}} is not defined \texttt{obsignatio.\allowbreak{}cu} compiles as C++. The functions the host and the device share are marked \texttt{OBSIGNATIO\_\allowbreak{}SHARED}, and with no CUDA toolchain \texttt{obsignatio\_\allowbreak{}many} and \texttt{obsignatio\_\allowbreak{}bits} refuse every request as \texttt{ENGINE\_\allowbreak{}ERROR\_\allowbreak{}RESOURCE}. A POSIX part with no nvcc builds the seal as C++ and links the daemon and \texttt{tessera\_\allowbreak{}run} with c++; the measure and job tests, which take device memory, are not built there, and say so. \textbf{Measured:} on a Raspberry Pi 5 (Linux aarch64, no CUDA toolchain, jobs given 2 of its 4 cores), \texttt{obsignatio\_\allowbreak{}test} host only, host vectors 210 cases, level keys 18, seal 778, device refused 4, each 0 failed; tessera's tests exit 0, and \texttt{tessera\_\allowbreak{}run} 15 checks, 0 failed. On Windows, each a tessera job, \texttt{obsignatio\_\allowbreak{}test}'s 15 tallies at 0 failed, measure 9 cases and job 25 checks at 0 failed, \texttt{tessera\_\allowbreak{}run} 14 checks, 0 failed. On WSL with no nvcc on the non-interactive path, \texttt{obsignatio\_\allowbreak{}test} host only at 0 failed, \texttt{tessera\_\allowbreak{}run} 15 checks, 0 failed; on a loaded machine two peaks of 2.149 and 2.116 processors against the 2.100 bound give 15 checks, 2 failed, and on a quieter machine they read 1.979 and 1.915.
\item \textbf{Tessera's ledger on the device, built} (there should be two). The ledger's decisions are one source (\texttt{tessera\_\allowbreak{}ledger\_\allowbreak{}core.\allowbreak{}h}): the host daemon's ledger grows its rooms and runs the core, with no CUDA context and the Pi build unchanged; the device's tessera (\texttt{tessera\_\allowbreak{}device.\allowbreak{}\{h,cu\}}) runs the same calls in one thread over rooms in device memory. \textbf{Measured}, each a tessera job on Windows: the daemon's \texttt{run.\allowbreak{}sh} exit 0 (\texttt{tessera\_\allowbreak{}ledger\_\allowbreak{}test} through the core 0 failed, job test 25/0, \texttt{tessera\_\allowbreak{}run} 14/0); \texttt{tessera\_\allowbreak{}device\_\allowbreak{}test} 601 checks, 0 failed, 200 rounds of 400 calls, every round growing the device's rooms.
\end{itemize}
\end{itemize}
\item \textbf{(b) Other targets}, each its own L: x86 (Intel, AMD), the AMD GPU instruction set, ARM (AArch64), RISC-V, and VHDL, where the program becomes a circuit.
\item \textbf{(c) The universal compiler and its proof.} Given L, the program compiles to it. The proof is the same program on every target matching the host word for word, the way the record tests hold the compiled path to the host and the interpreter (item 9).
\item \textbf{(d) Deriving L.} The engine helps a target derive its own ruleset, and the target's questions are then coherent.
\item \textbf{(e) Unknown L.} Where ruleset L is unknown, more and more basic questions are asked, fully unbound. A small instance is being built: the emitter holds no constant for PTX's \texttt{.\allowbreak{}version} or \texttt{.\allowbreak{}target}. It compiles an empty kernel with NVRTC and reads back the header, and the part of L it does not know it gets by asking the toolchain.

\begin{itemize}
\item Either a target's registers and their operators, its opcode or assembly, or a partial description of what the hardware is believed to be.
\item Given a simple set of math, the answers are checked and how the target works is known at once.
\item Illegal operations can be asked about: what crashes the process, since some systems stop a process that performs an operation their system holds absurd.
\item \textbf{Three ways into L, a reading.} The first statement names three ways in: registers with their operators, the opcode list or assembly, or a partial description of the hardware. PTX gives the emitter registers, operators and assembly. The question it still asks is the NVRTC header probe above. Status: not built.
\item \textbf{The membership query, a reading.} The second statement is a membership query: a small sum answered by the target's own hardware returns its rule. Take q = x / 0 on integers. x86 raises \#DE, delivered on Linux as SIGFPE, and INT\_\allowbreak{}MIN / −1 raises it too. AArch64's SDIV and UDIV return 0. RISC-V's M extension returns all ones for the quotient and x for the remainder. PTX's div leaves the value unspecified. One question separates the three rulesets whose manuals define the answer. PTX's manual leaves it undefined, and the emitter treats undefined as refused. The value a given GPU returns can equal one of the other three, and an observed value cannot pick PTX out. Cited from knowledge, with the sources unread: the Intel SDM, the Arm ARM, the RISC-V unprivileged specification's division table and the PTX ISA. The record machine refuses a zero divisor (\texttt{cycle\_\allowbreak{}record\_\allowbreak{}kernel.\allowbreak{}cu} and \texttt{cycle.\allowbreak{}c}). The emitter's own test for a zero divisor before each \texttt{div.\allowbreak{}u32} and \texttt{div.\allowbreak{}u64} is not built. Measured (11(f) 3): 7 / 0 faults on x86 and answers 0 on AArch64, the signed division SDIV; RISC-V is not measured.
\item \textbf{Illegal operations, a reading.} The third quote is the negative half of L: the operations the target refuses by killing the asker. The record machine holds its own half already, since a refused lane is a question L does not answer. The target's half has to be asked from a process that can die. A crash is an answer, and the signal or exit status names the rule broken. A tessera job has the same structure. On CUDA an illegal address leaves the context unusable, and the unit of a probe has to be a process. SIGILL on an instruction the part lacks is how an unknown CPU reports which extension sets it has. Cited from knowledge, with the sources unread.
\item \textbf{More basic questions, a reading.} When a question kills the process, the next question has to be more basic and asked from a fresh process. The cell that watches the probes, and the limit on what it knows, are posit 12 (\hyperref[chap:wants]{wants.md}).
\end{itemize}
\item \textbf{(f) The cell's machinery, in build order.} Step 1's stage 1 and step 3 are built, each with its status below. The cellular machinery for this universal compiler is called wilder than what von Neumann envisioned by decades of orders of magnitude, a comparison no figure in the tree measures.

\begin{enumerate}
\item \textbf{The PTX emitter} (item 11). Stage 1 is built and proved. Every PTX launch is checked against the interpreter on the record tests, and the PTX is then compared with the C path. \textbf{Both paths held, a reading.} The PTX path and the C source path both stay, and every program is measured on both. The relation between the two is the data, and rule (i) is set from it. Owed from this stage:

\begin{itemize}
\item (i) a rule that routes each program by the measured relation of the two paths and keeps both paths in the engine. The first quantity measured to separate them is the local frame against the device's stack limit ((iii), the frame against the stack limit). After the PTX build and its hold, the rule reads the kernel's local frame and the stack limit (\texttt{cuda\allowbreak{}Limit\allowbreak{}Stack\allowbreak{}Size}). A frame within the limit runs as PTX. Past the limit the rule builds the C source as well, and the smaller frame runs: the C source's wherever it fits and the PTX's does not. The other build's hold is given back (\texttt{cycle\_\allowbreak{}program\_\allowbreak{}release}). The PTX runs where the C source does not build, where its frame cannot be read, or where the two frames are equal. Routing a launch by its lanes is later work. Status: built (\texttt{cycle\_\allowbreak{}record\_\allowbreak{}route} and \texttt{cycle\_\allowbreak{}program\_\allowbreak{}release} in \texttt{cycle\_\allowbreak{}compile\_\allowbreak{}route.\allowbreak{}cu}) and proved. \textbf{Measured.} The eleven record tests ran, each a tessera job under \texttt{CYCLE\_\allowbreak{}RECORD\_\allowbreak{}CHECK=1} against a fresh cache, once as PTX and once as C source. Both had 0 failed (guide 12, coherence 21, boundary 59, gaussian 16, bitwise 25, divide 20, order 18, table 16, tower 210, speed 12 and lane 17 checks), and in each all 687 checked launches gave the same records and refusals as the interpreter. Rule (i) sent 7 program loads to C, on boundary (1,992 steps) and tower (4,537 to 6,510 steps): their PTX frames were 1,120 to 1,576 bytes against the 1,024-byte limit, and their C frames 176 to 184 bytes. The routing is proved. Its time is not measured apart from the device: it saves the stack growth of 7.449 to 10.130 ms a run ((iii), the frame against the stack limit), paid at a run's first launch and outside the kernel sums;
\item (ii) a thread count chosen per launch, for a launch of a few thousand lanes to reach every SM and not only a few blocks' worth. Status: built (\texttt{cycle\_\allowbreak{}record\_\allowbreak{}resident} in \texttt{cycle\_\allowbreak{}record\_\allowbreak{}launch.\allowbreak{}cu}) and proved with stage 1;

\begin{itemize}
\item \textbf{The launch rule against the device's state} (Measured). The boundary test ran as C (\texttt{CYCLE\_\allowbreak{}RECORD\_\allowbreak{}NVRTC=1}, \texttt{CYCLE\_\allowbreak{}RECORD\_\allowbreak{}CHECK=1}, a fresh cache each run), alternating the launch rule with (ii)'s: old, new, old, new. The compiled launches' host-timed total over the 19 checked launches, in ms, was 123.635 and 172.857 under the old rule and 149.581 and 164.806 under the new, with 121.307, 123.950 and 129.188 in three runs of the new before them. Every run passed its 59 checks. The two rules' totals overlap, and at 1,048,576 lanes both print the same grid, 92 thread blocks of 256 threads at once. (ii) did not slow the C path. The 21:24 run that gave the C column below its 41.793 ran the interpreter in 375.689 ms. The eight C runs since took 553.082 to 844.936 ms on the interpreter, and the two of them on that code took 553.082 and 842.056. The device ran faster at 21:24 than in any run since, a reading, and what changed in its state is not known. The RTX 3070 also drives the display, and nvidia-smi showed no other compute process when checked. Times from different runs are read against each run's own interpreter total, a reading: the C column below and the C times in the two readings under (iii) come from the 21:24 run.
\end{itemize}
\item (iii) the early store: each record word stored once its last put lays it, not every word at the lane's end. Built and proved with stage 1. \textbf{Measured.} Each record test ran as its own tessera job under \texttt{CYCLE\_\allowbreak{}RECORD\_\allowbreak{}CHECK=1} against a fresh cache, lanes as PTX, and all eleven passed: guide 12, coherence 21, boundary 59, gaussian 16, bitwise 25, divide 20, order 18, table 16, tower 210, speed 12 and lane 17 checks, 0 failed. All 687 checked launches gave the same records and refusals as the interpreter. The kernels' time in ms, summed over each test's checked launches:

{\scriptsize
\setlength{\tabcolsep}{4pt}
\begin{longtable}{>{\RaggedRight\arraybackslash}p{\dimexpr0.1495\linewidth-2\tabcolsep\relax}>{\RaggedRight\arraybackslash}p{\dimexpr0.1359\linewidth-2\tabcolsep\relax}>{\RaggedRight\arraybackslash}p{\dimexpr0.1463\linewidth-2\tabcolsep\relax}>{\RaggedRight\arraybackslash}p{\dimexpr0.1325\linewidth-2\tabcolsep\relax}>{\RaggedRight\arraybackslash}p{\dimexpr0.2176\linewidth-2\tabcolsep\relax}>{\RaggedRight\arraybackslash}p{\dimexpr0.2182\linewidth-2\tabcolsep\relax}}
\toprule
test & PTX, the early store & PTX, the lane-end store & C source & PTX, both paths [its interpreter] & C, both paths [its interpreter] \\
\midrule
\endhead
tower & 147.265 & 172.080 & 912.861 & 138.824 [2397.359] & 286.545 [2151.482] \\
order & 25.414 & 61.358 & 276.034 & 6.825 [146.721] & 86.893 [250.027] \\
gaussian & 138.086 & 193.010 & 315.176 & 99.448 [2663.893] & 141.549 [2719.230] \\
speed & 17.656 & 19.948 & 65.380 & 17.313 [238.495] & 63.052 [252.659] \\
lane & 2.868 & 4.001 & 31.517 & 2.671 [54.044] & 9.916 [61.539] \\
coherence & 10.970 & 13.629 & 30.300 & 9.888 [133.787] & 14.070 [152.511] \\
bitwise & 1.870 & 3.257 & 23.563 & 1.606 [21.927] & 4.070 [22.768] \\
guide & 0.264 & 0.208 & 0.484 & 0.216 [1.461] & 0.652 [1.811] \\
boundary & 24.636 & 50.019 & 41.793 & 55.927 [872.781] & 126.399 [638.947] \\
divide & 106.697 & 479.441 & 111.690 & 102.468 [156.013] & 105.988 [167.289] \\
table & 2.044 & 3.849 & 2.096 & 1.838 [2.207] & 2.059 [2.103] \\
\bottomrule
\end{longtable}
}

PTX with the early store is faster than the C source in every test, and faster than the lane-end store in every test but guide. The C column is run in the faster device state (ii)'s comparison found. The two both-paths columns are stage 1's proof run (Measured), both paths under the same code, each with its run's interpreter total in brackets. There PTX is faster than C in all eleven tests, by 1.03 to 12.7 times (Derived: divide 1.034, table 1.120, gaussian and coherence 1.423, tower 2.064, boundary 2.260, bitwise 2.534, guide 3.019, speed 3.642, lane 3.712, order 12.732). Divide and table, which otherwise spill and run slower than C (479.441 against 111.690, 3.849 against 2.096), do not with the early store. The device's state moves between runs on the same code: boundary's interpreter total is 593.667 ms in the early store's run and 872.781 in the both-paths run. The columns therefore compare the two paths and do not isolate rule (i) from the device. The divide test's largest program, 5 steps over 512 lanes, holds 241 registers, a 616-byte local frame and 96.901 ms with the early store, against 255 registers, a 1,256-byte frame and 467.075 ms with the lane-end store. The order test's 12,293-step program holds at most 56 live words a lane, 68 registers and a 136-byte local frame, and does not spill, where the lane-end store holds 255 registers and a 16,696-byte frame. Its assembly and link take 70.8 s, against 440.5 s for the lane-end store.

\begin{itemize}
\item \textbf{The live estimate} (built). The most 32-bit words a lane holds live at once, reckoned from the lifetimes the PTX text gives ptxas: each value's limbs and sign from its step to its last reader, each record word from its first put to its last, each atom word from its first reader to its last, each step's own temporaries at that step, and the words every lane holds throughout. It is reported with each program, and the live-word counts in this item are its output.
\item \textbf{Live words against the two paths} (one row a launch, each launch's live-word estimate, registers, frame and lanes beside its PTX and C times). Of the 11 launches past 320 live words, PTX is slower on 4: bitwise at 348 (0.557 against 0.515 ms), divide at 385 (7.326 against 7.086), boundary at 369 (8.342 against 3.491) and tower at 370 (11.137 against 3.898). The other 7 run PTX 1.04 to 14.5 times as fast (Derived). Two tower launches alike in live words (370), registers (255), frame (1,144 bytes) and lanes (256) split: one ran 16.970 ms as PTX against 54.642 as C, the other 11.137 against 3.898. No budget of live words separates the two paths yet.
\item \textbf{The frame against the stack limit} (Measured: the early run's logs, one row a launch, its frame, registers and lanes beside its host-timed and device times and the gap between the two). On the 680 launches whose local frame is 960 bytes or less, the gap runs 0.127 to 1.630 ms. On the 7 whose frame is 1,120 to 1,576 bytes, it runs 7.449 to 10.130 ms. The device's stack limit, 1,024 bytes a thread (the tower test prints it before and after its blocks), falls between the two. A frame past the limit makes the runtime raise it at a run's first launch and reserve that frame for every resident thread, and \texttt{cycle\_\allowbreak{}stack\_\allowbreak{}return} gives the reservation back once the run is done, by design (\texttt{cycle\_\allowbreak{}launch.\allowbreak{}cu}: a 7,600-byte frame holds 467,668,992 bytes past the 1 KiB default, and setting the limit back returns them in 5.0 ms). The gap's jump at the limit is that growth and return, paid once a run, a reading. Boundary's 1,992-step program shows both sides: as PTX it ran 0.893 ms on the device and 8.342 ms host-timed, and as C 2.920 and 3.491.
\end{itemize}
\end{itemize}
\item \textbf{The ruleset as data.} PTX's rules move out of the emitter's code into a table the emitter reads: carry chains, register classes, the calling convention and the header. Each operation's emitter becomes a functional over L. This is (a)'s split.

\begin{itemize}
\item PTX output is one of the n rulesets the engine supports.
\item \textbf{Stage 2's plan, a reading} (a plan; 2a, 2b and 2c build it, below, except the canonical form and its signum identity). The emitter splits into a core that knows no target and a ruleset L it reads at run time. The core holds what every target shares: each record word's first and last put and the early store; the refusal sites and the words in flight; the value and atom-word lifetimes and the live-word estimate; the unrolling by width and the branch-free selects. L holds what a target owns: mnemonics, register classes and their declarations, labels, address syntax, the calling convention, the header, and the toolchain that builds it (nvJitLink for PTX, NVRTC for C). The C source path is a ruleset too, and stage 2 brings the two paths under one core. A ruleset maps each operation to a native form or to a construction from more basic ones: the carry is native in PTX (\texttt{add.\allowbreak{}cc}, \texttt{addc}), x86 (\texttt{adc}) and AArch64 (\texttt{adds}, \texttt{adcs}), and on RISC-V it is built from \texttt{sltu}, the compare stage 5 uses to the same end. The core moves to \texttt{src/\allowbreak{}cu/\allowbreak{}scaffolding/\allowbreak{}}, each L to a data file in \texttt{src/\allowbreak{}engine/\allowbreak{}compiler/\allowbreak{}codegen/\allowbreak{}rulesets/\allowbreak{}}, and the \texttt{cycle/\allowbreak{}} files keep the runtime: launch, share, check, block and cache. Rule (i) becomes routing among n rulesets, its first quantity the frame against the stack limit. The proof: the PTX ruleset writes PTX byte-identical to stage 1's, every program of the eleven record tests a cache hit, and the C ruleset passes the checked run.
\item \textbf{The ruleset file} (a plan, not built). The extension is \texttt{.\allowbreak{}krs}, for kolmogorov ruleset. A ruleset is plain ASCII, one entry a line, with the extension \texttt{.\allowbreak{}krs}. The first line is \texttt{krs 1}, the format and its version. The first word of each line after it names its kind: \texttt{header}, \texttt{toolchain}, \texttt{class}, \texttt{form}, or \texttt{construct}, whose block runs to a line \texttt{end}. \texttt{\#} starts a comment, and a template is the rest of its line, verbatim. The engine writes each file's canonical form itself, and a ruleset's identity is the obsignatio signum over the canonical bytes. The first two files are \texttt{ptx.\allowbreak{}krs} and \texttt{c.\allowbreak{}krs}. JSON, YAML, TOML and S-expressions are ruled out, and \texttt{.\allowbreak{}reg}, the Windows registry's extension.
\item \textbf{Stage 2a, the ruleset as data} (built and proved). There are n rulesets, each a \texttt{.\allowbreak{}krs} file whose first line is \texttt{krs 1}, read once a process from \texttt{src/\allowbreak{}engine/\allowbreak{}compiler/\allowbreak{}codegen/\allowbreak{}rulesets/\allowbreak{}} or from the folder \texttt{CYCLE\_\allowbreak{}RULESETS} names. \texttt{CYCLE\_\allowbreak{}RULESETS} names a folder, as \texttt{CYCLE\_\allowbreak{}CACHE} does. The six switches stay \texttt{CYCLE\_\allowbreak{}RECORD\_\allowbreak{}CHECK}, \texttt{CYCLE\_\allowbreak{}RECORD\_\allowbreak{}INTERPRET}, \texttt{CYCLE\_\allowbreak{}RECORD\_\allowbreak{}LTO}, \texttt{CYCLE\_\allowbreak{}RECORD\_\allowbreak{}NVRTC}, \texttt{CYCLE\_\allowbreak{}RECORD\_\allowbreak{}REPORT} and \texttt{CYCLE\_\allowbreak{}RECORD\_\allowbreak{}TTL}.

\begin{itemize}
\item The entry kinds. \texttt{ruleset <name>} names the target the file describes. \texttt{toolchain <name>} names what turns the text into a program the target loads (\texttt{nvjitlink} in \texttt{ptx.\allowbreak{}krs}). \texttt{header <name>} names where the text's opening lines come from (\texttt{probe\_\allowbreak{}nvrtc} in \texttt{ptx.\allowbreak{}krs}, the header asked of NVRTC, (e)). \texttt{bank <name> = <definition>} is a bank of registers, \texttt{\{n\}} the register's number in it. \texttt{fixed <name> = <definition>} is one register every lane holds throughout. \texttt{form <name> <parameter>.\allowbreak{}.\allowbreak{}.\allowbreak{} = <text>} is a template. \texttt{construct <name> <parameter>.\allowbreak{}.\allowbreak{}.\allowbreak{}} is a form built from more basic ones, its lines running to \texttt{end}, and the reader reads it (\texttt{ruleset\_\allowbreak{}core\_\allowbreak{}read.\allowbreak{}h} reads the construct kind).
\item The text rules. A form's text is the rest of its line after \texttt{=}, byte for byte, with \texttt{\textbackslash{}t} a tab, \texttt{\textbackslash{}n} a line's end and \texttt{\textbackslash{}\textbackslash{}} a backslash. \texttt{\{p\}} is parameter p's argument where p is one of the form's parameters, and any other brace stands for itself. A line that begins with \texttt{\#} is a comment, a blank line is nothing, and a carriage return at a line's end is dropped. A form is kept cut at its parameters, and writing one appends its pieces with each argument between them.
\item \texttt{ptx.\allowbreak{}krs} holds 9 banks (file, sign, out, atom, temporary, wide, predicate, member, immediate), 12 fixed registers and 88 forms in 170 lines.
\item \textbf{Stage 2c, the C source as a ruleset} (built and proved with 2a). \texttt{c.\allowbreak{}krs} holds 83 forms in 161 lines, with toolchain \texttt{nvrtc} and header \texttt{prelude} (the operator block's prelude, \texttt{s\_\allowbreak{}cycle\_\allowbreak{}prelude}, \texttt{target.\allowbreak{}h}). One reader serves n rulesets through a schema each language reads against. A refused \texttt{c.\allowbreak{}krs} leaves a program the PTX does not hold on the interpreter, with the reason “its C ruleset, c.krs, was refused”, and rule (i) keeps the PTX. \texttt{utils/\allowbreak{}test/\allowbreak{}src/\allowbreak{}cu/\allowbreak{}transpiler/\allowbreak{}codegen/\allowbreak{}ruleset\_\allowbreak{}read\_\allowbreak{}test.\allowbreak{}cu} reads every ruleset against its schema and holds the files each is read from against each other for absurd relations (\texttt{ruleset\_\allowbreak{}core\_\allowbreak{}relation.\allowbreak{}h}). The emitter writes through the ruleset, and the byte-for-byte match of a separate C printf is theory. 2a and 2c are proved together on the device (the proof, below).
\item \textbf{SASS} (built). \texttt{sass.\allowbreak{}krs} is NVIDIA SASS for sm\_\allowbreak{}86, the machine code the device runs, read back by the cell's probes. It holds 9 banks, 12 fixed registers, 69 forms and 2 constructs in 301 lines.
\item \textbf{Yosys} (built). \texttt{yosys.\allowbreak{}krs} is Yosys's generic synthesis of a lane GHDL wrote as Verilog, a script Yosys reads a command a line (\texttt{yosys -\allowbreak{}s}). It holds 5 forms in 22 lines.
\item The emitter writes no PTX of its own. It decides what a step does, and the ruleset defines it for the target. The PTX left in \texttt{cycle\_\allowbreak{}internal.\allowbreak{}h} is the interpreter's inline multiply-add (\texttt{mad.\allowbreak{}lo.\allowbreak{}cc.\allowbreak{}u32}, \texttt{addc.\allowbreak{}u32}), device code the emitter does not write.
\item Refusal. A ruleset that lacks a bank, register or form the emitter names, holds one it does not name, or gives a form other parameters is refused whole. A form given the wrong number of arguments marks the program broken. Either sends the program to the C source, and the reason is reported.
\item \textbf{The proof} (Measured). The eleven record tests ran as tessera jobs with \texttt{CYCLE\_\allowbreak{}RECORD\_\allowbreak{}CHECK=1}: once as PTX against a cache seeded from stage 1's final PTX run, and once as C source against one seeded from its final C run. A program whose text differs from stage 1's by a byte misses the cache and is written anew. None was, as PTX or as C: the PTX run read 165 programs from the cache and the C run 196, and no ruleset was refused. Both modes had 0 failed: guide 12, coherence 21, boundary 59, gaussian 16, bitwise 25, divide 20, order 18, table 16, tower 210, speed 12 and lane 17 checks. The PTX coherence test's first build stopped in the host compiler (\texttt{cl.\allowbreak{}exe} exit 0xC0000142) before any stage 2 code ran. It was run again alone, from the same snapshot with a cache seeded from the first PTX run: 21 checks, 0 failed, its 38 programs read from the cache and none written anew. The PTX ruleset writes PTX byte-identical to stage 1's, and the C ruleset writes stage 1's C source.
\item The built reader departs from the plan above in three places: \texttt{class} is \texttt{bank} and \texttt{fixed}; a template is read with escapes where the plan reads it verbatim; and the canonical form with its signum identity is not built. \texttt{c.\allowbreak{}krs} carries out the plan's second file.
\end{itemize}
\item \textbf{Stage 2b, the emitter in \texttt{src/\allowbreak{}cu/\allowbreak{}transpiler/\allowbreak{}vendor\_\allowbreak{}bin\_\allowbreak{}layouts/\allowbreak{}}} (built and proved). The emitter is \texttt{src/\allowbreak{}cu/\allowbreak{}transpiler/\allowbreak{}vendor\_\allowbreak{}bin\_\allowbreak{}layouts/\allowbreak{}}: \texttt{container\_\allowbreak{}layout.\allowbreak{}\{h,c\}} holds a container's layout as rows in a file, the offsets, widths, tags and names a format uses, which the emitter carries none of; \texttt{container\_\allowbreak{}write.\allowbreak{}\{h,c\}} is one emitter for every container, which reads a layout file and writes what that layout describes, carrying over untouched what it does not understand; and \texttt{nvidia/\allowbreak{}elf64\_\allowbreak{}nvidia.\allowbreak{}tsv} is one such layout, beside the cubin writer that reads it. The rulesets and the code generation are \texttt{src/\allowbreak{}cu/\allowbreak{}transpiler/\allowbreak{}codegen/\allowbreak{}}, the rulesets at \texttt{codegen/\allowbreak{}rulesets/\allowbreak{}}. The emitter reads no device and loads no library. The \texttt{cycle/\allowbreak{}} runtime names the target through a \texttt{Cycle\allowbreak{}Target} built by \texttt{cycle\_\allowbreak{}target} (the operator block's hash, the device's compute capability, NVRTC's version and the prelude), and compiles, links, caches and launches what the emitter writes. Every build of the engine compiles the emitter with it: \texttt{build\_\allowbreak{}engine.\allowbreak{}sh}, \texttt{build\_\allowbreak{}driver.\allowbreak{}sh}, \texttt{examples/\allowbreak{}qasm/\allowbreak{}build.\allowbreak{}sh}, \texttt{sims/\allowbreak{}run.\allowbreak{}sh} and the twelve record and residual tests' scripts.

\begin{itemize}
\item \textbf{The move.} The emitter is cut from the record machine into \texttt{src/\allowbreak{}cu/\allowbreak{}scaffolding/\allowbreak{}}: \texttt{static} is dropped from the five calls the record machine makes, the target passed in place of its own statics, with the ruleset folder and the include. Every file is ASCII.
\item \textbf{The proof} (Measured). The eleven record tests ran as tessera jobs with \texttt{CYCLE\_\allowbreak{}RECORD\_\allowbreak{}CHECK=1}: once as PTX against a cache seeded from stage 2's PTX runs, and once as C source against one seeded from its C run. No program was written anew, as PTX or as C: the PTX run read 203 programs from the cache, the 165 of stage 2's PTX run and the 38 of its coherence run again, and the C run read 196, and no ruleset was refused. Both had 0 failed: guide 12, coherence 21, boundary 59, gaussian 16, bitwise 25, divide 20, order 18, table 16, tower 210, speed 12 and lane 17 checks. The emitter moves without changing a byte of what it writes.
\item \textbf{The build check} (Measured). \texttt{examples/\allowbreak{}qasm/\allowbreak{}test/\allowbreak{}run.\allowbreak{}sh} checks qasm: \texttt{qasm\_\allowbreak{}test}, its exact test 39 checks and its self test 44 checks, 0 failed. \textbf{In the driver's build.} \texttt{examples/\allowbreak{}cell\_\allowbreak{}tracking/\allowbreak{}build\_\allowbreak{}driver.\allowbreak{}sh}'s \texttt{FUNCTIONALS} (\texttt{build\_\allowbreak{}driver.\allowbreak{}sh}) names \texttt{c/\allowbreak{}types/\allowbreak{}file\_\allowbreak{}defs/\allowbreak{}krep} and \texttt{c/\allowbreak{}transpiler/\allowbreak{}codegen}, with no emit entry and no \texttt{base/\allowbreak{}} path.
\end{itemize}
\item \textbf{The code generator and its language classes} (built). The code generator is a base class, \texttt{Code\allowbreak{}Generator} (\texttt{code\_\allowbreak{}generator.\allowbreak{}h}, the register lane that names no language), and each language inherits it and names its ruleset: \texttt{Ptx\allowbreak{}Target} (\texttt{ptx\_\allowbreak{}target.\allowbreak{}h}, \texttt{ptx.\allowbreak{}krs}), \texttt{CTarget} (\texttt{c\_\allowbreak{}target.\allowbreak{}h}, \texttt{c.\allowbreak{}krs}) and \texttt{Vhdl\allowbreak{}Target} (\texttt{vhdl\_\allowbreak{}target.\allowbreak{}h}, \texttt{vhdl.\allowbreak{}krs}). The \texttt{cycle/\allowbreak{}} module is cut into functional chunks, the split item 11's first statement asks for. \textbf{Measured}: the eleven record tests pass as PTX, 11 of 11 with 203 programs read from the cache, and as C source, 11 of 11 with 196, and no program is written anew.
\item \textbf{The test folders.} The engine's test folders sit out of \texttt{src/\allowbreak{}engine} under \texttt{test/\allowbreak{}engine}, which mirrors the category tree: \texttt{test/\allowbreak{}engine/\allowbreak{}runtime/\allowbreak{}obsignatio}, \texttt{examples/\allowbreak{}qasm/\allowbreak{}test}, \texttt{test/\allowbreak{}engine/\allowbreak{}runtime/\allowbreak{}daemon} and \texttt{test/\allowbreak{}engine/\allowbreak{}nbody/\allowbreak{}max\_\allowbreak{}tree} each run from there with 0 failed. The cell\_\allowbreak{}tracking README states that \texttt{.\allowbreak{}kcr} is its own compression format.
\item \textbf{The invariance theorem, a reading.} K\_\allowbreak{}U(x) ≤ K\_\allowbreak{}V(x) + c\_\allowbreak{}\{U,V\}, with c\_\allowbreak{}\{U,V\} the length of an interpreter for V written for U (Solomonoff 1964, Kolmogorov 1965, Chaitin 1969; Li and Vitányi, Theorem 2.1.1; cited from knowledge). The engine's device interpreter, \texttt{cycle\_\allowbreak{}record\_\allowbreak{}kernel} (\texttt{cycle\_\allowbreak{}record\_\allowbreak{}kernel.\allowbreak{}cu}), has that form for the record machine on the device. The emitter core with a ruleset L runs on the host and translates each record program into an L program whose length grows with the record program's, which bounds K\_\allowbreak{}L(x) by m K\_\allowbreak{}rec(x) + h (Derived), m the longest template a step and h the header: a multiplicative bound. The \texttt{.\allowbreak{}krs} file's canonical length enters the additive constant when the core and the ruleset are carried as an L program. Carried as an L program that runs the code it writes, the core and the ruleset give K\_\allowbreak{}L(x) ≤ K\_\allowbreak{}rec(x) + |core as L| + |krs| + c, c the driver that runs the emitter and jumps to its output (Derived), and the bootstrap's fixed point (11(a)) shows the core running as an L program. On NVIDIA the jump needs the host's assembler and loader, and the device interpreter keeps the additive form with no host step. A ruleset derived by probes is a measured description of an unknown machine.
\item \textbf{Across machines, a reading} (a plan, not built). Cross-hardware communication is then no problem from zero. Values: every ruleset compiles the same record program, and each target is proved by matching the host word for word. Two targets each proved on a program give the same bits for it, and equal obsignatio signums show it. Wire order: the record format fixes the byte order once, as little-endian 32-bit limbs, the way the seal already writes its lengths and offsets as LE64 (κ and Σ in the engine in math). Ordering: every \texttt{.\allowbreak{}krs} carries a memory-model section beside its arithmetic, defining acquire, release and a full fence at each scope. The core emits “publish the block” once, and each target defines it (from knowledge): x86 is ordered as TSO, where acquire and release need no fence and a full fence (\texttt{mfence} or a locked instruction) orders a store before a later load; AArch64 uses \texttt{dmb ish}, or \texttt{ldar} and \texttt{stlr}; RISC-V uses \texttt{fence} with its predecessor and successor sets; PTX uses scoped fences, \texttt{fence.\allowbreak{}acq\_\allowbreak{}rel.\allowbreak{}gpu} and \texttt{fence.\allowbreak{}acq\_\allowbreak{}rel.\allowbreak{}sys}. This carries item 10's rule, “A block's data is fenced before its state is written”, across targets. Matching the host proves a program's values; the order between programs rests on each target's fences, and a run that matches does not prove it, a reading. Posit 19 replaces this as the direction (stage 4).
\item \textbf{Prior art, cited from knowledge.} GCC's machine descriptions, each target's instruction patterns held as data the compiler's core reads. LLVM's TableGen target descriptions and its legalization, which expands an operation a target lacks into ones it has (an add with carry out becomes an add and \texttt{sltu} on RISC-V). C. W. Fraser and D. R. Hanson, \emph{A Retargetable C Compiler: Design and Implementation} (Addison-Wesley, 1995), lcc, whose back ends are tree grammars compiled by lburg. N. Ramsey and M. F. Fernández, “Specifying representations of machine instructions”, ACM TOPLAS 19 (1997) 492–524, the New Jersey Machine-Code Toolkit. For the memory models: L. Lamport, “How to make a multiprocessor computer that correctly executes multiprocess programs”, IEEE Transactions on Computers C-28 (1979) 690–691; P. Sewell, S. Sarkar, S. Owens, F. Zappa Nardelli and M. O. Myreen, “x86-TSO: a rigorous and usable programmer's model for x86 multiprocessors”, Communications of the ACM 53 (2010) 89–97; J. Alglave, L. Maranget and M. Tautschnig, “Herding cats: modelling, simulation, testing, and data mining for weak memory”, ACM TOPLAS 36 (2014), article 7; the RISC-V Instruction Set Manual, Volume I, its RVWMO memory model; D. Lustig, S. Sahasrabuddhe and O. Giroux, “A formal analysis of the NVIDIA PTX memory consistency model”, ASPLOS 2019.
\end{itemize}
\item \textbf{The cell, a probe runner} (posit 12). Each ribosome, a small probe program, runs in a child process the cell can lose. The cell records how each one ends (exit status, signal or fault, and output) and sees everything inside it. A tessera job has this shape already. Status: built (\texttt{src/\allowbreak{}cu/\allowbreak{}transpiler/\allowbreak{}lstar/\allowbreak{}interface/\allowbreak{}interface.\allowbreak{}\{h,c\}}, \texttt{interface\_\allowbreak{}probe\_\allowbreak{}run}, error module 23, \texttt{ENGINE\_\allowbreak{}MODULE\_\allowbreak{}INTERFACE}), not in the engine library build: its test builds it. \texttt{interface\_\allowbreak{}probe\_\allowbreak{}run} starts a probe, a small program that asks the target one question, in a child process, and records how it ended. On Windows the probe starts suspended in a job object that ends every process in it when the cell lets it go, with the fault dialogs off (\texttt{Set\allowbreak{}Error\allowbreak{}Mode}, inherited). On POSIX it runs in a process group of its own, and a pipe that closes on exec tells a probe that never started from one that exited. Its output and errors go to one file the cell reads back, its input is empty, and a probe past its limit is ended with every process it started. The answer is how it ended (exited, signaled, faulted, out of time, not started), its code, the rule a signal or a Windows NTSTATUS names (arithmetic, address, instruction, stack, trap, abort, other), its output and its wall time. An exit status with its top bit set on Windows is read as the fault that ended the process.

\begin{itemize}
\item \textbf{Measured.} \texttt{test/\allowbreak{}engine/\allowbreak{}compiler/\allowbreak{}interface/\allowbreak{}interface\_\allowbreak{}test.\allowbreak{}sh} builds the cell and \texttt{test/\allowbreak{}engine/\allowbreak{}compiler/\allowbreak{}interface/\allowbreak{}interface\_\allowbreak{}probe.\allowbreak{}c}, one question a process, built with no optimization and volatile operands, and holds each ending to the host's rules: 17 checks, 0 failed on Windows (MSVC, x86-64) and on WSL (gcc, x86\_\allowbreak{}64), each a tessera job. On Windows 7 / 0 faulted 0xC0000094, INT\_\allowbreak{}MIN / −1 0xC0000095, a read of address 16 0xC0000005, \texttt{ud2} 0xC000001D, \texttt{\_\allowbreak{}\_\allowbreak{}debugbreak} 0x80000003, unbounded recursion 0xC00000FD, and \texttt{abort()} 0xC0000409 (the C runtime's fail-fast). On WSL both divisions gave SIGFPE, the address and the recursion SIGSEGV, \texttt{ud2} SIGILL, \texttt{int3} SIGTRAP and \texttt{abort()} SIGABRT. A probe that hangs was ended at 0.522506 s on Windows and at 0.501591 s on WSL, by SIGKILL, against a 0.5 s limit. A program that is not there was not started: code 2 on Windows, errno 2 on WSL. On both, exit 0 and exit 3 exit as asked, the output and then the errors come back in one file (50 bytes on Windows, whose C runtime ends a line with two, 48 on Linux), a room of 8 keeps 7 bytes and counts the whole output, and after each death a fresh probe exits 0. The test holds each division to the part's own rule, chosen when it is built (\texttt{INTERFACE\_\allowbreak{}TEST\_\allowbreak{}DIVISION\_\allowbreak{}FAULTS}, \texttt{INTERFACE\_\allowbreak{}TEST\_\allowbreak{}ZERO\_\allowbreak{}QUOTIENT} and \texttt{interface\_\allowbreak{}test\_\allowbreak{}answered}, \texttt{test/\allowbreak{}engine/\allowbreak{}compiler/\allowbreak{}interface/\allowbreak{}interface\_\allowbreak{}test.\allowbreak{}c}): x86 faults on both, AArch64 answers 0 and INT\_\allowbreak{}MIN, and RISC-V answers −1 and INT\_\allowbreak{}MIN. A part the test does not know stops its build. 17 checks, 0 failed on three hosts: Windows and WSL as before, and a Raspberry Pi 5 (Linux aarch64, Cortex-A76, gcc; CPU part 0xd0b, 4 cores and 8 GB, as read from its \texttt{/\allowbreak{}proc/\allowbreak{}cpuinfo} and \texttt{free}). The Pi ran the test as a plain process at nice 10, since tessera does not build there yet: its seal, \texttt{obsignatio.\allowbreak{}cu}, and its link need nvcc. On the Pi 7 / 0 exited 0 printing 0, and INT\_\allowbreak{}MIN / −1 exited 0 printing −2147483648; neither trapped. The rest matched WSL: SIGSEGV for the address and the recursion, SIGILL for \texttt{udf \#0}, SIGTRAP for \texttt{brk \#0}, SIGABRT for \texttt{abort()}, the hang ended at 0.500207 s, and the program that is not there not started, errno 2.
\item \textbf{One signal for two rules, a reading.} POSIX gives the cell one signal for two rules twice over: SIGFPE for a zero divisor and for INT\_\allowbreak{}MIN / −1, and SIGSEGV for an address out of range and for a stack past its limit. Windows names each apart (0xC0000094 and 0xC0000095, 0xC0000005 and 0xC00000FD). The cell reports a stack overflow on POSIX as an address fault. Telling the two apart there needs the faulting address, read on an alternate signal stack (\texttt{sigaltstack}, \texttt{si\_\allowbreak{}addr}), and that is not built. x86 raises \#DE for both divisions (cited from knowledge); Windows reports the two as 0xC0000094 and 0xC0000095, and Linux as SIGFPE for both. AArch64 answers both with no trap, 0 and INT\_\allowbreak{}MIN, measured on the Pi. One question, 7 / 0, tells x86 from AArch64. RISC-V's −1 is in the test from knowledge and not measured.
\end{itemize}
\item \textbf{Membership queries and illegal operations, run.} A membership query for every instruction the ruleset uses, each answer checked against the host's exact integers. Probes of illegal operations: division by zero, INT\_\allowbreak{}MIN / −1, an address out of range, and an instruction the part lacks. These are (e)'s readings, run.

\begin{itemize}
\item \textbf{Built and measured.} \texttt{ruleset\_\allowbreak{}reader.\allowbreak{}h} gives a ruleset's forms to anything that writes the target's text. \texttt{test/\allowbreak{}engine/\allowbreak{}compiler/\allowbreak{}interface/\allowbreak{}interface\_\allowbreak{}ptx\_\allowbreak{}probe\_\allowbreak{}main.\allowbreak{}cu} writes kernels from \texttt{ptx.\allowbreak{}krs} through them and runs one question a process under the cell: membership, where each of 55 forms answers 65,536 cases against the host's integers (25 questions), and the illegal operations: a load from address 16, a misaligned 32-bit load, trap, and \texttt{elect.\allowbreak{}sync}, which PTX gives sm\_\allowbreak{}90 and later. \texttt{interface\_\allowbreak{}ptx\_\allowbreak{}test} holds each answer and asks again from a fresh process after every death. \textbf{Measured}, a tessera job, both rulesets (the engine's own and the flagless one, item 5): 10 checks, 0 failed each. Membership: 25 questions, 0 with a case that differs; \texttt{word\_\allowbreak{}div} 62,093 defined cases and \texttt{wide\_\allowbreak{}div} 65,345 agree, the rest (a zero divisor) undefined, the device answering ffffffff. Illegal address 700, misaligned 716, trap refused with a CUDA error, \texttt{elect.\allowbreak{}sync} refused by the toolchain for this device (sm\_\allowbreak{}86), and a fresh probe after each answers 6 + 7 = 0x0000000d.
\item \textbf{Not proved:} \texttt{cell\_\allowbreak{}ptx} gives 2 failed in the matrix run, undiagnosed (item 13).
\item \textbf{The wire, a want} (not built; \hyperref[chap:wants]{wants.md} posit 19). The probes of stages 4 and 5 turn on a transport, a bus, a network interface or a protocol's framing, each derived as a \texttt{.\allowbreak{}krs}. Listening is T⁻¹ (E4), and a peer is recognized by the saving I(s : self) = K(s) − K(s | self), a reading. This replaces stage 2's byte order and fences as the direction. Boundary, proposed, Doug's to rule on: hardware we own, and wires we are entitled to hear.
\item \textbf{Identity and the forgery guard, wants} (not built; answers to posit 19's reading, \hyperref[chap:wants]{wants.md}). Identity is proved where a time series is encoded into it. Every cell carries an identity from its spawn: a secret drawn from the noise, with keys derived from it and its lineage. A hash chain over the cell's series, its head signed under the spawn key, proves the series' holder and order. It does not prove that the series is true, a reading. The forgery guard is not needed yet. The forgery guard, meaning a check that refuses bytes or a claim of identity made by anyone but the cell named, is a want. It is not scheduled.
\end{itemize}
\item \textbf{More basic constructions} ((e)). When a question kills its process, or its answer disagrees with the host, that instruction leaves L. Its operation is built from more basic ones and asked again: a quotient from subtraction and shifts, a product from sums, a carry from a compare. This is (e)'s “more and more basic, fully unbound”.

\begin{itemize}
\item \textbf{Built and measured.} A ruleset entry may be a construct, \texttt{construct <name> <parameters>} through \texttt{end}, each line a form of the same ruleset, with \texttt{\{p\}} a parameter and \texttt{\{bank:n\}} the n-th scratch register of a bank, fresh for each use. A construct is expanded wherever its name is written, recursively, and its scratch comes from the step's own temporaries. \texttt{test/\allowbreak{}engine/\allowbreak{}compiler/\allowbreak{}codegen/\allowbreak{}rulesets/\allowbreak{}flagless/\allowbreak{}\{ptx,c\}.\allowbreak{}krs} is \texttt{ptx.\allowbreak{}krs} with no condition code: the 13 carry and borrow forms rebuilt as constructs on a \texttt{\%carry} register, RISC-V's shape, which has no carry flag. \textbf{Measured}, each a tessera job under \texttt{CYCLE\_\allowbreak{}RECORD\_\allowbreak{}CHECK=1}: the eleven record tests on the engine's own rulesets, the caches seeded from 2b's, as PTX 11 of 11 at 0 failed, read from the cache 203, written anew 0, and as C source 11 of 11, read from the cache 196, written anew 0; on the flagless ruleset against a fresh cache, 11 of 11 at 0 failed, 151 programs written as PTX and 5 as C, every launch held to the interpreter.
\end{itemize}
\item \textbf{A second target}, AArch64 first, then RISC-V, each as its own L ((b)). It is proved by the same record programs matching the host word for word ((c)).

\begin{itemize}
\item \textbf{The host oracle on a second part, built and measured} (stage 6's first step). \texttt{test/\allowbreak{}engine/\allowbreak{}compiler/\allowbreak{}cycle/\allowbreak{}record\_\allowbreak{}host\_\allowbreak{}test.\allowbreak{}c} imprints, lays and runs (\texttt{cycle\_\allowbreak{}record\_\allowbreak{}run\_\allowbreak{}host}) programs over all 18 record operations on 4,096 lanes, and each program prints an FNV-1a digest of its inputs and of every output word, taken from each word's lowest byte. The part's byte order therefore does not enter. It also holds the reused register layout to the unreused one, and a zero divisor and an inexact exact quotient to a refused run. It builds with no CUDA toolchain where there is none. \textbf{Measured}, each a tessera job, the same lines on three parts, 9 checks, 0 failed on each: Windows (MSVC, x86-64), WSL (gcc, x86\_\allowbreak{}64) and a Raspberry Pi 5 (gcc, AArch64, under the tessera built there with no CUDA toolchain), \texttt{ANCHOR\_\allowbreak{}EXACT\_\allowbreak{}LIMBS} 128. Records: arithmetic fefaad30b21bf7be (the same with registers reused), division 21c04f0b581ef777, bitwise f1bc8fdc0316c355, tables 39b3224a57bbc124, members 468a0e0e48772ff3. The Pi's build gave no warnings (gcc -Wall -Wextra). The device's record tests hold the device to the x86 host word for word. The Pi's host therefore agrees with the device too.
\item \textbf{Plan} (not built). AArch64 runs natively on the Raspberry Pi 5, with \texttt{cycle\_\allowbreak{}record\_\allowbreak{}run\_\allowbreak{}host} (\texttt{cycle.\allowbreak{}c}, portable C) as the host oracle there. RISC-V follows on an ESP32-C6 and an ESP32-P4, each a 32-bit RISC-V part (from knowledge). An ESP32-S3 is Xtensa.
\item \textbf{The ESP32's JTAG, a reading} (cited from knowledge). The ESP32 has JTAG. Over JTAG the host holds the cell and the part is the probe. A fault halts the core, and the host reads its cause and address from it: \texttt{mcause}, \texttt{mepc} and \texttt{mtval} on RISC-V. The C6, the P4 and the S3 carry a USB-Serial-JTAG bridge on the chip, Espressif's OpenOCD drives it, ESP-IDF's panic handler halts when a debugger is attached, and \texttt{ebreak} enters debug mode when \texttt{dcsr.\allowbreak{}ebreakm} is set. Read from the halted core, a stack overflow can be told from an address fault.
\item \textbf{The form limit} (built). An overflow in keymath's linear forms is caught: \texttt{keymath\_\allowbreak{}sum\_\allowbreak{}held} tests a sum before forming it. \texttt{record\_\allowbreak{}host\_\allowbreak{}test} carries a check, “form limit”. \textbf{Measured}: 11 passed and 0 failed on Windows and on the Pi. The form limit records \texttt{8e0b1ff8af17a325}.
\item \textbf{VHDL, a target on the Pi} (built). The register lane is a base class, \texttt{Code\allowbreak{}Generator} (\texttt{code\_\allowbreak{}generator.\allowbreak{}h}), and \texttt{Ptx\allowbreak{}Target} (\texttt{ptx\_\allowbreak{}target.\allowbreak{}h}) and \texttt{Vhdl\allowbreak{}Target} (\texttt{vhdl\_\allowbreak{}target.\allowbreak{}h}) are languages over it. The rulesets gain the bank \texttt{immediate}, \texttt{vhdl.\allowbreak{}krs} is the VHDL ruleset (\texttt{vhdl\_\allowbreak{}target.\allowbreak{}h}, GHDL analyzes and runs it), and \texttt{record\_\allowbreak{}vhdl\_\allowbreak{}test.\allowbreak{}cpp} holds its programs to the host. \textbf{Measured}: as PTX, 11 of 11 with 203 programs read from the cache and none written anew; with no flags set, 11 of 11; the cell's PTX test, 10 passed and 0 failed; the host test, 11 passed and 0 failed. On the Pi under GHDL 5.0.1, bitwise, tables, members and form limit equal the host. 5 programs are not held, since they call the operator block.

\begin{itemize}
\item \textbf{The lane in clock states} (built). State forms cut the lane into clock states, and \texttt{vhdl.\allowbreak{}krs} writes a clocked entity; the tree holds \texttt{vhdl\_\allowbreak{}target.\allowbreak{}h}'s \texttt{scheduled()} and \texttt{program\_\allowbreak{}schedule\_\allowbreak{}model()} and \texttt{test/\allowbreak{}engine/\allowbreak{}compiler/\allowbreak{}codegen/\allowbreak{}vhdl\_\allowbreak{}construction\_\allowbreak{}set.\allowbreak{}cpp}. \textbf{Measured, on the Pi:} 12 passed and 0 failed, and every program equals the host both cut only where it must be (4 states) and with every form alone (141 to 289 states). \texttt{ghdl -\allowbreak{}-\allowbreak{}synth} accepts it. For form limit with every form alone over 256 words, Yosys gives 167,697 generic cells and a longest path of 38 cells, and infers no RAM yet: the memory becomes flip-flops.
\end{itemize}
\end{itemize}
\item \textbf{Clustering.} The footprint of one cell, its bytes and its processes, and how many cells run at once, is worked as arithmetic first. The cells cluster, since each occupies very little memory.
\item \textbf{Vertical growth, then the bootstrap.} First a plan for the transform in which more than one program occupies a register vertically and never horizontally (posit 3, \hyperref[chap:wants]{wants.md}, theory), and the plan comes before any build. Then (a)'s bootstrap fixed point.

\begin{itemize}
\item \textbf{How a cell's gain survives, a question for Doug before the plan.} A gain escapes the noise floor and keeps enough energy to propagate, not to fall back down and dissipate through localized cancellation. Any tipped scale above the mean-field noise energy is irrevocable.
\item \textbf{Two thresholds, a reading.} The first quote names two. Escaping the noise floor is nucleation: a gain has to pass a barrier before it can grow at all (Volmer and Weber 1926; Becker and Döring 1935). Propagating is criticality: a line whose members leave a mean of k copies a step dies out for certain at k ≤ 1, unless every member leaves exactly one, and survives with a chance above 0 only at k > 1 (Harris 1963).
\item \textbf{Irrevocable, a reading.} A path's reversal is e\textasciicircum{}(−ΔS/k\_\allowbreak{}B) times as likely as the path, ΔS the entropy the path produces (Crooks 1999). A tipped scale is irrevocable by degree: the odds of undoing it fall exponentially in its margin over the noise and reach 0 at no finite margin, and the mean field's noise energy marks no step in them.
\item \textbf{A new gain's survival, a reading.} A variant with a small advantage s, starting as one copy, survives its early rare phase with probability about 2s where offspring counts are Poisson (Haldane 1927). Started as k independent copies, it is lost with probability about (1 − 2s)\textasciicircum{}k ≈ e\textasciicircum{}(−2sk). At s = 0.01 one copy is lost about 98 times in 100, and about 230 copies bring the loss to 1 in 100 (Derived).
\item \textbf{The question for the plan.} How many copies a new variant starts with, or a protected nursery that holds it until it is established. The plan waits on Doug's answer.
\item \textbf{Prior art, cited from knowledge.} J. B. S. Haldane, “A mathematical theory of natural and artificial selection, Part V: Selection and mutation”, Proc. Camb. Phil. Soc. 23 (1927) 838–844. G. E. Crooks, “Entropy production fluctuation theorem and the nonequilibrium work relation for free energy differences”, Phys. Rev. E 60 (1999) 2721–2726. T. E. Harris, The Theory of Branching Processes (Springer, 1963). M. Volmer and A. Weber, “Keimbildung in übersättigten Gebilden”, Z. phys. Chem. 119 (1926) 277–301. R. Becker and W. Döring, “Kinetische Behandlung der Keimbildung in übersättigten Dämpfen”, Ann. Phys. 24 (1935) 719–752.
\end{itemize}
\end{enumerate}

\begin{itemize}
\item \textbf{The question the cells are selected for} (a want, not built). Entropy is not simulated but adopted as a core principle, the cells are selected for one fixed thing, and a wider question can be asked that makes them evolve to answer it.

\begin{itemize}
\item \textbf{A moving target, a reading.} The cells are selected now for one fixed thing. The want is a target that moves or stays open, a question no cell answers once and for all. Ω is the candidate: each of its bits is a finite fact, and no single program computes them all (Chaitin 1975). The engine's own Ω run is item 12.
\item \textbf{Prior art, cited from knowledge.} Novelty search drops the objective and selects for behavior not seen before (J. Lehman and K. O. Stanley, “Abandoning objectives: evolution through the search for novelty alone”, Evolutionary Computation 19 (2011) 189–223). POET evolves the environments beside the agents that solve them (R. Wang, J. Lehman, J. Clune and K. O. Stanley, “Paired Open-Ended Trailblazer (POET): endlessly generating increasingly complex and diverse learning environments and their solutions”, arXiv:1901.01753, 2019). Tierra and Avida are the counter-reading, systems built for open-ended evolution whose novelty leveled off in the runs reported (T. S. Ray, “An approach to the synthesis of life”, Artificial Life II, 1992, 371–408; C. Ofria and C. O. Wilke, “Avida: a software platform for research in computational evolutionary biology”, Artificial Life 10 (2004) 191–229), and open-ended evolution is still counted an open problem (T. Taylor et al., “Open-ended evolution: perspectives from the OEE workshop in York”, Artificial Life 22 (2016) 408–423). G. J. Chaitin, “A theory of program size formally identical to information theory”, J. ACM 22 (1975) 329–340.
\end{itemize}
\end{itemize}
\item \textbf{The name.} The library is orior.

\begin{itemize}
\item \textbf{The name, a reading.} The name is Doug's, and it is Latin (Telos, τέλος, “end” or “purpose”, is Greek). The conclusion that this leads inevitably to AGI is Doug's expectation, and no figure in the tree bears on it.
\end{itemize}
\end{itemize}
\item \textbf{Ω's device phase has no budget} (\texttt{chaitin\_\allowbreak{}omega.\allowbreak{}cu}). Recorded only. A budget would change the sim's behavior, and that is Doug's call.

\begin{itemize}
\item \textbf{The exits.} The engine's loop runs while a job waits or a term is live, and no count of rounds, steps or tokens ends it: “There is no step budget and no token budget”. A term leaves the device by one of four exits: a normal form halts it, Brent's watcher proves a loop, \texttt{omega\_\allowbreak{}grows\_\allowbreak{}forever} proves growth, or the device parks it as outgrown.
\item \textbf{Parking is memory pressure.} While the device cannot hold a round's 19 rooms, the stepping term with the largest reduct is marked grown and the reducts are laid again without it, one term at a time. A term that keeps growing with no growth proof, and still fits, keeps the run going.
\item \textbf{The 32-bit index.} A round whose pool passes 2\textasciicircum{}32 − 1 tokens fails a check and stops the run (\texttt{OMEGA\_\allowbreak{}ENGINE\_\allowbreak{}INDEX\_\allowbreak{}MOST}). It is a hard stop, and it settles no term.
\item \textbf{The cross check.} The terms the engine settled are run again by \texttt{omega\_\allowbreak{}run} through 24 bits (\texttt{OMEGA\_\allowbreak{}ENGINE\_\allowbreak{}CROSS\_\allowbreak{}MOST}). At the default L of 28 lengths 25 to 28 are not cross-checked.
\item \textbf{Measured, the record of runs} (every printed run found on this machine):

\begin{itemize}
\item through 20 bits: finished, “5 rounds, 8010 records swept, 0 parked as outgrown, 118 ms” (16 checks, 0 failed);
\item 20 jobs: round 4096 with 5 live, 18 of 20 jobs done, 194,378,537 records swept, the largest term 57,330 tokens, 25,283 ms, with 5 live from round 16 on and no final line;
\item the default L of 28 with the cross check raised to 28 (a tessera job): 88 live from round 64 on, 18 of 24 jobs done from round 8 on, and the largest term 33,787 tokens at round 64, 1,081,339 at 128 and 138,412,027 at 256. Round 256 took 1,427,203 ms and swept 6,038,215,855 records. It is stopped there with 1,398 CPU s spent and no check reached, and the cross check returns to 24.
\end{itemize}
\item No run at the default L of 28 has finished.
\end{itemize}
\item \textbf{The engine in categories.} The engine files sit in category folders, 300 to 500 lines each, with no leftover tokens or mangled names and no British definition in the file names.

\begin{itemize}
\item \textbf{The categories.} \texttt{src/\allowbreak{}engine} is in category folders, and the test tree mirrors it under \texttt{test/\allowbreak{}engine}:

\begin{itemize}
\item arithmetic: \texttt{no\_\allowbreak{}rounding}, \texttt{decimal\_\allowbreak{}double}, \texttt{double\_\allowbreak{}fields};
\item analysis: \texttt{compression}, \texttt{entropy\_\allowbreak{}history}, \texttt{golden\_\allowbreak{}bands}, \texttt{residual}, \texttt{residual\_\allowbreak{}survey}, \texttt{shift\_\allowbreak{}agreement}, \texttt{unit\_\allowbreak{}sweep}, \texttt{tower}, \texttt{period} and \texttt{noise\_\allowbreak{}detector};
\item codecs: \texttt{blosc}, \texttt{deflate}, \texttt{inflate}, \texttt{lz4}, \texttt{snappy}, \texttt{zstd}, \texttt{crc} and \texttt{zip};
\item compiler: \texttt{cycle}, \texttt{codegen}, \texttt{emit}, \texttt{keymath}, \texttt{key\_\allowbreak{}schedule}, \texttt{cell}, \texttt{bootstrap}, \texttt{cubin} and \texttt{gnascor};
\item formats: \texttt{cfg\_\allowbreak{}json}, \texttt{hdf5}, \texttt{nifti}, \texttt{npy}, \texttt{nrrd}, \texttt{stack}, \texttt{tiff}, \texttt{krep}, \texttt{dicom} and \texttt{zarr};
\item quantum: \texttt{qasm};
\item runtime: \texttt{daemon} (its \texttt{service} included), \texttt{device\_\allowbreak{}pool}, \texttt{schedule}, \texttt{scriptura}, \texttt{obsignatio} and \texttt{radix\_\allowbreak{}keys};
\item \texttt{nbody}, \texttt{prg\_\allowbreak{}sch}, \texttt{render} and \texttt{sims} are categories of their own.
\end{itemize}
\item \textbf{Splits.} No C, C++ or CUDA file under \texttt{src/\allowbreak{}engine} or \texttt{test/\allowbreak{}engine} is over 500 lines. The pieces of a file share an \texttt{*\_\allowbreak{}internal.\allowbreak{}h}. The cycle module is sixteen files in \texttt{compiler/\allowbreak{}cycle/\allowbreak{}} (the compile's cache, route and toolchain, the launch, the prelude, the record kernel and launch, the sweep, the host's \texttt{cycle.\allowbreak{}c}, and their headers).
\item \textbf{Names.} The emitter is the code generator (\texttt{codegen/\allowbreak{}}: \texttt{Target}, \texttt{Code\allowbreak{}Generator}, \texttt{Machine\allowbreak{}Instr}, \texttt{Opcode}, \texttt{Register\allowbreak{}Class}, \texttt{Schedule}, \texttt{Asm\allowbreak{}Printer}). The text program is the assembly printer (\texttt{asm\_\allowbreak{}printer}), and \texttt{EMIT\_\allowbreak{}DEVICE} is \texttt{CODEGEN\_\allowbreak{}DEVICE}. \texttt{ENGINE\_\allowbreak{}RECORD\_\allowbreak{}LIMBS\_\allowbreak{}MOST} is \texttt{ENGINE\_\allowbreak{}RECORD\_\allowbreak{}LIMBS\_\allowbreak{}MAX}, most is max, tally is results, reading is measurement, good is ok, \texttt{X\_\allowbreak{}REFUSED} is \texttt{X\_\allowbreak{}ERROR}, and \texttt{X\_\allowbreak{}TOOK} is \texttt{X\_\allowbreak{}STATUS\_\allowbreak{}CHECK}. Kept: tessera, scriptura, obsignatio, signum, crystal, steer, climb, lane, ruleset, the \texttt{.\allowbreak{}krs} keywords, and tessera's admit, held and hold\_\allowbreak{}until.
\item \textbf{Definition.} American throughout, file names and outputs included: \texttt{noise\_\allowbreak{}detector\_\allowbreak{}neighbors.\allowbreak{}cu}, \texttt{maint/\allowbreak{}chain/\allowbreak{}fetch\_\allowbreak{}labeled.\allowbreak{}py} and \texttt{blocks\_\allowbreak{}labeled.\allowbreak{}json} carry American definition, and the engine writes \texttt{noise\_\allowbreak{}neighbors.\allowbreak{}tsv}, \texttt{pi\_\allowbreak{}rings\_\allowbreak{}center.\allowbreak{}png} and \texttt{shuffled\_\allowbreak{}rings\_\allowbreak{}center.\allowbreak{}png}, and the flicker table's columns \texttt{neighbor\_\allowbreak{}pairs} and \texttt{neighbor\_\allowbreak{}squares}. The signed \texttt{maint/\allowbreak{}signing/\allowbreak{}manifest.\allowbreak{}tsv} still names the two chain files by their British definitions under \texttt{tools/\allowbreak{}chain/\allowbreak{}}, and re-signing it is Doug's.
\item \textbf{Format.} clang-format with the repository's \texttt{.\allowbreak{}clang-\allowbreak{}format} over all 604 engine files.
\item \textbf{Not proved.} The whole test matrix has not run. Three failures stand undiagnosed: the qasm build, \texttt{cell\_\allowbreak{}ptx} 2 failed, and a \texttt{record\_\allowbreak{}bitwise\_\allowbreak{}c} hang. The bootstrap route (item 11 (a)) has not run.
\end{itemize}
\item \textbf{The ruleset asked of the part.} The compiler derives a target's ruleset by putting questions to the target and reading what comes back. sm\_\allowbreak{}86 is the first subject instead of the target. nvdisasm and ptxas are the oracle: no vendor source is read and no shipped ELF is taken apart to learn a rule.

\begin{itemize}
\item \textbf{Three searches, each wider than the last.} A listing gives the forms some compiler happened to write. Widening turns each of a form's 128 bits over and keeps every encoding the disassembler names whole. The sweep cuts the opcode out of a carrier the part ran and puts all 4096 values back, from each of 8 carriers, which no compiler's output bounds.
\item \textbf{A form is its operation and the kind of each operand together.} \texttt{SHF.\allowbreak{}L.\allowbreak{}U32} with a register in the shift slot and \texttt{SHF.\allowbreak{}L.\allowbreak{}U32} with a number there carry one name and are two forms, and the assembler puts a number in a different place for each. The search had kept a decode only where it named an operation other than the one it came from, and had widened only the 118 forms a listing gave, never the 433 it found.
\item \textbf{The kind of a slot is a three bit field.} On sm\_\allowbreak{}86 the low word reads 0x2 for a register in the second and the third operand, 0xa and 0x8 for a constant and a number in the second, 0x6 and 0x4 for a constant and a number in the third. A register to a number is two bits, through the constant that lies between. One round reaches the constant and the round after reaches the number. \texttt{SHF.\allowbreak{}L.\allowbreak{}S32}, \texttt{SHF.\allowbreak{}L.\allowbreak{}S64}, \texttt{SHF.\allowbreak{}L.\allowbreak{}U64} and \texttt{SHF.\allowbreak{}L.\allowbreak{}W.\allowbreak{}U32} had their number forms only because the sweep's carriers happened to hold those mode bits, and the sweep holds the high word fixed.
\item \textbf{Measured} (RTX 3070, sm\_\allowbreak{}86, CUDA 13.3): 118 forms listed widen to 2436 over 2 rounds, the sweep adds 373 on top, and the part's machine file holds 2927 forms. 0 without fields, 0 differing when read back, 0 past what it holds. Of about 3200 generated lines across 16 programs, 88 find no form.
\item \textbf{The walk is held at 2 rounds.} Run to its own end it does not close: 3128 forms from 118 in 19 minutes, flat at 177 a minute, no plateau. The component reachable a bit at a time is most of what the part decodes. It would pass the 16384 a machine holds and end in \texttt{refused}, reporting nothing, and a form reached far out carries the operand bits of a chain of forms it has nothing to do with, which is \texttt{PLOP3.\allowbreak{}LUT} at depth.
\item \textbf{The 88 that are left are not a search gap.} All 34 forms the part has with a number in the first source slot are single source: BREV, FLO, IABS, I2I, LDC, B2R, BAR. No operation taking three sources has one there. \texttt{SEL Rd, Ra, Rb, P} takes a register at Ra and no number. These are sass.krs's \texttt{wide\_\allowbreak{}select} and \texttt{word\_\allowbreak{}select} handed a number for \texttt{chosen}, and the fix is the ruleset's or the generator's.
\item \textbf{Registers.} R0 through R254 hold and R255 is RZ. What holds is gated by the count the ELF declares: R238 and R254 error under a frame declaring fewer. \texttt{register\_\allowbreak{}file\_\allowbreak{}holds()} is 238.
\item \textbf{Three comparisons no compiler wrote.} ptxas reads zero and below off the negations of NE and GE. \texttt{ISETP.\allowbreak{}EQ.\allowbreak{}U32.\allowbreak{}AND}, \texttt{ISETP.\allowbreak{}EQ.\allowbreak{}U32.\allowbreak{}AND.\allowbreak{}EX} and \texttt{ISETP.\allowbreak{}LT.\allowbreak{}AND} were in no listing and nothing had run them. Each was reached one bit out, and each answered right on the part.
\item \textbf{The rotate the part refuses.} The funnel idiom every compiler writes was asked first on \texttt{SHF.\allowbreak{}R.\allowbreak{}U32.\allowbreak{}HI} and \texttt{SHF.\allowbreak{}L.\allowbreak{}U32} with one register handed to both halves, and the part answered wrong: \texttt{.\allowbreak{}U32} has no funnel. \texttt{SHF.\allowbreak{}R.\allowbreak{}U64} and \texttt{SHF.\allowbreak{}L.\allowbreak{}U64.\allowbreak{}HI} answer right, and both came out of widening instead of out of any compiler.
\item \textbf{The resident is not asked of the ruleset.} ptxas turns ptx.krs's 110 PTX instructions into 248 SASS that runs. A program is put together by writing into that cubin's \texttt{cycle\_\allowbreak{}lane}, leaving \texttt{cycle\_\allowbreak{}program} as ptxas wrote it, and 247 of the 248 survive. \texttt{program\_\allowbreak{}unit\_\allowbreak{}written()} answers 0 for SASS and sass.krs is never asked for a resident. Its 6 refusals are fail closed and correct: a constant bank past 0, \texttt{PLOP3.\allowbreak{}LUT} operand 6 with no field, \texttt{BAR.\allowbreak{}SYNC.\allowbreak{}DEFER\_\allowbreak{}BLOCKING} operand 0 with no field, and a call to a symbol its form was not learned with.
\item \textbf{The listing is not the instruction.} \texttt{LDG}'s bits 32-39 change nothing nvdisasm prints. Two instructions that print the same can differ there. 968 of 976 match byte for byte and all 976 read back right.
\item \textbf{A loop asked of the part} (Q16, A18). \texttt{loop\_\allowbreak{}back\_\allowbreak{}if} is derived by asking: every form in the machine file is put in its place with its operands filled by their kinds, at the end of a body that counts N down with \texttt{word\_\allowbreak{}add} and sets the flag with \texttt{test\_\allowbreak{}word\_\allowbreak{}nonzero}. On sm\_\allowbreak{}86, 2533 of the 2927 forms assemble, 1715 fall through and 8 come back on every N, every one a \texttt{BRA}. None is cheaper than the next above the spread of its own runs, and \texttt{sass.\allowbreak{}krs} keeps its \texttt{BRA}. \texttt{interface\_\allowbreak{}sass\_\allowbreak{}probe} given \texttt{loop} and a machine file puts this ask alone against the cubins of an earlier run (\texttt{SASS\_\allowbreak{}PATTERN} in \texttt{interface\_\allowbreak{}sass\_\allowbreak{}test.\allowbreak{}sh}).
\item \textbf{Read back with no disassembler.} \texttt{sass\_\allowbreak{}encoding\_\allowbreak{}read} reads an encoding back into its instruction through the machine file's forms, and \texttt{sass\_\allowbreak{}loop\_\allowbreak{}walk} checks one emitted instruction as a loop's way back, on or off, and never changes it (\texttt{compiler/\allowbreak{}cubin/\allowbreak{}sass\_\allowbreak{}assemble.\allowbreak{}h}). The walk holds that the guard is the flag alone, that one label or number operand added to the address after the instruction lands on the label, and that the first operand writes nothing the loop keeps. A caller can stop it at any step, and with it off what was emitted is what runs, as code written to run in constant time needs. It gives the expected step on ten known cases out of ten, and agrees with all 8 forms above.
\item \textbf{A branch's distance starts at bit 34}, in four-byte steps. Bits 32 and 33 are the operation's own, as \texttt{BRA}, \texttt{BRA.\allowbreak{}U} and \texttt{BRA.\allowbreak{}DIV} show the same distance with 0, 1 and 2 there. Of the 2927 forms, 2383 read back as they were listed, and the rest are known limits: two names the disassembler prints for one encoding (\texttt{IMAD.\allowbreak{}MOV}), addresses with a uniform register, absolute 64-bit \texttt{CALL.\allowbreak{}ABS} and \texttt{JMP} targets, and 380 forms whose operands the assembler cannot place.
\item \textbf{Not proved.} The suite has not run since the machine file was replaced. The run that wrote it passed every check except the comparison against the committed file it replaces.
\end{itemize}
\end{enumerate}
